\documentclass[10pt,a4paper]{book}
\usepackage[utf8]{inputenc}
\usepackage[spanish]{babel}
\usepackage{amsmath}
\usepackage{amsfonts}
\usepackage{amssymb}
\usepackage{graphicx}
\usepackage[hidelinks]{hyperref}

\usepackage[left=2cm,right=2cm,top=2cm,bottom=2cm]{geometry}

\title{\textbf{ANÁLISIS DE SERIES DE TIEMPO NO LINEALES}}
\author{HOLGER KANTZ Y THOMAS SCHREIBER\\\textit{traducido por \textbf{natits}}}
\date{}

\begin{document}
\maketitle
\newpage
\tableofcontents
\newpage
\chapter{Temas básicos}
\section{Introducción: ¿por qué métodos no lineales?}
Probablemente estés leyendo este libro porque tienes una fuente interesante de datos y sospechas que no es lineal. O bien sabes positivamente que no es lineal porque tienes alguna idea de lo que está ocurriendo en la parte del mundo que estás observando, o bien sospechas que lo es porque has intentado un análisis de datos lineal y no quedaste satisfecho con sus resultados.

Los métodos lineales interpretan toda estructura regular en un conjunto de datos, como una frecuencia dominante, a través de correlaciones lineales (que se definirán en el Capítulo 2). Esto significa, en resumen, que la dinámica intrínseca del sistema está gobernada por el paradigma lineal según el cual causas pequeñas producen efectos pequeños. Dado que las ecuaciones lineales solo pueden dar lugar a soluciones que decaen (o crecen) exponencialmente o que oscilan periódicamente de manera amortiguada, todo comportamiento irregular del sistema debe atribuirse a algún tipo de entrada externa aleatoria. Ahora bien, la teoría del caos nos ha enseñado que la entrada aleatoria no es la única fuente posible de irregularidad en la salida de un sistema: los sistemas caóticos no lineales pueden generar datos muy irregulares con ecuaciones de movimiento puramente deterministas, de manera autónoma, es decir, sin términos dependientes del tiempo. Por supuesto, un sistema que posea tanto no linealidad como entradas aleatorias muy probablemente producirá datos irregulares también.

Aunque aún no hemos introducido las herramientas necesarias para hacer afirmaciones cuantitativas, veamos algunos ejemplos de conjuntos de datos reales. Ellos representan varios problemas diferentes de análisis de datos de los que se podría beneficiar un tratamiento con métodos no lineales, mientras que los métodos lineales resultarían inapropiados.

\subsection*{Ejemplo 1.1 (Datos de láser NMR)}

En un experimento de laboratorio realizado en 1995 por Flepp, Simonet y Brun en el Departamento de Física de la Universidad de Zúrich, se operó un láser de Resonancia Magnética Nuclear bajo condiciones tales que la amplitud de la salida del láser (en radiofrecuencia) varía de manera irregular en el tiempo. 

A partir de la configuración del experimento queda claro que el sistema es altamente no lineal y que el ruido de entrada aleatorio tiene una amplitud muy pequeña en comparación con la amplitud de la señal. Por lo tanto, no se asume que la irregularidad de la señal se deba únicamente al ruido de entrada. De hecho, ha sido posible modelar el sistema mediante un conjunto de ecuaciones diferenciales que no involucran ningún componente aleatorio en absoluto; véase Flepp \emph{et al.} (1991). El Apéndice B.2 contiene más detalles sobre este conjunto de datos.

Los valores sucesivos de la señal parecen ser muy erráticos, como puede verse en la Fig.~1.1. No obstante, como veremos más adelante, es posible realizar predicciones bastante precisas de valores futuros de la señal utilizando un algoritmo de predicción no lineal. La Figura~1.2 muestra el error medio de predicción en función de qué tan lejos en el futuro se realizan los pronósticos. De manera intuitiva, mientras más lejos en el futuro se hagan las predicciones, mayor será la incertidumbre. Después de aproximadamente 35 pasos temporales, la predicción se vuelve peor que si simplemente se utilizara el valor medio como predicción (línea horizontal). 

En horizontes de predicción cortos, el crecimiento del error de predicción puede aproximarse bien mediante una exponencial con un exponente de 0.3, lo cual se indica como línea discontinua. Usamos el método de predicción simple que se describirá en la Sección~4.2. Para comparación, también mostramos el resultado del mejor predictor \emph{lineal} que pudimos ajustar a los datos (cruces). Observamos que la capacidad de predicción debida a las correlaciones lineales en los datos es mucho más débil que la debida a la estructura determinista, en particular para horizontes de predicción cortos. La predictibilidad de la señal puede considerarse una firma de la naturaleza determinista del sistema. Véase la Sección~2.5 para detalles sobre el método de predicción lineal utilizado.

La estructura no lineal que conduce a la predictibilidad a corto plazo en este conjunto de datos no es evidente en una representación como la de la Fig.~1.1. Sin embargo, podemos hacerla visible graficando cada punto de datos contra su predecesor, como se hace en la Fig.~1.3. A esta gráfica se le denomina \emph{retrato de fases}. Esta representación es una aplicación particularmente simple de la \emph{incrustación con retardo temporal}, una herramienta básica que se usa con frecuencia en el análisis de series temporales no lineales. Este concepto será introducido formalmente en la Sección~3.2. En el presente caso, solo necesitamos una representación de los datos que sea imprimible en dos dimensiones.

\subsection*{Ejemplo 1.2 (Frecuencia respiratoria humana)}

Uno de los conjuntos de datos utilizados para la competencia de series temporales del Instituto Santa Fe en 1991 92 \textit{WeigendGershenfeld1993} fue proporcionado por A.~Goldberger del Beth Israel Hospital en Boston \textit{Rigney1993}; véase también el Apéndice B.5. De entre varios canales seleccionamos un registro de 16 minutos del flujo de aire a través de la nariz de un sujeto humano. Un gráfico de este segmento de datos se muestra en la Fig.~1.4.

En este caso se sabe muy poco sobre el origen de las fluctuaciones de la frecuencia respiratoria. El único indicio de que el comportamiento no lineal desempeña un papel proviene de los propios datos: la señal no es compatible con el supuesto de que haya sido generada por un proceso aleatorio gaussiano con solo correlaciones lineales (posiblemente distorsionado por una función de medición no lineal). Mostramos esto creando un conjunto de datos artificial que tiene exactamente las mismas propiedades \emph{lineales}, pero sin ningún determinismo adicional incorporado. Este conjunto de datos consiste en números aleatorios que han sido reescalados para que la distribución de los valores coincida con la del original (por lo tanto, la media y la varianza también son idénticas) y filtrados de modo que el espectro de potencias sea el mismo. (Cómo se hace esto, y otros aspectos adicionales, se discuten en el Capítulo~4 y la Sección~7.1). Si los datos medidos se describieran adecuadamente mediante un proceso lineal, no deberíamos encontrar diferencias significativas respecto a los artificiales.

Usemos nuevamente una \emph{incrustación con retardo temporal} para visualizar y comparar las series temporales original y artificial. Simplemente graficamos cada valor de la serie frente al valor tomado un tiempo de retardo $t$ antes. Encontramos que los dos \emph{retratos de fases} resultantes se ven cualitativamente diferentes. Parte de la estructura presente en los datos originales (Fig.~1.5, panel izquierdo) no se reproduce en la serie artificial (Fig.~1.5, panel derecho). Dado que ambas series tienen las mismas propiedades lineales, la diferencia probablemente se deba a la no linealidad en el sistema. Más significativamente aún, el conjunto de datos original es estadísticamente asimétrico bajo inversión temporal, lo cual se refleja en el hecho de que la Fig.~1.5, panel izquierdo, es no simétrica bajo reflexión respecto a la diagonal. Esta observación hace muy poco probable que la serie de datos represente una oscilación armónica ruidosa.

\subsection*{Ejemplo 1.3 (Datos de cuerda vibrante)}

Nick Tufillaro y Timothy Molteno, en el Departamento de Física de la Universidad de Otago, Dunedin, Nueva Zelanda, proporcionaron un par de series temporales (véase el Apéndice B.3) de una cuerda vibrante en un campo magnético.

La envolvente de la elongación del alambre experimenta oscilaciones que pueden ser caóticas, dependiendo de los parámetros elegidos. Uno de estos conjuntos de datos está dominado por un ciclo de período cinco, lo que significa que después de cada quinta oscilación el registro se (aproximadamente) repite. En la Fig.~1.6 mostramos las elongaciones $x$ e $y$ registradas simultáneamente. Trazamos una línea continua a lo largo de las primeras 1000 mediciones y colocamos un punto en cada una de las siguientes 4000. Observamos que el ciclo de período cinco no se mantiene perfecto durante el intervalo de medición. Esto se hace más evidente al graficar un punto por cada ciclo contra el número de ciclos. (Estos puntos se obtienen de manera sistemática mediante la llamada \emph{sección de Poincaré}; véase la Sección~3.5.) En la Fig.~1.7 se aprecia que el ciclo de período cinco se interrumpe ocasionalmente. Una explicación natural podría ser la presencia de influencias externas sobre el equipo que provocan que el sistema abandone el estado periódico. Sin embargo, una explicación mucho más apropiada de este conjunto de datos se da en el Ejemplo~8.3 de la Sección~8.4. Los episodios irregulares se deben en realidad a la \emph{intermitencia}, un fenómeno típico encontrado en sistemas dinámicos no lineales.

\subsection*{Ejemplo 1.4 (Electrocardiograma fetal)}

Consideremos el problema de procesar una señal para extraer el diminuto componente del electrocardiograma fetal a partir de los (pequeños) potenciales eléctricos en el abdomen de una mujer embarazada (trazo superior en la Fig.~1.8). La señal fetal, más rápida y más pequeña, no puede distinguirse con técnicas lineales clásicas porque tanto los electrocardiogramas maternos como los fetales tienen espectros de potencia de banda ancha. Sin embargo, utilizando una técnica no lineal de proyección en el espacio de fases, que originalmente fue desarrollada para la reducción de ruido en datos caóticos (Sección~10.3.2), es posible realizar la separación de manera automática. El trazo inferior en la Fig.~1.8 muestra el componente fetal extraído. Se ofrecen más explicaciones en la Sección~10.4. Nótese que no necesitamos suponer (ni en realidad esperamos) que el corazón sea un sistema caótico determinista. 

\bigskip

Existe una amplia gama de cuestiones que debemos abordar al hablar del análisis de series temporales no lineales. ¿Cómo podemos obtener los resultados más precisos y significativos para un conjunto de datos limpio y claramente determinista, como el del primer ejemplo? ¿Qué modificaciones son necesarias si el caso es menos claro? ¿Qué puede hacerse si, como en el ejemplo de la frecuencia respiratoria, todo lo que tenemos son indicios de que los datos no se describen adecuadamente mediante un modelo lineal con entradas gaussianas? Dependiendo del conjunto de datos y de la tarea de análisis que tengamos en mente, elegiremos diferentes enfoques. No obstante, hay una serie de ideas generales con las que conviene familiarizarse, independientemente del aspecto de los datos. Los conceptos más importantes para analizar datos complejos se presentarán en la Primera Parte de este libro. Los temas que son más teóricos, requieren datos de mayor calidad o tienen un interés más especializado se tratarán en la Segunda Parte.

\bigskip

Obviamente, en un campo tan diverso como el análisis de series temporales no lineales, cualquier selección de temas para un solo volumen debe ser incompleta. Sería ingenuo afirmar que todas las elecciones que hemos hecho se basan exclusivamente en criterios objetivos. Existe un fuerte sesgo hacia métodos que hemos encontrado conceptualmente interesantes o útiles en el trabajo práctico, o ambos. Qué métodos resulten más útiles depende del tipo de datos a analizar y, por tanto, parte de nuestro sesgo ha sido determinado por nuestros contactos con grupos experimentales. Finalmente, no somos excepción a la regla de que la gente prefiere los métodos en los que ha estado involucrada en su desarrollo.

Aunque creemos que algunos enfoques presentados en la literatura son efectivamente inútiles para el trabajo científico (como determinar exponentes de Lyapunov a partir de tan solo 16 puntos de datos), queremos enfatizar que otros métodos que mencionamos solo brevemente o no mencionamos en absoluto, pueden ser muy útiles para un problema concreto de series temporales. A continuación listamos algunas omisiones importantes (aparte de las que fueron dejadas de lado por ignorancia nuestra). Generalizaciones no lineales de los modelos ARMA y otros métodos populares entre los estadísticos se presentan en Tong (1990). Generalizando la función usual de autocorrelación de dos puntos se llega al \emph{bispectrum}. Este y modelos de series temporales relacionados se discuten en Subba Rao \& Gabr (1984). Dentro de la teoría de sistemas dinámicos, el análisis de órbitas periódicas inestables juega un papel importante. Por un lado, las órbitas periódicas son importantes para el estudio de la topología de un atractor. Se han obtenido resultados muy interesantes sobre la estructura plantilla, números de enlace, etc., pero el enfoque se limita a aquellos atractores de dimensión dos o mayor que pueden ser embebidos en un espacio de fases tridimensional. (Las variedades unidimensionales, es decir, trayectorias de sistemas dinámicos, no forman nudos en más de tres dimensiones). Remitimos al lector interesado a Tufillaro \emph{et al.} (1992). Por otro lado, las expansiones de órbitas periódicas constituyen una manera poderosa de calcular cantidades características que se definen como promedios sobre la medida natural, tales como dimensiones, entropías y exponentes de Lyapunov. 

En el caso de que un sistema sea \emph{hiperbólico}, es decir, que las direcciones de expansión y contracción no sean tangentes entre sí en ningún punto, se pueden derivar expresiones con convergencia exponencial para tales cantidades. Sin embargo, cuando la hiperbolicidad falta (el caso genérico), es necesario un número muy grande de órbitas periódicas para las expansiones de ciclo. Hasta ahora no se ha demostrado que este tipo de análisis sea factible con datos experimentales de series temporales. La teoría se explica en Artuso \emph{et al.} (1990) y en \emph{The Webbook} \textit{Cvitanovic2001}.

Después de una breve revisión de los conceptos básicos del análisis lineal de series temporales en el Capítulo~2, se introducirán las ideas más fundamentales del enfoque de la dinámica no lineal en los capítulos sobre el espacio de fases (Capítulo~3) y sobre predictibilidad (Capítulo~4). Estos dos capítulos son esenciales para la comprensión del resto del texto; de hecho, el concepto de una representación en espacio de fases, en lugar de una en el dominio del tiempo o de la frecuencia, es la característica distintiva del análisis de series temporales dinámicas no lineales. Otro concepto fundamental de la dinámica no lineal es la sensibilidad de los sistemas caóticos a los cambios en las condiciones iniciales, lo cual se discute en el capítulo sobre inestabilidad dinámica y el exponente de Lyapunov (Capítulo~5). Para que un sistema disipativo sea simultáneamente acotado e inestable, su trayectoria debe residir en un conjunto con propiedades geométricas inusuales. El estudio de cómo se investigan estas series temporales se aborda en el capítulo sobre la geometría de atractores y dimensiones fractales (Capítulo~6). Cada uno de estos dos últimos capítulos contiene los fundamentos de su tema, mientras que material adicional se presentará más adelante en el Capítulo~11. Relajaremos el requisito de determinismo en el Capítulo~7 (y más tarde en el Capítulo~12). Propondremos ideas generales para inspeccionar y estudiar datos complejos, incluyendo enfoques visuales y simbólicos. Además, estableceremos métodos estadísticos para caracterizar datos que carecen de evidencia sólida de determinismo, como el escalamiento o la autosimilitud. Finalmente, situaremos nuestras consideraciones dentro del contexto más amplio de la teoría de sistemas dinámicos no lineales en el Capítulo~8.

\bigskip

La segunda parte del libro contendrá material avanzado que puede ser de interés una vez que uno de los algoritmos básicos haya sido aplicado con éxito; obviamente no tiene sentido estimar todo el espectro de exponentes de Lyapunov si ni siquiera el mayor de ellos ha sido determinado de manera confiable. Como regla general, refiérase a la Primera Parte para obtener resultados iniciales. Después, consulte la Segunda Parte para obtener resultados \emph{óptimos}. Esta regla se aplica a las incrustaciones (Capítulo~9), al tratamiento de ruido (Capítulo~10), así como al modelado y predicción (Capítulo~12). Allí se incluyen métodos para tratar procesos no lineales con forzamiento estocástico. Un nivel matemático avanzado es necesario para el estudio de cantidades invariantes como espectros de Lyapunov y dimensiones generalizadas (Capítulo~11). Los avances recientes en el tratamiento de señales no estacionarias se presentan en el Capítulo~13. También se incluirá una breve descripción de la sincronización caótica en el Capítulo~14.

\bigskip

El control del caos es un campo enorme por sí mismo, pero dado que está estrechamente relacionado con el análisis de series temporales, se discutirá para concluir el texto en el Capítulo~15.

A lo largo del texto haremos referencia a implementaciones específicas de los métodos discutidos, que están incluidas en el paquete de software TISEAN.\footnote{Información de contexto sobre la implementación y sugerencias para el uso de los programas se ofrecen en el Apéndice~A.} El paquete TISEAN contiene una gran cantidad de rutinas, algunas de las cuales son implementaciones de técnicas estándar o incluidas por conveniencia (análisis espectral, números aleatorios, etc.), y no se discutirán más aquí. Nos centraremos en aquellas rutinas que implementan ideas esenciales del análisis de series temporales no lineales. Para más información, véase la documentación distribuida con el software.

El análisis de series temporales no lineales no está tan bien establecido ni es tan comprendido como su contraparte lineal. Aunque haremos todo lo posible por explicar las perspectivas y limitaciones de los métodos que introduciremos, será necesario familiarizarse con los algoritmos mediante el uso de datos artificiales donde se conozcan los resultados correctos. Aunque en los ejemplos hemos utilizado casi exclusivamente datos experimentales para ilustrar los conceptos discutidos en el texto, introduciremos varios modelos numéricos en los ejercicios. Invitamos al lector a resolver algunos de los problemas y a utilizar los datos artificiales antes de analizar datos medidos. Es una mala idea aplicar algoritmos poco familiares directamente a datos con propiedades desconocidas; es mejor practicar primero con datos que tengan propiedades conocidas, como el número de puntos de datos y la tasa de muestreo, el número de grados de libertad, la cantidad y naturaleza del ruido, etc. Para dar una idea de modelos populares, introduciremos el mapa logístico (Ejercicio~2.2), el mapa de Hénon (Ejercicio~3.1), el mapa de Ikeda (Ejercicio~6.3), un mapa tridimensional de Baier y Klein (Ejercicio~5.1), las ecuaciones de Lorenz (Ejercicio~3.2) y la ecuación diferencial retardada de Mackey–Glass (Ejercicio~7.2). Diversos modelos de ruido pueden considerarse para comparación, incluyendo el \emph{promedio móvil} (Ejercicio~2.1) y el \emph{movimiento browniano} (Ejercicio~5.3). Estos se utilizan únicamente como referencias en los ejercicios donde se introducen los métodos.

\section{Lecturas adicionales}

Evidentemente, cualquiera que desee aprovechar los métodos de series temporales basados en la teoría del caos mejorará su comprensión de los resultados aprendiendo más sobre dicha teoría. Existe un pequeño y recomendable libro acerca del azar y el caos escrito por Ruelle (1991). Una introducción accesible y sin demasiada carga matemática es el libro de Kaplan \& Glass (1995). Otros textos sobre el tema incluyen a Bergé \emph{et al.} (1986), Schuster (1988) y Ott (1993). Una presentación muy clara es la de Alligood \emph{et al.} (1997). Material más avanzado se encuentra en Eckmann \& Ruelle (1985) y en Katok \& Hasselblatt (1996). Muchos artículos relevantes para los aspectos prácticos del caos se reproducen en Ott \emph{et al.} (1994), complementados con una introducción general. Los aspectos teóricos de la dinámica caótica relevantes para el análisis de series temporales se discuten en Ruelle (1989).

Existen varios libros que cubren el análisis de series temporales caóticas, cada uno con un énfasis distinto. Abarbanel (1996) discute aspectos del análisis de series temporales con métodos de la teoría del caos, con un énfasis particular en el trabajo de su grupo sobre técnicas de vecinos más cercanos. El pequeño volumen de Diks (1999) se centra en varios aspectos particulares, incluyendo el efecto del ruido en datos caóticos, los cuales se discuten de manera bastante detallada. También hay monografías vinculadas a situaciones experimentales específicas, pero que contienen material de mayor interés general: Pawelzik (1991), Buzug (1994) (ambos en alemán), Tufillaro \emph{et al.} (1992) y Galka (2000). Muchos artículos interesantes sobre métodos no lineales en series temporales pueden encontrarse en los siguientes volúmenes de actas de conferencias: Mayer-Kress (1986), Casdagli \& Eubank (1992), Drazin \& King (1992) y Gershenfeld \& Weigend (1993). Artículos relevantes sobre dinámica no lineal y fisiología en particular se encuentran en Bélair \emph{et al.} (1995). Revisiones importantes, cada una con puntos de vista algo diferentes, se hallan en Grassberger \emph{et al.} (1991), Casdagli \emph{et al.} (1991), Abarbanel \emph{et al.} (1993), Gershenfeld \& Weigend (1993) y Schreiber (1999). El paquete de software TISEAN está acompañado por un artículo de Hegger \emph{et al.} (1999), que también puede usarse como una visión general de los métodos de series temporales.

Para un estadístico interesado en la visión de las series temporales no lineales, véanse los libros de Priestley (1988) o de Tong (1990). Además, ha habido una serie de talleres en los que especialistas de diferentes campos se reúnen. Las actas de estos encuentros también contienen valiosos artículos de referencia.

\newpage
\chapter{Herramientas lineales y consideraciones generales}

\section{Estacionariedad y muestreo}

De manera general, una medición científica de cualquier tipo es en principio más útil cuanto más reproducible sea. Necesitamos saber que los números que medimos corresponden a propiedades del objeto estudiado, hasta cierto margen de error de medición. En el caso de las mediciones de series temporales, la reproducibilidad está estrechamente vinculada a dos nociones diferentes de estacionariedad.

La forma más débil pero más evidente de \emph{estacionariedad} requiere que todos los parámetros relevantes para la dinámica de un sistema sean fijos y constantes durante el periodo de medición (y que dichos parámetros se mantengan iguales cuando se reproduce el experimento). Este requisito debe cumplirse no solo por el diseño experimental, sino también por el propio proceso que tiene lugar en este entorno fijo. En un inicio esto puede parecer desconcertante, ya que uno suele suponer que parámetros externos constantes inducen un proceso estacionario, pero de hecho veremos en este libro varios casos donde esto no se cumple.  

Si el proceso bajo observación es probabilístico, se caracterizará mediante distribuciones de probabilidad para las variables involucradas. En un proceso estacionario, estas probabilidades no deben depender del tiempo. Lo mismo ocurre si el proceso se describe mediante un conjunto de \emph{probabilidades de transición} entre distintos estados. Si existen reglas deterministas que gobiernan la dinámica, estas reglas no deben cambiar durante el intervalo de tiempo cubierto por una serie temporal.

En algunos casos, podemos manejar un cambio simple de un parámetro una vez que este cambio ha sido detectado. Si, por ejemplo, la calibración del aparato de medición deriva, podemos intentar reescalar los datos continuamente para mantener constante la media y la varianza. Sin embargo, esto puede ser peligroso, a menos que estemos \emph{seguros} de que solo la escala de medición, y no la dinámica, se está modificando. Una modulación estrictamente periódica de un parámetro puede interpretarse como una variable dinámica más bien que como un parámetro, y no destruye la estacionariedad.

\bigskip

Desafortunadamente, en la mayoría de los casos no tenemos acceso directo al sistema que produce una señal y no podemos establecer evidencia de que sus parámetros sean efectivamente constantes. Por ello debemos formular un segundo concepto de \emph{estacionariedad}, basado en los propios datos disponibles. Este concepto debe ser distinto, ya que existen muchos procesos que son formalmente estacionarios en el límite de tiempos de observación infinitamente largos, pero que se comportan efectivamente como procesos no estacionarios cuando se estudian en intervalos finitos. Un fenómeno destacado que pertenece a esta clase se denomina \emph{intermitencia}, y será discutido en la Sección~8.4.

Una serie temporal, al igual que cualquier otra medición, debe aportar suficiente información para determinar de manera inequívoca la cantidad de interés. Todos los algoritmos que se presentarán más adelante tienen ciertos requisitos mínimos sobre la duración de la serie temporal, su precisión y la frecuencia con que deben realizarse las mediciones, de modo que el resultado sea significativo. Si esto no puede lograrse, debe considerarse un subfenómeno más reducido para su estudio.

\newpage

\subsubsection*{Ejemplo 2.1}

Supongamos que deseas estudiar cómo cambia la temperatura en el techo de tu casa. ¿Cuántas mediciones individuales necesitarás? La respuesta es: depende. Si lo que quieres es estudiar lo que ocurre en el transcurso de un día, puedes subir al techo para medir una vez cada hora durante una semana, y así obtener un promedio razonable del perfil de temperatura. Si te interesan las fluctuaciones estacionales, quizá subas al techo solo una vez al día, siempre a la misma hora, pero lo harás durante todo un año. Si lo que quieres es hacer una afirmación sobre cómo fluctúa la temperatura de un año a otro, tendrás que realizar un par de mediciones cada año. Y, finalmente, si tu objetivo es una descripción completa del fenómeno, necesitarás subir al techo una vez por hora todos los días del año durante un par de años. Nótese que el requisito de que el fenómeno bajo estudio deba ser muestreado lo suficiente no es exclusivo del análisis de series temporales, sino que es de carácter general.  

\bigskip

De manera más formal, una señal se denomina estacionaria si todas las probabilidades conjuntas de encontrar al sistema en un estado en cierto momento y más tarde en otro estado son independientes del tiempo \emph{dentro del periodo de observación}, es decir, cuando se calculan a partir de los datos. Esto incluye la constancia de los parámetros relevantes, pero también exige que los fenómenos pertenecientes a la dinámica estén contenidos en la serie temporal con suficiente frecuencia, de manera que las probabilidades u otras reglas puedan inferirse adecuadamente. Si la señal observada es bastante regular la mayor parte del tiempo, pero contiene algún evento muy irregular de forma ocasional, entonces la serie temporal debe considerarse no estacionaria para nuestros propósitos, incluso si todos los parámetros relevantes permanecieron constantes: la señal es intermitente. Solo si los eventos raros (por ejemplo, los brotes irregulares mencionados antes) también aparecen varias veces en la serie temporal, podemos hablar de una independencia efectiva de las probabilidades conjuntas observadas y, por tanto, de estacionariedad.

\bigskip

Las señales no estacionarias son muy comunes, en particular al observar fenómenos naturales o culturales. Tales señales pueden ser valiosas de estudiar e, incluso, se puede ganar dinero prediciendo los precios de las acciones en situaciones que no tienen paralelo en la historia. En este libro adoptamos el punto de vista de que una serie temporal es solo una forma especial de medición científica (es decir, reproducible), la cual se produce para mejorar nuestro conocimiento sobre cierto fenómeno. Ten en cuenta que casi todos los métodos y resultados del análisis de series temporales que se presentarán más adelante asumen la validez de dos condiciones: que los parámetros del sistema permanecen constantes y que el fenómeno se encuentra suficientemente muestreado. Como discutiremos en el Capítulo~13, el requisito de estacionariedad es una consecuencia del hecho de que intentamos aproximarnos a un fenómeno dinámico a través de una serie temporal finita, y por ello constituye un requisito para casi todas las herramientas estadísticas de análisis de series temporales, incluidas las lineales. 

Claramente, dado que los métodos de análisis de series temporales en última instancia se reducen a algoritmos que comprimen los datos en un conjunto de pocos números, se pueden aplicar a cualquier secuencia de datos, incluso a datos no estacionarios. Sin embargo, los resultados no pueden asumirse como caracterización del sistema subyacente. Por tanto, desaconsejamos firmemente ignorar una no estacionariedad conocida. Una salida posible es la segmentación de la serie temporal en segmentos casi estacionarios (véase más adelante para métodos). Somos muy conscientes del problema de que muchas (quizá las más) señales interesantes son no estacionarias. Por ello, en el Capítulo~13 volveremos sobre el problema de la no estacionariedad a un nivel más avanzado y propondremos extensiones del concepto de reconstrucción del espacio de fases para procesos deterministas y estocásticos con parámetros que derivan lentamente.

Debemos mencionar que no consideramos la noción de \emph{estacionariedad débil}, la cual puede encontrarse en la literatura sobre análisis lineal de series temporales. La estacionariedad débil solo requiere que cantidades estadísticas hasta segundo orden sean constantes. En un contexto no lineal, esto es claramente inadecuado.

\newpage

\section{Pruebas de estacionariedad}

Después de haber puesto tanto énfasis en el problema de la estacionariedad, debemos abordar la cuestión de cómo puede detectarse la no estacionariedad en un conjunto de datos dado. Obviamente, la estacionariedad es una propiedad que nunca puede establecerse de manera completamente positiva. Resulta que el asunto es aún más complicado, ya que el requisito de estacionariedad difiere según la aplicación. Por lo tanto, en esta sección solo daremos algunos conceptos generales y discutiremos técnicas particulares en los lugares apropiados dentro del contexto de las aplicaciones. En el Capítulo~13 discutiremos algunas formas de tratar con series temporales provenientes de procesos no estacionarios específicos. Véase también la sección de “Lecturas adicionales” al final de este capítulo para referencias.

Como primer requisito, la serie temporal debe cubrir un intervalo de tiempo mucho mayor que la escala de tiempo característica más larga que sea relevante para la evolución del sistema.

\bigskip

Por ejemplo, la concentración de azúcar en la sangre de un ser humano está impulsada por el consumo de alimentos y, por tanto, sigue aproximadamente un ciclo de 24 horas. Si esta cantidad se registra durante 24 horas o menos, el proceso debe considerarse no estacionario, sin importar cuántos puntos de datos se hayan tomado durante ese tiempo. Información cuantitativa puede obtenerse a partir del espectro de potencia, que se introducirá más adelante. La escala de tiempo relevante más larga puede estimarse como el inverso de la frecuencia más baja que aún contiene una fracción significativa de la potencia total de la señal. Una serie temporal solo puede considerarse estacionaria en escalas de tiempo mucho mayores. Véase también la Sección~2.3.1.

Violaciones fuertes del requisito básico de que las propiedades dinámicas del sistema subyacente a una señal no deben cambiar durante el período de observación pueden comprobarse simplemente midiendo tales propiedades en varios segmentos de los datos. Probabilidades de transición, correlaciones, etc., calculadas, por ejemplo, para la primera y segunda mitad de los datos disponibles, no deben diferir más allá de sus fluctuaciones estadísticas. Para este propósito se prefieren características con fluctuaciones estadísticas conocidas o despreciables. Las cantidades estadísticamente más estables son la media y la varianza. Cantidades más sutiles, como los componentes espectrales, correlaciones o estadísticas no lineales, pueden ser necesarias para detectar formas menos evidentes de no estacionariedad.

\subsubsection*{Ejemplo 2.2 (Varianza móvil de datos de voz)}

Realizamos una prueba muy simple de estacionariedad en la fonación de un paciente que sufría de \emph{laringitis aguda}. Calculamos la media y la varianza muestral para 50 segmentos consecutivos del conjunto de datos, cada uno con 2048 puntos de datos. Dentro de cada segmento, el error estándar de la media estimada $\langle s \rangle = \sum_{n=1}^N s_n / N$ está dado por\footnote{La desviación estándar dividida por un $\sqrt{N}$ adicional se conoce como el error al estimar el valor medio de números no correlacionados distribuidos gausianamente, mientras que la desviación estándar en sí misma describe la dispersión típica de los valores individuales respecto a su media.}

\begin{equation}
	\sqrt{ \frac{\sum_{n=1}^N (s - \langle s \rangle)^2}{N(N-1)} }
\end{equation}


omitiendo el hecho de que las muestras en la serie temporal no son independientes. Las fluctuaciones observadas de la media móvil de los 50 segmentos estaban dentro de estos errores. Sin embargo, la varianza móvil mostró variabilidad más allá de las fluctuaciones estadísticas, como puede verse en la Fig.~2.1. En general, es casi imposible obtener registros estacionarios de seres vivos. Aunque se debería tratar de mantener las condiciones experimentales lo más constantes posible, los cambios temporales en la dinámica espontánea son inevitables y, después de todo, muy naturales. Esta grabación se realizó con el fin de estudiar trastornos del sistema fonatorio. Dado que en este contexto la amplitud (instantánea) no es la cantidad más interesante, uno podría decidir que un cambio lento no tiene mayor importancia. Es entonces necesario utilizar métodos de análisis de series temporales que sean robustos frente a cambios en la amplitud, como el espectro de potencia. 

Sin embargo, puede ser preferible realizar un análisis dependiente del tiempo para evitar artefactos y aprovechar la información sobre el sistema contenida en estos cambios. Más adelante en este capítulo retomaremos este ejemplo y calcularemos un \emph{espectrograma}, es decir, un \emph{espectro de potencia móvil}, que proporciona información mucho más interesante sobre los cambios que ocurren en las frecuencias dominantes.

En sistemas caóticos experimentales, no es raro que una deriva de un parámetro no produzca un cambio visible en la media o en la distribución de valores. Las correlaciones lineales y el espectro también pueden permanecer sin alteraciones. Solo las relaciones dinámicas no lineales y las probabilidades de transición cambian apreciablemente. Esto puede detectarse comparando los errores de predicción con respecto a un modelo no lineal en diferentes secciones del conjunto de datos. Esta posibilidad se discute en el contexto de predicciones no lineales en la Sección~4.4.

La cuestión de si un conjunto de datos es una muestra suficiente para una aplicación particular, como la estimación de una cantidad característica, puede comprobarse observando la convergencia de esa cantidad cuando se utilizan fracciones cada vez mayores de los datos disponibles en el cálculo. Una dimensión de atractor obtenida a partir de la primera mitad de los datos no debe diferir sustancialmente del valor determinado usando la segunda mitad, y debería coincidir con el valor calculado para todo el conjunto de datos dentro de los errores estadísticos estimados. Un ejemplo de esta estrategia se presenta en el Ejemplo~7.6 de la Sección~7.2. Nótese que esta prueba es bastante burda, ya que la convergencia de estadísticas no lineales puede ser muy lenta y, de hecho, se sabe poco sobre su velocidad.

Una cantidad que sufre considerablemente por la no estacionariedad en los datos es la \emph{dimensión de correlación} (que se introducirá en el Capítulo~6). Mientras que una simple deriva de calibración suele aumentar la dimensión (pues elimina cualquier estructura fractal), la mayoría de otros tipos de no estacionariedad y de muestreo insuficiente producen estimaciones de dimensión espuriamente bajas. Por ejemplo, si el proceso no es recurrente, como ciertos procesos de ruido coloreado con un espectro $f^{-\alpha}$ con $1 < \alpha < 3$ (el espectro de potencia en función de la frecuencia diverge como una ley de potencia a bajas frecuencias), la dimensión debería ser infinita, mientras que los algoritmos estándar de dimensión arrojan valores finitos. Los detalles se discutirán en la Sección~6.5, donde introduciremos el \emph{diagrama de separación temporal} \textit{Provenzale1992}, que es una prueba de estacionariedad particularmente útil para estimar la suma de correlación. Sin duda, esta prueba también puede usarse en aplicaciones más generales para distinguir correlaciones geométricas de temporales.

\newpage

\section{Correlaciones lineales y el espectro de potencia}

La inferencia estadística lineal es un campo extremadamente bien desarrollado. Además de la definición de varias herramientas estadísticas para caracterizar los datos, su principal valor reside en el hecho de que los conceptos pueden derivarse de una manera rigurosa. Por ejemplo, es relevante distinguir entre una cantidad como un valor medio y la forma en que se obtiene un número para representarlo a partir de una muestra finita. A esto se le llama una \emph{estimación}, y puede haber diferentes estimaciones de la misma cantidad en el mismo conjunto de datos, dependiendo de los supuestos realizados. Así, distinguiremos las estimaciones de cualquier cantidad colocando un “sombrero” sobre el símbolo. A diferencia de la teoría lineal, las estimaciones de cantidades no lineales que haremos más adelante usualmente no pueden estudiarse lo suficiente como para demostrar rigurosamente su corrección y convergencia. En este libro introduciremos únicamente los conceptos de procesos estadísticos lineales y del espectro de potencia y funciones de correlación; cualquier otro aspecto de los métodos lineales de series temporales no será considerado explícitamente.

Las soluciones de ecuaciones de movimiento deterministas lineales estables con coeficientes constantes son superposiciones lineales de oscilaciones armónicas amortiguadas exponencialmente (la estabilidad aquí implica que las soluciones permanecen acotadas; de lo contrario, también son posibles soluciones en crecimiento exponencial). Por lo tanto, los sistemas lineales siempre necesitan entradas irregulares para producir señales irregulares. El sistema más simple que produce señales no periódicas es un \emph{proceso estocástico lineal}. Por esta razón, supondremos momentáneamente que el observable es una variable aleatoria. Sin embargo, este punto de vista también será razonable si hablamos de sistemas deterministas, por varias razones. En primer lugar, el ruido de medición siempre introduce algo de aleatoriedad. En segundo lugar, la condición inicial desconocida del sistema convierte la observación actual en un evento aleatorio. Finalmente, la inestabilidad de los sistemas caóticos sugiere una descripción probabilística en términos de medidas invariantes (teoría ergódica). Así, una medición $s_n$ del estado en el tiempo $n$ de un proceso estocástico puede considerarse como tomada de una distribución de probabilidad subyacente para observar diferentes valores o secuencias de valores.

Un proceso estocástico se denomina \emph{estacionario} si estas probabilidades son constantes en el tiempo: $p(s_n) = p(s)$

Dado que no se conoce de antemano de qué distribución de probabilidad $p(s)$ se extraen las mediciones, la información sobre $p(s)$ solo puede inferirse a partir de la serie temporal. La \emph{media} de la distribución de probabilidad $p(s)$, definida como

\begin{equation}
\langle s \rangle := \int_{-\infty}^{\infty} ds' \, s' \, p(s'), 
\end{equation}

puede estimarse mediante la media de la serie temporal finita

\begin{equation}
\langle \hat{s} \rangle_N = \frac{1}{N} \sum_{n=1}^N s_n ,
\end{equation}

donde $\langle \cdot \rangle_n$ denota el promedio en el tiempo y $N$ es el número total de mediciones en la serie temporal. La \emph{varianza} de la distribución de probabilidad se estimará mediante la varianza de la serie temporal:

\begin{equation}
\sigma^2 := \langle (s - \langle s \rangle)^2 \rangle = \int_{-\infty}^{\infty} ds' \, (s' - \langle s \rangle)^2 p(s'),
\end{equation}

\begin{equation}
\hat{\sigma}^2 = \frac{1}{N-1} \sum_{n=1}^N (s_n - \langle \hat{s} \rangle_N)^2 
= \frac{1}{N-1} \left( \sum_{n=1}^N s_n^2 - N (\langle \hat{s} \rangle_N)^2 \right).
\end{equation}

Nos referiremos a $\sigma$ como la desviación estándar o la amplitud rms (\emph{root mean squared}), en contraste con la \emph{varianza} $\sigma^2$.

Para estas dos cantidades, el orden temporal de las mediciones es irrelevante y, por lo tanto, no pueden aportar información sobre la evolución temporal de un sistema. Dicha información puede obtenerse a partir de las \emph{autocorrelaciones} de una señal. La autocorrelación con retardo $\tau$ está dada por

\begin{equation}
c_{\tau} = \frac{1}{\sigma^2} \langle (s_n - \langle s \rangle)(s_{n-\tau} - \langle s \rangle) \rangle 
= \frac{\langle s_n s_{n-\tau} \rangle - \langle s \rangle^2}{\sigma^2}.
\tag{2.5}
\end{equation}

La estimación de las autocorrelaciones a partir de una serie temporal es sencilla siempre que el retardo $\tau$ sea pequeño en comparación con la longitud total de la serie. Por lo tanto, las estimaciones $\hat{c}_\tau$ son razonables únicamente para $\tau \ll N$.  

Si graficamos los valores $s_n$ frente a los correspondientes valores con un retardo fijo $\tau$, es decir, $s_{n-\tau}$, la autocorrelación $c_\tau$ cuantifica cómo se distribuyen estos puntos. Si se dispersan uniformemente en el plano, entonces $c_\tau = 0$. Si tienden a agruparse a lo largo de la diagonal $s_n = s_{n-\tau}$, entonces $c_\tau > 0$, y si están más cerca de la recta $s_n = -s_{n-\tau}$ tenemos $c_\tau < 0$. Estos dos últimos casos reflejan cierta tendencia de $s_n$ y $s_{n-\tau}$ a ser proporcionales entre sí, lo que hace plausible que la función de autocorrelación refleje únicamente correlaciones lineales.

Si la señal se observa en tiempo continuo, se puede introducir la \emph{función de autocorrelación} $c(\tau)$, y las correlaciones de la Ec.~(2.5) son estimaciones de $c(\tau = \tau \Delta t)$. Obviamente, $c_{\tau} = c_{-\tau}$ y $c_0 = 1$.

Un caso particular se presenta si las mediciones provienen de una distribución gaussiana:

\begin{equation}
p(s) = \frac{1}{\sigma \sqrt{2\pi}} \, e^{-(s - \langle s \rangle)^2 / 2\sigma^2},
\tag{2.6}
\end{equation}

y, además, las distribuciones conjuntas de múltiples mediciones también son gaussianas. Como consecuencia, sus segundos momentos están dados por las autocorrelaciones. Dado que una distribución gaussiana multivariada está completamente especificada por sus primeros y segundos momentos, tal proceso aleatorio se describe estadísticamente de manera completa mediante la media, la varianza $\sigma^2$ y la función de autocorrelación $c_\tau$, y se denomina un \emph{proceso gaussiano lineal aleatorio}.

Obviamente, si una señal es periódica en el tiempo, entonces la función de autocorrelación es periódica en el retardo $\tau$. Los procesos estocásticos tienen autocorrelaciones decrecientes, pero la tasa de decaimiento depende de las propiedades del proceso. Las autocorrelaciones de señales provenientes de sistemas caóticos deterministas típicamente también decaen de forma exponencial con el incremento del retardo. Las autocorrelaciones no son lo suficientemente características como para distinguir señales aleatorias de señales caóticas deterministas.

\bigskip

En lugar de describir las propiedades estadísticas de una señal en el espacio real, se pueden estudiar sus propiedades en el espacio de Fourier. La \emph{transformada de Fourier} establece una correspondencia uno a uno entre la señal en ciertos instantes de tiempo (\emph{dominio temporal}) y cómo ciertas frecuencias contribuyen a la señal, así como la forma en que las fases de las oscilaciones se relacionan con las fases de otras oscilaciones (\emph{dominio frecuencial}). Aparte de diferentes convenciones respecto a la normalización, la transformada de Fourier de una función $s(t)$ está dada por

\begin{equation}
\tilde{s}(f) := \frac{1}{\sqrt{2\pi}} \int_{-\infty}^{\infty} s(t) e^{2\pi i f t} dt,
\tag{2.7}
\end{equation}

y la de una serie temporal discreta y finita por

\begin{equation}
\tilde{s}_k = \frac{1}{\sqrt{N}} \sum_{n=1}^N s_n e^{2\pi i k n / N}.
\tag{2.8}
\end{equation}

Aquí, las frecuencias en unidades físicas son $f_k = k / (N \Delta t)$, donde $k = -N/2, \ldots, N/2$ y $\Delta t$ es el intervalo de muestreo. La normalización de ambas transformadas es tal que la transformada inversa solo implica un cambio de signo en el exponente y un intercambio de $f$ (o, respectivamente, $k$) y $t$ (respectivamente $n$). Ambas operaciones son por tanto invertibles y constituyen transformaciones lineales.

El \emph{espectro de potencia} de un proceso se define como el módulo al cuadrado de la transformada de Fourier continua, $S(f) = |\tilde{s}(f)|^2$. Es el cuadrado de la amplitud, por el cual la frecuencia $f$ contribuye a la señal. Si queremos estimarlo a partir de una serie temporal discreta, podemos usar el \emph{periodograma}, $S_k = |\tilde{s}_k|^2$. Este no es el estimador más útil para $S(f)$ por dos razones. Primero, la resolución de frecuencia finita de la transformada de Fourier discreta conduce a \emph{fugas} (\emph{leakage}) hacia bins de frecuencia adyacentes. Por ejemplo, en el caso de una señal armónica pura, donde la transformada continua arroja una $\delta$-distribución, la transformada discreta solo produce un pico finito que se ensancha para longitudes de serie finitas. Segundo, sus fluctuaciones estadísticas son del mismo orden que $S(f)$ mismo. Esto puede remediarse promediando sobre bins de frecuencia adyacentes o, más eficientemente, promediando sobre ventanas en el dominio temporal.  

Si no se utiliza un paquete para la estimación del espectro de potencia, se recomienda consultar el \emph{Numerical Recipes} de Press et al.~(1992), o textos similares para más detalles sobre estimación numérica del espectro de potencia y técnicas de ventana (y rutinas). El paquete de software TISEAN contiene una rutina relativamente sencilla para la estimación espectral, además de otras herramientas lineales. Para series temporales largas de tamaño finito, puede usarse software especializado como el \emph{toolbox} de procesamiento de señales de Matlab.

La relación importante entre el espectro de potencia y la función de autocorrelación es conocida como el \emph{teorema de Wiener–Khinchin} y establece que la transformada de Fourier de la función de autocorrelación es igual al espectro de potencia. La potencia total crece linealmente con la longitud de la serie temporal y puede calcularse tanto en el dominio temporal como en el frecuencial, mediante el \textit{teorema de Parseval}
\begin{equation}
S = \sum_{n=1}^N |s_n|^2 = \sum_{k=-N/2}^{N/2} |\tilde{s}_k|^2 
\end{equation}



El espectro de potencia es particularmente útil para estudiar la oscilación de un sistema. Habrá picos más nítidos o más anchos en las frecuencias dominantes y en sus múltiplos enteros, los armónicos. El ruido de medición añade un piso continuo al espectro: cuanto más intensa la señal, más visibles son las líneas espectrales. Los sistemas deterministas caóticos también pueden mostrar líneas espectrales nítidas incluso en ausencia de un piso continuo en el espectro. Esto es una consecuencia inmediata de una función de autocorrelación exponencialmente decreciente. Sin información adicional, no es posible decidir a partir del espectro si la señal continua proviene de ruido puro o de un sistema (casi) periódico o caótico. Un ejemplo es el mapa de Golub \& Swinney (1975), que muestra cómo la amplitud continua en el espectro de potencia aumenta a medida que la dinámica cambia de cuasi-periódica a caótica.

\subsubsection*{Ejemplo 2.3 (Láser de realimentación no lineal)}
Veamos un espectro típico de potencia de un sistema caótico continuo en el tiempo, el láser óptico de realimentación no lineal descrito en el Apéndice B.10. Estimamos el espectro de potencia de los 128 000 puntos de datos deslizando una ventana cuadrada de longitud del 40\% sobre los datos y promediando el módulo al cuadrado de las transformadas de Fourier de todas las ventanas (Fig.~2.2).

Cuando estudiamos una señal en el dominio de la frecuencia promediando en el tiempo, como en las autocorrelaciones, perdemos toda la información temporal. A veces, en particular cuando las propiedades del sistema pueden cambiar con el tiempo (lo cual conduce a señales no estacionarias), queremos preservar esta información. En este caso, el análisis espectral puede realizarse en segmentos consecutivos de una serie temporal más larga para rastrear cambios temporales. Esto se denomina usualmente un \emph{espectrograma}.

\subsubsection*{Ejemplo 2.4 (Espectrograma de datos de voz)}
Usemos el espectro de potencia para el análisis de la fonación de un paciente (Apéndice B.8) que en el Ejemplo 2.2 se mostró que es no estacionaria. Primero, calculamos un espectro de potencia promedio de todo el conjunto de datos. Aparecen picos más nítidos en frecuencias discretas, en su mayoría múltiplos enteros: este fenómeno se denomina \emph{bifonación}; véase Herzel \& Reuter (1996).  

Luego calculamos espectros de potencia para 100 segmentos de 2048 puntos de datos cada uno. El traslape de segmentos consecutivos es ahora de 1024 puntos. Tal espectro dependiente del tiempo se denomina habitualmente \emph{espectrograma} y se utiliza de manera rutinaria en la investigación de la voz. En la Fig.~2.3, para cada uno de los segmentos colocamos un símbolo en las frecuencias donde la potencia alcanza un valor relativo máximo. A la izquierda también mostramos las frecuencias que contribuyen al espectro de potencia total, promediado en toda la serie. Además, trazamos líneas de puntos en las dos frecuencias fundamentales, 306 Hz y 426 Hz. Hacia el final del registro aparecen subarmónicos de la mitad o un tercio de la frecuencia fundamental más baja.

\newpage

\subsubsection{Estacionariedad y el componente de baja frecuencia en el espectro de potencia}

Como mencionamos en la Sección 2.1, la no estacionariedad efectiva de una serie temporal finita está relacionada con el hecho de que las escalas de tiempo de la dinámica son del orden del tiempo total de observación. El espectro de potencia refleja la contribución de todos los posibles periodos, desde el intervalo de muestreo hasta el tiempo total cubierto por la serie, y por lo tanto debería contener la información relevante.  

Si hay demasiada potencia en las bajas frecuencias, la serie temporal debe considerarse no estacionaria, ya que los modos de Fourier correspondientes solo tienen muy pocas oscilaciones durante el tiempo de observación. Lo que no está del todo claro es qué significa exactamente “demasiada” potencia. Recordemos que el ruido blanco (o mejor dicho, en tiempo discreto, un proceso \emph{iid}: números aleatorios independientes e idénticamente distribuidos), que es un prototipo de proceso estacionario, tiene potencia igual en todos los bins de frecuencia.  

\bigskip

Además, la potencia para la frecuencia $f=0$ puede modificarse arbitrariamente mediante un simple desplazamiento en los datos, y está además relacionada con la integral sobre la función de autocorrelación (una consecuencia del teorema de Wiener–Khinchin). Por lo tanto, debemos esperar un problema real solo si el espectro de potencia diverge alrededor de cero, como $f^{-\alpha}$ con $\alpha$ positivo. En particular, el movimiento browniano tiene un espectro $f^{-2}$ y se sabe que no es recurrente en dos o más dimensiones.  

La decisión sobre si los datos son estacionarios o no es mucho más difícil si el espectro de potencia se aproxima a una constante distinta de cero para frecuencias pequeñas, lo cual es mucho mayor que el fondo. Solo podemos sugerir que esto se interprete como una advertencia.

\newpage

\section{Filtros lineales}

Aunque la teoría de la Sección 2.3 fue desarrollada a partir de procesos estocásticos lineales —los cuales generan oscilaciones amortiguadas irregularmente distorsionadas—, aquí reinterpretaremos algunos aspectos del tratamiento de señales (casi) periódicas que son distorsionadas por ruido de medición. Es decir, para una situación donde el espectro de potencia de la señal no perturbada consiste en picos pronunciados en ciertas frecuencias. Mostraremos cómo una \emph{señal} y un \emph{ruido aleatorio} pueden distinguirse sobre la base del espectro de potencia.  

Debido a la linealidad de la transformada de Fourier, el espectro de potencia $S_k^{(s)}$ de una superposición aditiva de una señal $x$ y un ruido $\eta$ es la superposición del espectro de la señal $S_k^{(x)}$ y del ruido $S_k^{(\eta)}$:

\begin{equation}
S_k^{(s)} = S_k^{(x)} + S_k^{(\eta)}.
\end{equation}

Esto puede transformarse fácilmente en un filtro en el dominio de Fourier. Es decir, asumimos que la secuencia observada también es tal superposición, $s_n = x_n + \eta_n$, y queremos estimar la señal pura, $\hat{x}_n \approx x_n$, filtrando la secuencia medida de manera apropiada: $\hat{x}_k = \phi_k \tilde{s}_k$.  

Si requerimos que la distancia cuadrática media $e^2$ entre $\hat{x}_n$ y $x_n$ sea mínima, el \emph{filtro óptimo} puede construirse resolviendo un problema de minimización. En el dominio de Fourier debemos minimizar:

\begin{equation}
e^2 = \sum_{k=-N/2}^{N/2} |\hat{x}_k - \tilde{x}_k|^2
= \sum_{k=-N/2}^{N/2} |\phi_k(\tilde{x}_k + \tilde{\eta}_k) - \tilde{x}_k|^2
= \sum_{k=-N/2}^{N/2} \Big[ (\phi_k - 1)^2 |\tilde{x}_k|^2 + \phi_k^2 |\tilde{\eta}_k|^2 \Big]
\end{equation}

(asumiendo independencia estadística de $x$ y $\eta$), lo cual puede hacerse anulando la derivada respecto a $\phi_k$. La solución es

\begin{equation}
\phi_k = \frac{S_k^{(x)}}{S_k^{(s)}}
\end{equation}

Esto se denomina el \emph{filtro de Wiener} (\emph{Wiener filter}). Obviamente, no podemos calcular $S_k^{(x)}$ directamente, pero podemos intentar estimarlo mediante la inspección del espectro de potencia total.  

Aquí es donde necesitamos un criterio para distinguir la señal del ruido a partir de sus propiedades respectivas. Como hemos mencionado, esto es posible con frecuencia para señales con oscilaciones dominantes. Sin embargo, para señales con espectros de banda ancha intrínsecos, como aquellas provenientes de sistemas caóticos deterministas, esto suele ser imposible.

\newpage

\subsubsection*{Ejemplo 2.5 (Fracaso en la reducción de ruido en electrocardiogramas usando un filtro de Wiener)}

La Fig.~2.4 muestra un electrocardiograma (ECG) humano relativamente ruidoso. Los paneles superior e inferior muestran las gráficas antes y después del filtrado de Wiener, respectivamente. (Véase el Apéndice B.7 para más información sobre los datos). $S_k^{(x)}$ ha sido estimado mediante $S_k^{(s)} - \bar{S}^{(\eta)}$, donde la potencia del ruido $S_k^{(\eta)}$ se ha supuesto independiente de $k$ (ruido blanco) e idéntica al piso de ruido $\bar{S}^{(\eta)}$, que persiste hasta altas frecuencias. Observamos que el filtro de Wiener es incapaz de limpiar el ECG, lo cual está relacionado con el hecho de que el espectro de potencia del ECG libre de ruido también es de banda ancha debido a los picos pronunciados. Véase Schreiber \& Kaplan (1996a, b) y la Sección 10.4 para un método de filtrado de ECGs con técnicas no lineales.

\bigskip

Existen situaciones, sin embargo, en las que contamos con información externa acerca de la contribución de la señal o del ruido al espectro de potencia. Un ejemplo es la contaminación por corriente alterna de 50 Hz (o 60 Hz en EE.\,UU.) en señales registradas con circuitos electrónicos.  

En este caso sabemos que la mayor parte de la potencia en 50 Hz se debe a esta contaminación y podemos construir un filtro interpolando entre los bins de frecuencia adyacentes no afectados.

\subsubsection*{Ejemplo 2.6 (Eliminación de la contaminación de 50 Hz AC con un filtro de Wiener)}

Un ejemplo en el que los filtros lineales son completamente adecuados incluso para señales no lineales es el ruido de 50 Hz (o 60 Hz) en la salida de amplificadores eléctricos. La contaminación está bien localizada en el dominio de Fourier y el espectro de la señal puede estimarse fácilmente interpolando sobre el pico espurio en el espectro de potencia.   La Fig.~2.5 muestra un electrocardiograma humano (ECG) antes y después del filtrado. (Véase el Apéndice B.7 para más información sobre los datos). Un filtro que elimina esencialmente un solo componente de frecuencia se denomina en la literatura \emph{filtro notch}. Para detalles técnicos puede consultarse la literatura estándar, por ejemplo, Hamming (1983).

\bigskip

Vale la pena mencionar que ciertos filtros lineales, que a veces están incorporados en el hardware de adquisición de datos, pueden añadir un grado de libertad determinista y, por lo tanto, aumentar la dimensión del atractor. La razón es que, para el suavizado de datos en línea, se puede utilizar un filtro con realimentación que introduce otro grado de libertad en el sistema. Véanse Badii \& Politi (1986), Badii \emph{et al.} (1988), Mitschke \emph{et al.} (1988), Mitschke (1990), Broomhead \emph{et al.} (1992), Sauer \& Yorke (1993), Theiler \& Eubank (1993), y la nota al pie en la página 231.

\newpage

\section{Predicciones lineales}

Estrechamente relacionado con el filtrado óptimo está el problema de la predicción. Para esta discusión será más natural trabajar en el dominio del tiempo. Supongamos que tenemos una secuencia de mediciones $s_n$, $n = 1, \ldots, N$, y queremos predecir el resultado de la siguiente medición, $s_{N+1}$.  A menudo se desea encontrar la predicción $\hat{s}_{N+1}$ que minimiza el valor esperado del error cuadrático de predicción $\langle (\hat{s}_{N+1} - s_{N+1})^2 \rangle$.  

\bigskip

Cuando asumimos estacionariedad, podemos estimar este valor esperado mediante el promedio sobre los valores medidos disponibles. Si restringimos aún más la minimización a modelos lineales de series temporales que incorporan las últimas $m$ mediciones, podemos expresar esto como:

\begin{equation}
\hat{s}_{N+1} = \sum_{j=1}^{m} a_j s_{n-m+j}
\end{equation}

y minimizar

\begin{equation}
\sum_{n=m}^{N-1} \left( \hat{s}_{n+1} - s_{n+1} \right)^2
\end{equation}

respecto a los parámetros $a_j$, $j=1,\ldots,m$.  

\bigskip

Aquí hemos supuesto que la media de la serie temporal ya ha sido restada de las mediciones. Exigiendo que las derivadas respecto a todos los $a_j$ sean cero, obtenemos la solución resolviendo el siguiente sistema lineal de ecuaciones:

\begin{equation}
\sum_{j=1}^{m} C_{ij} a_j = \sum_{n=m}^{N-1} s_{n+1} s_{n-m+i}, 
\quad i = 1, \ldots, m.
\end{equation}

Aquí $C_{ij}$ es la matriz de autocovarianza de tamaño $m \times m$:

\begin{equation}
C_{ij} = \sum_{n=m}^{N-1} s_{n-m+i} s_{n-m+j}.
\end{equation}

La relación lineal, Ec.~(2.12), está justificada tanto para el movimiento armónico como para procesos estocásticos lineales. Aunque pospondremos la discusión de los distintos modelos lineales de series temporales a la Sección 12.1.1, queremos señalar aquí que todos estos modelos requieren entradas aleatorias que no observamos. Así, la llamada \textit{media móvil} de estos modelos, que describe la respuesta al ruido externo, debe ser reemplazada aquí por la \textit{respuesta media}, la cual se asume nula ya que el ruido tiene media cero. Además, nótese que el modelo de la Ec.~(2.12) no puede esperarse que sea estable bajo iteración. Realizar predicciones a más pasos mediante iteración significaría usar los coeficientes $a_j$, que han sido optimizados para predicciones basadas en las mediciones (ruidosas), en predicciones basadas en los mejores estimadores $\hat{s}_n$. Sus errores poseen propiedades diferentes a las del ruido $\eta_n$. La manera más sencilla de evitar este problema es modificar la Ec.~(2.12) para hacer predicciones más allá del futuro inmediato en una única iteración.  

\bigskip

Un ejemplo de predicción lineal de series temporales se da en el Ejemplo 1.1, donde se compara también con un método no lineal sencillo. Un segundo ejemplo será presentado en el Capítulo 4, Ejemplo 4.2. Es siempre importante reportar los logros de un nuevo enfoque en comparación con los métodos estándar. Así, si logras predecir algunas series temporales con un algoritmo no lineal, debes asegurarte de que no solo has capturado las autocorrelaciones lineales en los datos, las cuales pueden ser mejor explotadas por un predictor lineal. Todavía no hemos dicho nada sobre cómo elegir el orden correcto $m$ de la relación lineal, Ec.~(2.12). Algunos comentarios generales se darán en la Sección 12.5; sin embargo, remitimos al lector a la literatura especializada sobre análisis lineal de series temporales para afirmaciones más rigurosas. Por ahora, solo sugerimos que aunque el modelo ha sido construido sobre el error promediado sobre los datos disponibles, no debemos citar este error como un error de predicción ya que ha sido determinado \textit{in-sample}. En su lugar, debemos retener parte de los datos mientras construimos el modelo para poder calcular el error relevante \textit{out-of-sample} una vez que se ha establecido un modelo.

\newpage

\subsection*{Lecturas adicionales}

Material práctico sobre análisis lineal de series temporales, filtrado y predicción, incluyendo los programas esenciales, puede encontrarse en el famoso \textit{Numerical Recipes}, Press \textit{et al.} (1992). Existen muchos libros de texto sobre análisis espectral de series temporales; ejemplos clásicos son Box \& Jenkins (1976), Priestley (1988), Jenkins \& Watts (1986) y Brockwell \& Davis (1987). Los filtros digitales se tratan, por ejemplo, en Hamming (1983). Las pruebas clásicas de estacionariedad se describen en Priestley (1992). El Capítulo 13 aborda series temporales no estacionarias y ofrece más referencias. Otras referencias incluyen Isliker \& Kurths (1993), Kennel (1997), así como Schreiber (1997).

\newpage

\section{Ejercicios}

\begin{enumerate}
    \item[2.1] Crea una muestra $\{x_n\}$ de números aleatorios gaussianos no correlacionados. Puedes usar una rutina de algún paquete de software [por ejemplo, \textit{Numerical Recipes}, Press \textit{et al.} (1992)] o la rutina \texttt{makenoise} de TISEAN. Ahora aplica el \textit{filtro de media móvil}
    \[
    s_n = \frac{1}{15} \sum_{j=-7}^{7} x_{n+j}
    \]
    para obtener 1024 variables gaussianas correlacionadas. Obtén una rutina de algún paquete de software (por ejemplo, TISEAN) que estime el espectro de potencia de un conjunto de datos. Aplica la rutina a las secuencias $\{x_n\}$ y $\{s_n\}$. ¿Qué diferencia observas?
	\newpage
    \item[2.2] Crea dos series temporales artificiales. La primera, $\{\eta_n, \; n = 1, \dots, 4096\}$, contiene números aleatorios distribuidos uniformemente (usa tu generador favorito de números aleatorios) en el intervalo $[0,1]$. La segunda serie, $\{s_n, \; n = 1, \dots, 4096\}$, se basa en la evolución determinista de $x_n$, que sigue las reglas $x_0 = 0.1$ y $x_{n+1} = 1 - 2x_n^2$ (el \textit{mapa de Ulam}, un caso especial de la \textit{ecuación logística}). Los valores $x_n$ no se miden directamente, sino a través de la función de observación no lineal
    \[
    s_n = \arccos(-x_n)/\pi.
    \]
    Compara la media, la varianza y los espectros de potencia de las dos series temporales. 
\end{enumerate}

\newpage

\chapter{Métodos de espacio de fases}
\section{Determinismo: unicidad en el espacio de fases}

Los métodos de series temporales no lineales discutidos en este libro están motivados y basados en la teoría de los \textit{sistemas dinámicos}; es decir, la evolución temporal se define en algún espacio de fases. Dado que tales sistemas no lineales pueden exhibir caos determinista, este es un punto de partida natural cuando aparece irregularidad en una señal. Eventualmente, uno podría pensar en incorporar también un componente estocástico en la descripción. Sin embargo, hasta ahora debemos asumir que este componente estocástico es pequeño y esencialmente no cambia las propiedades no lineales. Así, todos los enfoques exitosos que conocemos o bien asumen que la no linealidad es una pequeña perturbación de un proceso estocástico lineal esencialmente, o consideran al elemento estocástico como una pequeña contaminación de un proceso no lineal esencialmente determinista. Si se supone que un conjunto de datos proviene de procesos estocásticos genuinamente \textit{no lineales}, las herramientas de análisis de series temporales son todavía muy limitadas y su discusión se pospondrá a la Sección 12.1.

Consideremos por un momento un sistema puramente determinista. Una vez que su estado presente está fijado, los estados en todos los tiempos futuros quedan también determinados. Por tanto, será importante establecer un espacio vectorial (llamado \textit{espacio de estados} o \textit{espacio de fases}) para el sistema, de modo que especificar un punto en este espacio defina el estado del sistema y viceversa. Entonces podemos estudiar la dinámica del sistema estudiando la dinámica de los puntos de fase correspondientes. En teoría, los sistemas dinámicos suelen estar definidos por un conjunto de ecuaciones diferenciales ordinarias de primer orden (véase más adelante) que actúan en un espacio de fases. La teoría matemática de las ecuaciones diferenciales ordinarias asegura la existencia y unicidad de las trayectorias, si se cumplen ciertas condiciones. No obstante, no debemos esperar una distinción académica estricta entre el estado y el espacio de fases, pero sí debemos remarcar que, salvo modelos matemáticos excepcionales, con ecuaciones dadas no siempre habrá una única elección de lo que el espacio de fases de un sistema debe contener.

\bigskip

El concepto de \textit{estado de un sistema} es poderoso incluso para sistemas no deterministas. Una gran clase de sistemas puede describirse por un conjunto (posiblemente infinito) de estados y algún tipo de reglas de transición que especifican cómo el sistema puede pasar de un estado a otro. Los miembros más destacados de esta categoría son los \textit{procesos de Markov}, para los cuales las reglas de transición se dan en forma de un conjunto de probabilidades de transición, y el estado futuro se selecciona aleatoriamente de acuerdo con estas probabilidades. La característica esencial de estos procesos es su estricta memoria finita: las probabilidades de transición a los estados futuros solo pueden depender del estado presente, no del pasado.\footnote{Si se quiere, puede considerarse un sistema puramente determinista como un caso límite de un proceso de Markov en un continuo de estados. La transición al estado especificado por la regla determinista ocurre con probabilidad 1 y todas las demás transiciones tienen probabilidad 0.} Mencionamos este enfoque porque trata las incertidumbres en la regla de transición —lo que llamamos \textit{ruido dinámico}— como el caso genérico de una distribución con picos de las probabilidades de transición. El caso determinista libre de ruido se obtiene en el límite de un pico delta. Esta no es la formulación más útil para el propósito de este libro. En su lugar, formulamos el caso particular de un sistema determinista en sí mismo como un punto de partida que nos permite establecer algunas formas interesantes de entender señales y sistemas. Tengamos en cuenta el hecho de que podemos tratar el ruido dinámico no solo como una complicación adicional en una situación limpia, sino que también podemos considerar al determinismo estricto como un caso límite de una clase muy general de modelos.

\newpage

Introduzcamos ahora algo de notación para sistemas dinámicos deterministas en el espacio de fases. Las modificaciones necesarias para sistemas no exactamente deterministas se discutirán más adelante. Para simplificar, nos restringiremos al caso en el que el espacio de fases es un espacio vectorial finito dimensional $\mathbb{R}^m$ (las ecuaciones en derivadas parciales como la ecuación de Navier–Stokes para el flujo hidrodinámico conducen a sistemas dinámicos altamente interesantes que viven en espacios de fases de dimensión infinita). Un estado está especificado por un vector $\mathbf{x} \in \mathbb{R}^m$. Entonces podemos describir la dinámica ya sea mediante una aplicación $m$-dimensional o mediante un sistema explícito de $m$ ecuaciones diferenciales ordinarias de primer orden. En el primer caso, el tiempo es una variable discreta:  

\begin{equation}
\mathbf{x}_{n+1} = \mathbf{F}(\mathbf{x}_n), \quad n \in \mathbb{Z}, 
\end{equation}

y en el segundo caso es una variable continua:  

\begin{equation}
\frac{d}{dt}\mathbf{x}(t) = \mathbf{f}(\mathbf{x}(t)), \quad t \in \mathbb{R}.
\end{equation}

La segunda situación suele denominarse \textit{flujo}. El campo vectorial $\mathbf{f}$ en la Ecuación (3.2) está definido de modo que no depende explícitamente del tiempo, y por lo tanto se denomina \textit{autónomo}.

\bigskip

Si $\mathbf{f}$ contiene una dependencia explícita del tiempo, por ejemplo, a través de algún término externo de excitación, la literatura matemática ya no considera esto como un sistema dinámico, puesto que la invariancia bajo traslaciones temporales se rompe. El vector de estado por sí solo (es decir, sin la información sobre el tiempo actual $t$) no define la evolución de manera única. En muchos casos, como fuerzas periódicas de excitación, el sistema puede hacerse autónomo mediante la introducción de grados de libertad adicionales (por ejemplo, una excitación sinusoidal puede generarse mediante un oscilador armónico autónomo adicional con un acoplamiento unidireccional). Entonces, típicamente se puede definir un espacio de fases extendido en el cual la evolución temporal vuelve a ser una función única de los vectores de estado, incluso sin introducir grados de libertad auxiliares; basta con introducir, por ejemplo, un ángulo de fase de la fuerza excitadora. 

Sin embargo, al admitir dependencias arbitrarias en el tiempo de $\mathbf{f}$, esto incluye también el caso de un sistema estocástico impulsado por ruido, el cual no es un sistema dinámico en el sentido de una dependencia única del futuro respecto de un vector de estado actual.

En el caso autónomo, se sabe que la solución del problema de valor inicial de la 
Ecuación (3.2) existe y es única si el campo vectorial $\mathbf{f}$ es continuo en el sentido de Lipschitz.  
Una secuencia de puntos $\mathbf{x}_n$ o $\mathbf{x}(t)$, resolviendo las ecuaciones anteriores, se denomina 
\textit{trayectoria} del sistema dinámico, con $\mathbf{x}_0$ o $\mathbf{x}(0)$, respectivamente, 
como la \textit{condición inicial}.  
Las trayectorias típicas o bien divergen hacia el infinito conforme avanza el tiempo, o permanecen en una región acotada para siempre, que es el caso que aquí nos interesa.\footnote{En este libro no discutimos problemas de dispersión, donde las trayectorias se aproximan a una región de interacción que proviene del infinito y, después de dispersarse, desaparecen nuevamente hacia el infinito. Sin embargo, existe una clase de fenómenos muy interesantes llamados \textit{dispersión caótica}. La dinámica determinista pero caótica conduce a una dependencia irregular del ángulo de dispersión respecto al parámetro de impacto.}  
El comportamiento observado depende tanto de la forma de $\mathbf{F}$ (o, respectivamente, $\mathbf{f}$) como de la condición inicial; muchos sistemas permiten ambos tipos de solución.  
El conjunto de condiciones iniciales que conducen al mismo comportamiento asintótico de la trayectoria se denomina \textit{cuenca de atracción} para ese movimiento en particular.  

En muchas ocasiones encontraremos más conveniente la formulación en tiempo discreto.  
La solución formal de la ecuación diferencial Eq.~(3.2) que relaciona una condición inicial con el final de la trayectoria un intervalo de tiempo después se denomina a veces \textit{mapa de un paso} de $\mathbf{f}$.  
Nos referiremos a él de este modo, ya que, después de todo, las series temporales con las que trabajamos sólo se nos entregan en pasos de tiempo discretos.  
Además, la integración numérica de las ecuaciones diferenciales, Eq.~(3.2), con un paso de tiempo finito $\Delta t$ genera un mapa.  
Por ejemplo, el esquema de integración de Euler produce  

\begin{equation}
\mathbf{x}(t+\Delta t) \approx \mathbf{x}(t) + \Delta t \, \mathbf{f}(\mathbf{x}(t)).
\end{equation}

Cuando el paso de tiempo $\Delta t$ es pequeño, la diferencia entre los valores consecutivos $\mathbf{x}(t), \mathbf{x}(t+\Delta t)$ también lo es, lo cual es característico del tipo particular de mapa que surge de las ecuaciones diferenciales.  
Nos referiremos a tales series temporales como “tipo flujo”.  
La diferencia fundamental entre datos tipo flujo y tipo mapa será discutida más adelante y en la Sección~3.5.

Los sistemas dinámicos usados como ejemplos en este libro no solo tienen soluciones acotadas, sino que usualmente también son \textit{disipativos}, lo que significa que, en promedio, un volumen en el espacio de fases que contiene condiciones iniciales se contrae bajo la dinámica. Entonces, tenemos en promedio $|\det J_{\mathbf{F}}| < 1$ o $\operatorname{div} \mathbf{f} < 0$, respectivamente. ($J_{\mathbf{F}}$ es la matriz Jacobiana de derivadas de $\mathbf{F}$: $(J_{\mathbf{F}})_{ij} = \partial / \partial x^{(j)} F_i(\mathbf{x})$). Para tales sistemas, un conjunto de condiciones iniciales de medida positiva será, después de algún tiempo transitorio, atraído hacia un subconjunto del espacio de fases. Este conjunto es invariante bajo la evolución dinámica y se denomina el \textit{atractor} del sistema. Ejemplos simples de atractores (no caóticos) son puntos fijos (después de que el sistema se estabiliza en un estado estacionario tras el tiempo transitorio) y ciclos límite (el sistema se aproxima a un movimiento periódico). Véase Fig. 3.1. Para un sistema autónomo con dos grados de libertad y tiempo continuo, estas (junto con órbitas homoclínicas y heteroclínicas conectando puntos fijos, las cuales no se discuten aquí) son las únicas posibilidades (teorema de Poincaré–Bendixson). Atractores más interesantes pueden ocurrir en flujos en al menos tres dimensiones (véase Fig. 3.2), donde el mecanismo de estiramiento y plegado puede producir \textit{caos}.  

\bigskip

\textbf{Ejemplo 3.3.} La divergencia del flujo introducido en el Ejemplo 3.1 es cero, por lo que la dinámica conserva el área. El mapa de Hénon, Ejemplo 3.2, es disipativo si el determinante independiente del Jacobiano, $|\det J| = |b|$, se toma menor que la unidad. El sistema de Lorenz (véase Ejercicio 3.2) es fuertemente disipativo, ya que su divergencia, $-b - \sigma - 1$, para la elección usual de parámetros, es aproximadamente $-10$. Esto significa que, en promedio, un volumen se reducirá a $e^{-10}$ de su volumen original en cada unidad de tiempo.

Característico de los sistemas caóticos es que los atractores correspondientes son objetos geométricos complicados, que típicamente exhiben una estructura \textit{fractal}. También se les llama \textit{atractores extraños}. Explicaremos en el Capítulo 6, y en particular en la Sección 6.1, lo que entendemos por un atractor extraño o fractal.

Hasta ahora hemos ilustrado que es natural describir un sistema dinámico determinista como un objeto en el espacio de fases, ya que esta es la manera óptima de estudiar sus propiedades dinámicas y geométricas. Dado que las ecuaciones dinámicas (o las ecuaciones de movimiento) están definidas en el espacio de fases, también es lo más natural utilizar una descripción en espacio de fases para aproximaciones a estas ecuaciones. Dichas dinámicas aproximadas serán importantes para las predicciones (Capítulos 4 y 12), la determinación de exponentes de Lyapunov (Capítulos 5 y 11), el filtrado de ruido (Capítulo 10) y la mayoría de las demás aplicaciones (Capítulo 15). Para una comprensión más profunda de la naturaleza del sistema subyacente, es deseable conocer la geometría del atractor. La estrecha relación entre dinámica y geometría se expresa mediante resultados teóricos que relacionan los exponentes de Lyapunov (dinámica) y las dimensiones (geometría). Sin embargo, los algoritmos prácticos a veces también pueden formularse de manera equivalente en términos de geometría o dinámica. (Véase la discusión sobre métodos de reducción de ruido, Sección 10.3).

\newpage

\section{Reconstrucción con retardos}

Habiendo resaltado la importancia del espacio de fases para el estudio de sistemas con propiedades deterministas, debemos enfrentar el primer problema: lo que observamos en un experimento no es un objeto en el espacio de fases, sino una serie temporal, muy probablemente solo una secuencia de mediciones escalares. Por lo tanto, debemos convertir las observaciones en vectores de estado. Este es el problema importante de la \textit{reconstrucción del espacio de fases}, que se resuelve técnicamente mediante el método de retardos (o construcciones relacionadas).

Lo más común es que la serie temporal sea una secuencia de mediciones escalares de alguna cantidad que depende del estado actual del sistema, tomadas en múltiplos de un tiempo de muestreo fijo:

\begin{equation}
s_n = s(\mathbf{x}(n\Delta t)) + \eta_n .
\end{equation}

Es decir, observamos el sistema a través de alguna función de medición $s$ y realizamos observaciones solo hasta cierto punto de fluctuaciones aleatorias $\eta_n$, el \textit{ruido de medición}. En este nivel de presentación, despreciaremos el efecto del ruido. (Más adelante discutiremos su efecto en la Sección 10.2.)

Una \textit{reconstrucción con retardos} en $m$ dimensiones se forma entonces mediante los vectores $\mathbf{s}_n$, dados por

\begin{equation}
\mathbf{s}_n = (s_{n-(m-1)\tau}, s_{n-(m-2)\tau}, \ldots, s_{n-\tau}, s_n) .
\end{equation}

La diferencia temporal en número de muestras $\tau$ (o en unidades de tiempo, $\tau \Delta t$) entre componentes de los vectores de retardo se denomina \textit{lag} o \textit{tiempo de retardo}. Nótese que para $\tau > 1$, solo aumenta la ventana temporal cubierta por cada vector, mientras que el número de vectores construidos a partir de la serie temporal escalar permanece aproximadamente igual. Esto se debe a que creamos un vector para cada observación escalar, $s_n$, con $n > (m-1)\tau$.

Varios teoremas de incrustación se ocupan de la cuestión de en qué circunstancias y hasta qué punto el objeto geométrico formado por los vectores $\mathbf{s}_n$ es equivalente\footnote{En el sentido de que pueden ser mapeados uno a otro mediante una aplicación diferenciable e invertible. Véase la Sección 9.1.} a la trayectoria original $\mathbf{x}_n$. De hecho, bajo circunstancias bastante generales, el atractor formado por $\mathbf{s}_n$ es equivalente al atractor en el espacio desconocido en el cual vive el sistema original si la dimensión $m$ del espacio de coordenadas con retardos es suficientemente grande. Más precisamente, esto está garantizado si $m$ es mayor que dos veces la \textit{dimensión de conteo de cajas} $D_F$ del atractor; es decir, hablando en términos generales, mayor que el doble del número de grados de libertad \textit{activos}, independientemente de cuán grande sea la dimensionalidad del verdadero espacio de estados. Dependiendo de la aplicación, incluso valores más pequeños de $m$ que satisfagan $m > D_F$ pueden ser suficientes. Para detalles teóricos véase Ding \textit{et al.} (1993).

Se puede esperar reconstruir el movimiento en atractores en sistemas como flujos hidrodinámicos o láseres, donde el número de partículas microscópicas es enorme, si solo unos pocos grados de libertad dominantes permanecen finalmente como resultado de algún comportamiento colectivo. En el Capítulo 9 discutiremos con mayor detalle el trasfondo de la incrustación y los teoremas.

\newpage

\section{Encontrar un buen embebido}

Cuando comenzamos a analizar una serie temporal escalar, no conocemos ni la \textit{dimensión de conteo de cajas}\footnote{La cual es formalmente necesaria para calcular $m$}, ni tenemos idea de cómo elegir el \textit{tiempo de retardo} $\tau$. El procedimiento depende de la dinámica subyacente en los datos y del tipo de análisis que se pretenda realizar. Lo más importante es que los teoremas de incrustación \textit{garantizan} que, para datos ideales libres de ruido, existe una dimensión $m$ tal que los vectores $\mathbf{s}_n$ son equivalentes a los vectores del espacio de fases. Usaremos este conocimiento para diseñar métodos que permitan determinar la \textit{dimensión} del atractor (Capítulo 6), el \textit{exponente de Lyapunov máximo} (Capítulo 5), y la \textit{entropía} (Capítulo 11).

En general, existen dos enfoques diferentes para optimizar los parámetros de embebido $m$ y $\tau$: 
\begin{enumerate}
    \item Explotar herramientas estadísticas específicas para su determinación y luego utilizar sus valores óptimos en análisis posteriores.
    \item Comenzar directamente con el análisis previsto y optimizar los resultados ajustando $m$ y $\tau$.
\end{enumerate}

Por ejemplo, las estimaciones de dimensión y de Lyapunov se obtendrán aumentando los valores de $m$ hasta que aparezca el comportamiento típico de datos deterministas.

Para muchos propósitos prácticos, el parámetro más importante de embebido es el producto $m\tau$ del tiempo de retardo y la dimensión de embebido, más que la dimensión $m$ o el retardo $\tau$ por separado. La razón es que $m\tau$ es el intervalo temporal representado por un vector de embebido (véase Kugiumtzis (1996) para una discusión).

Para mayor claridad, discutamos la elección de $m$ y $\tau$ por separado. Un conocimiento preciso de $m$ es deseable, ya que queremos explotar el determinismo con un esfuerzo computacional mínimo. Por supuesto, si un embebido de dimensión $m$ proporciona una representación fiel del espacio de estados, cualquier reconstrucción en $m' > m$ también lo hará. Sin embargo, elegir un valor demasiado grande de $m$ en datos caóticos introduce redundancia y degrada el rendimiento de muchos algoritmos, como predicciones y cálculos de exponentes de Lyapunov. Debido a la inestabilidad inherente de los datos caóticos, el primer y último elemento de un vector de retardo están más correlacionados cuanto mayor es la diferencia de tiempo. Por lo tanto, tomar un valor demasiado grande para $m$ no sería útil y correría el riesgo de ``confundir'' al algoritmo.

\newpage

\subsection{\textit{Vecinos falsos}}

Si asumimos que la dinámica en el espacio de fases está representada por un campo vectorial suave, entonces los estados vecinos deberían estar sujetos a casi la misma evolución temporal. Por lo tanto, después de un intervalo de tiempo corto hacia el futuro, las dos trayectorias que emergen de ellos aún deberían permanecer como vecinos cercanos, incluso si el caos puede introducir una divergencia exponencial de las dos. Este razonamiento se utilizará extensamente en el próximo capítulo sobre predicción. Aquí queremos referirnos a esta propiedad para discutir estadísticas que a veces ayudan a identificar si cierta dimensión de incrustación es suficiente para una reconstrucción de un espacio de fases.

El concepto, llamado \textit{vecinos falsos más cercanos}, fue introducido por Kennel, Brown y Abarbanel (1992). Lo presentamos aquí con algunas modificaciones menores que evitan ciertos resultados espurios debidos al ruido (Hegger \& Kantz (1999)).

La idea básica es buscar puntos en el conjunto de datos que son vecinos en el espacio de incrustación $m$, pero que no deberían ser vecinos, ya que su evolución temporal futura es demasiado diferente. Imagina que la dimensión de incrustación correcta para algún conjunto de datos es $m_0$. Ahora estudiamos los mismos datos en una incrustación de menor dimensión $m < m_0$. La transición de $m_0$ a $m$ es una proyección, eliminando ciertos ejes del sistema de coordenadas. Por lo tanto, puntos cuyas coordenadas se eliminan por la proyección y difieren fuertemente, pueden convertirse en “vecinos falsos” en el espacio $m$-dimensional.  

La estadística a estudiar ahora es obvia: para cada punto de la serie temporal, tomamos su vecino más cercano en $m$ dimensiones y calculamos la razón de las distancias entre estos dos puntos en $m+1$ dimensiones y en $m$ dimensiones. Si esta razón es mayor que un umbral $r$, el vecino era falso. Este umbral tiene que ser lo suficientemente grande como para permitir una divergencia exponencial debido al caos determinista.

Ahora, cuando denotamos la desviación estándar de los datos por $\sigma$ y usamos la norma máxima, la estadística a computar es

\begin{equation}
X_{fnn}(r) = \frac{\sum_{n=1}^{N-m-1} \Theta \left( \frac{|s_{n}^{(m+1)} - s_{k(n)}^{(m+1)}|}{|s_{n}^{(m)} - s_{k(n)}^{(m)}|} - r \right) \Theta \left( \frac{\sigma}{r} - |s_{n}^{(m)} - s_{k(n)}^{(m)}| \right)}{\sum_{n=1}^{N-m-1} \Theta \left( \frac{\sigma}{r} - |s_{n}^{(m)} - s_{k(n)}^{(m)}| \right)},
\end{equation}

donde $s_{k(n)}^{(m)}$ es el vecino más cercano a $s_n$ en $m$ dimensiones, es decir, $k(n)$ es el índice del elemento $k$ de la serie temporal diferente de $n$ para el cual $|s_n - s_k| = \min$. 

La primera función escalón en el numerador es uno si el vecino más cercano es falso, es decir, si la distancia aumenta en un factor mayor que $r$ al incrementar la dimensión de incrustación en una unidad, mientras que la segunda función escalón suprime todos aquellos pares cuya distancia inicial ya era mayor que $\sigma/r$. Los pares cuya distancia inicial es mayor que $\sigma/r$ por definición no pueden ser vecinos falsos, ya que, en promedio, no hay suficiente espacio para separarse más allá de $\sigma$. Por lo tanto, son candidatos inválidos para el método y no deben ser contados, lo que también se refleja en la normalización.  

Puede haber algunos vecinos falsos incluso cuando se trabaja en la dimensión de incrustación correcta. Paradójicamente, debido al ruido de medición, puede haber más vecinos falsos si se dan más datos. Con más datos, el vecino más cercano suele ser más próximo (hay más oportunidades de recurrencia en el atractor), mientras que la probabilidad de que la distancia de los $m+1$ componentes esté al menos en el orden del nivel de ruido permanece igual. Para datos con discretización muy gruesa (por ejemplo, 8 bits), incluso pueden existir vectores de retardo idénticos en $m$ dimensiones, que no son idénticos en $m+1$ dimensiones, de manera que la razón de distancias diverge. Estos aparecen como vecinos falsos para cualquier elección de $r$.  

No obstante, es razonable estudiar la razón de vecinos falsos $X_{fnn}$ como función de $r$. Los resultados del análisis de vecinos falsos pueden depender del retardo $\tau$. Por lo tanto, si se desea usar esta estadística para distinguir caos de ruido, es indispensable verificar los resultados con una prueba de datos sustitutos (véase el Capítulo 7).

\subsection*{Ejemplo 3.4 (Vecinos más cercanos falsos de datos de un circuito de resonancia)}
El circuito de resonancia eléctrica descrito en el Apéndice B.11 genera datos bastante libres de ruido en un atractor con una dimensión $D_F$ ligeramente superior a dos.  
Las estadísticas de vecinos falsos mostradas en la Fig. 3.3 sugieren que una incrustación de 5 dimensiones es claramente suficiente (y esto también está garantizado por la desigualdad $m > 2D_F$ del teorema de incrustación).  

También es evidente que los pocos vecinos falsos encontrados en 4 dimensiones (para $N=10000$ estos son alrededor de 3 en valor absoluto) son difíciles de interpretar. En Hegger et al.\ (1998) fue posible realizar predicciones bastante precisas y modelado en una incrustación de cuatro dimensiones, por lo que parece que los vecinos falsos son artefactos de ruido.

\newpage

\subsection{\textit{El retardo temporal}}

Una buena estimación del tiempo de retardo $\tau = \tau \Delta t$ es aún más difícil de obtener. 
El retardo no es objeto de los teoremas de incrustación, ya que consideran datos con precisión infinita. 
Las incrustaciones con el mismo $m$ pero diferente $\tau$ son equivalentes en el sentido matemático para datos libres de ruido, 
pero en la realidad una buena elección de $\tau$ facilita el análisis. 
Si $\tau$ es pequeño en comparación con las escalas de tiempo internas del sistema, 
los elementos sucesivos de los vectores de retardo están fuertemente correlacionados. 
Todos los vectores $s_n$ se agrupan entonces alrededor de la diagonal en $\mathbb{R}^m$, 
a menos que $m$ sea muy grande. 
Si $\tau$ es muy grande, los elementos sucesivos ya son casi independientes, 
y los puntos llenan una gran nube en $\mathbb{R}^m$, donde las estructuras deterministas se limitan a escalas muy pequeñas. 
El primer cero de la función de autocorrelación, Ec. (2.5), de la señal suele proporcionar un buen compromiso entre estos extremos. 
Un concepto más refinado se llama \textit{información mutua} y se presentará en la Sección 9.2. 
En este punto solo damos dos recetas. 
La primera es estudiar el Ejemplo 3.5 para tener una idea de cómo se puede verificar visualmente una elección razonable de $\tau$. 
La segunda es que para una señal con un componente (casi) periódico fuerte, un retardo temporal idéntico a un cuarto del período es una primera buena suposición. 
Todas las estadísticas no lineales en los próximos tres capítulos que se basan en el comportamiento de escalamiento sufren de rangos de escalamiento reducidos cuando $\tau$ ha sido elegido inapropiadamente, es decir, si $\tau$ es demasiado grande o demasiado pequeño. 

\bigskip

\textbf{Ejemplo 3.5 (ECG humano).} 
Aunque el electrocardiograma humano (ECG, véase Apéndice B.7) probablemente no sea una señal determinista, puede ser interesante visualizar tales señales en coordenadas de retardo. 
Comparemos tales representaciones para diferentes tiempos de retardo $\tau = \tau \Delta t$, donde $\Delta t$ es el tiempo de muestreo, ilustradas en la Fig. 3.4.

\newpage

\section{Inspección visual de datos}

Aunque este libro contiene muchos métodos refinados para la caracterización y análisis de datos, lo primero que debemos hacer con un nuevo conjunto de datos es mirarlo de varias maneras diferentes. Una gráfica de la señal como función del tiempo ya da las primeras pistas de posibles problemas de estacionariedad, tales como derivas, amplitudes o escalas de tiempo que varían sistemáticamente, y eventos raros. Esto nos permite seleccionar partes de la serie que parecen más estacionarias. Con bastante frecuencia, los datos experimentales contienen algunas fallas que pueden detectarse mediante inspección visual.

El siguiente paso sería una representación bidimensional como la mostrada en la Fig. 1.5. Incluso si no son visibles estructuras claras (como en el panel derecho), esto da una idea de qué desfase temporal puede ser razonable para un embebido del conjunto de datos.

\bigskip

La inspección visual también puede revelar simetrías en los datos o guiarnos hacia una representación más útil de los mismos. Un ejemplo de lo último es la medición de la potencia de salida de un láser. Contiene una distorsión no lineal de las variables físicamente más relevantes, la intensidad de los campos eléctrico y magnético dentro de la cavidad del láser: la potencia es proporcional al \textit{cuadrado} de estas cantidades. Por lo tanto, los datos muestran máximos pronunciados y mínimos suaves, y la señal cambia mucho más rápido alrededor de los máximos que de los mínimos. Tomar la raíz cuadrada de los datos los hace más convenientes para el análisis.

A veces se encuentran simetrías exactas en los datos, por ejemplo, bajo el cambio de signo del observable. En este caso, se puede ampliar la base de datos para un análisis puramente geométrico simplemente replicando cada punto de datos mediante la operación de simetría. Para un análisis de la dinámica, se puede aplicar la operación de simetría a la serie temporal en su conjunto. Como casi todos los métodos no lineales explotan relaciones locales de vecindad, ambos trucos duplican la base de datos. Al modelar la dinámica del sistema, la clase de funciones elegida debe respetar la simetría.

\newpage

\section{Superficie de sección de Poincaré}

Si consideramos el espacio de fases de un sistema de $m$ ecuaciones diferenciales autónomas, encontramos que, localmente, la dirección tangencial al flujo no aporta mucha información interesante. La posición del punto en el espacio de fases a lo largo de esta dirección puede cambiarse reparametrizando el tiempo. No tiene relación con la \textit{geometría} del atractor y no proporciona información adicional sobre la dinámica. Podemos usar esta observación para reducir en uno la dimensionalidad del espacio de fases y, al mismo tiempo, convertir el flujo continuo en el tiempo en un mapa discreto de tiempo.

El método, llamado \textit{sección de Poincaré}, es el siguiente. Primero formamos una superficie adecuada orientada en el espacio de fases (área sombreada en la Fig. 3.5). Podemos construir un mapa invertible en esta superficie siguiendo una trayectoria del flujo. Los iterados del mapa están dados por los puntos donde la trayectoria intersecta la superficie en una dirección especificada (desde arriba en la Fig. 3.5). Nótese que el “tiempo” discreto $n$ de este mapa es el conteo de intersecciones y usualmente no es simplemente proporcional al tiempo original $t$ del flujo. El tiempo que una trayectoria pasa entre dos puntos de intersección sucesivos variará, dependiendo tanto de la trayectoria en el espacio de estados reconstruido como de la superficie de sección elegida. A veces, después de aplicar esta técnica, permanece un número de puntos decepcionantemente pequeño. Los experimentadores tienden a ajustar la tasa de muestreo de modo que se realicen unas 10–100 observaciones por ciclo típico de la señal. Cada ciclo produce como máximo un punto en el mapa de Poincaré (si hay considerablemente menos intersecciones, la superficie es inapropiada). Sin embargo, si la base de datos en el mapa de Poincaré parece pobre, hay que tener en cuenta que el flujo de datos en sí no es mucho mejor, ya que la información adicional que contienen es en gran medida redundante respecto a las propiedades caóticas de la señal.

\bigskip

Además de la construcción de intersecciones, también podemos recolectar todos los mínimos o todos los máximos de una serie temporal escalar. Como discutiremos en la Sección 9.3, la derivada (numérica) temporal de la señal es una coordenada “válida” en un espacio de estados reconstruido. Por lo tanto, en el espacio reconstruido hipotético abarcado por vectores $(s(t), \dot{s}(t), \ddot{s}(t), \ldots)$, las intersecciones de la trayectoria con la superficie dada por $\dot{s}(t) = 0$ son precisamente los mínimos (o máximos) de la serie temporal. Los mínimos (o máximos) se interpretan como la función de medida especial que proyecta sobre el primer componente de un vector aplicado a los vectores de estado dentro de esta superficie. Por lo tanto, deben incrustarse con un retardo temporal con desfase unidad en sí mismos para formar el conjunto invariante de este mapa particular de Poincaré. El paquete TISEAN contiene una utilidad para formar secciones de Poincaré (llamada \texttt{poincare}) y otra que encuentra puntos extremos, llamada \texttt{extrema}. Encontrar buenas secciones a menudo requiere cierta experimentación con los parámetros de configuración.

La clase más simple de sistemas no autónomos que podemos manejar adecuadamente utilizando métodos de sistemas dinámicos no lineales consiste en aquellos que son impulsados periódicamente. Su espacio de fases es el espacio extendido que contiene la fase de la fuerza impulsora como variable adicional. Las superficies naturales de sección son aquellas de fase constante. El mapa resultante de Poincaré también se llama un \textit{mapa estroboscópico}. En este caso, la sección de Poincaré tiene el mismo lapso de tiempo entre intersecciones, lo cual es una característica muy útil para análisis cuantitativos.

\newpage

\textbf{Ejemplo 3.6 (Retrato de fase de datos de láser NMR).} Estos datos comprenden la salida caótica de un sistema periódicamente forzado (Apéndice B.2). El intervalo de muestreo es $\frac{1}{15}$ del período de la fuerza de conducción, $\Delta t = T_{\text{force}}/15$. Consideremos tres de los muchos tipos diferentes de mapas de Poincaré. La vista estroboscópica es la más natural en este caso y se ha utilizado para producir los datos mostrados en la introducción, Fig. 1.3. Aparte de una complicación que se explicará en el Ejemplo 9.4, $s_n = s(t = nT_{\text{force}} + t_0)$. La Figura 3.6 muestra una sección a través de un espacio de incrustación de retardo tridimensional utilizando interpolación lineal (izquierda) a través de un plano diagonal, y la colección de todos los mínimos de la serie temporal continua después de interpolación parabólica local (derecha). Para el primero, la sección se tomó de modo que $s_n^{(1)} = s(t_-) $, donde $t_-$ está definido por $s(t_-) = s(t_- + \tau)$. El retardo se tomó como $\tau = 4$. La segunda coordenada es $s_n^{(2)} = s(t_- - 2\tau)$. Así, la nueva serie temporal es una serie de vectores bidimensionales. Alternativamente, uno podría haber usado $s_n^{(1)}$ como una serie temporal escalar y graficarla como una representación de retardo bidimensional. La superficie que seleccionamos no es óptima ya que en la parte superior izquierda del atractor el ruido se amplifica considerablemente. En esta parte, la superficie de sección corta el atractor casi tangencialmente, de modo que las posiciones precisas de los puntos de intersección se ven fuertemente afectadas por el ruido y por la interpolación. El tiempo medio entre puntos de intersección es de aproximadamente 14. La desviación del valor esperado de 15 se debe a cruces falsos causados por el ruido. El panel derecho utiliza una sección en $\dot{s} = 0, \ddot{s} \geq 0$, es decir, cada mínimo de la serie interpolada parabólicamente. El gráfico muestra una incrustación de retardo de la serie temporal dada por los mínimos de estas parábolas locales. El tiempo medio entre elementos sucesivos de la nueva serie es de aproximadamente 15, y por lo tanto esta sección es tan adecuada como la vista estroboscópica. El ruido en los datos causa algunos extremos espurios (los puntos dispersos fuera de las estructuras principales). 

El atractor se ve diferente en los tres casos ya que sus detalles geométricos dependen de las propiedades precisas de la sección de Poincaré. No obstante, los tres atractores son equivalentes en el sentido de que se caracterizan por valores idénticos de las dimensiones, exponentes de Lyapunov y entropías. Las estimaciones numéricas de estas cantidades a partir de datos finitos y ruidosos pueden diferir considerablemente.

\newpage

\section{Gráficas de recurrencia}

El intento de análisis de series temporales para extraer información significativa de los datos depende en gran medida de las redundancias dentro de los datos. Al menos si los datos son aperiódicos y no se puede discernir una regla simple para su dependencia temporal, entonces las repeticiones aproximadas de ciertos eventos pueden ayudarnos a construir reglas más complicadas. Suponiendo determinismo y creyendo que una incrustación de retardo elegida representa un espacio de estados del sistema, una repetición aproximada se llama \textit{recurrencia}, es decir, el retorno de la trayectoria en el espacio de estados a una vecindad de un punto donde ya ha estado antes. Tales recurrencias existen en el espacio original para todos los tipos de movimiento que no son transitorios. Un sistema en un punto fijo es trivialmente recurrente para todos los tiempos. En un sistema en un ciclo límite, cada punto retorna exactamente a sí mismo después de una revolución. Un sistema en un atractor caótico retorna a una vecindad arbitrariamente pequeña de cualquiera de sus puntos. Esto está garantizado por la invarianza del conjunto que forma el soporte del atractor. Si, sin embargo, el sistema nunca retorna a todos los puntos que encontramos en la parte inicial de la serie temporal, entonces esto indica que esto fue un transitorio — la condición inicial estaba fuera del conjunto invariante y la trayectoria se relaja hacia este conjunto. Asimismo, no estacionariedades a través de parámetros del sistema dependientes del tiempo pueden causar una falta de recurrencia.

Un método muy simple para visualizar recurrencias se llama \textit{gráfica de recurrencia} y fue introducido por Eckmann \textit{et al.} (1987): Calcular la matriz

\[
M_{ij} = \Theta(\epsilon - |s_i - s_j|) ,
\]

donde $\Theta(.)$ es la función escalón de Heaviside, $\epsilon$ es un parámetro de tolerancia que debe elegirse, y $s_i$ son vectores de retardo de alguna dimensión de incrustación. Esta matriz es simétrica por construcción. Si la trayectoria en el espacio reconstruido retorna en el tiempo $i$ a la $\epsilon$-vecindad de donde estaba en el tiempo $j$ entonces $M_{ij} = 1$, de lo contrario $M_{ij} = 0$. Uno puede graficar $M_{ij}$ con puntos negros y blancos en el plano de índices para inspección visual. Esto es la gráfica de recurrencia. También se han propuesto esquemas numéricos para su caracterización cuantitativa. Estos son similares a muchas otras herramientas estadísticas basadas en vecindarios en el espacio de incrustación y por lo tanto no se discutirán aquí. Ver Casdagli (1997) para más detalles.

Entonces, ¿qué podemos aprender de una inspección visual de la matriz $M_{ij}$ mediante una gráfica de recurrencia? Nos da pistas sobre la serie temporal y sobre el espacio de incrustación en el cual estamos trabajando. En un sistema determinista, dos puntos que están cerca deberían tener imágenes bajo la dinámica que también estén cerca (incluso si no están tan cerca debido a inestabilidad dinámica). Por lo tanto, se espera que los puntos negros aparezcan típicamente como segmentos de línea cortos paralelos a la diagonal. Si además hay muchos puntos aislados, estos más bien indican cercanía coincidente y, por lo tanto, una fuerte componente de ruido en los datos, o una dimensión de incrustación insuficiente. Si hay principalmente puntos dispersos

donde $\Theta(.)$ es la función escalón de Heaviside, $\epsilon$ es un parámetro de tolerancia a elegir, y $s_i$ son vectores de retraso de alguna dimensión de embebido. Esta matriz es simétrica por construcción. Si la trayectoria en el espacio reconstruido retorna en el tiempo $i$ a la $\epsilon$-vecindad de donde estaba en el tiempo $j$ entonces $M_{ij} = 1$, de lo contrario $M_{ij} = 0$. Se puede graficar $M_{ij}$ con puntos negros y blancos en el plano de índices para inspección visual. Este es el \textit{recurrence plot}. También se han propuesto esquemas numéricos para su caracterización cuantitativa. Estos son similares a muchas otras herramientas estadísticas basadas en vecindades en el espacio de embebido y por lo tanto no serán discutidos aquí. Véase Casdagli (1997) para más detalles.

Entonces, ¿qué podemos aprender de una inspección visual de la matriz $M_{ij}$ mediante un \textit{recurrence plot}? Proporciona pistas sobre la serie temporal y sobre el espacio de embebido en el que estamos trabajando. En un sistema determinista, dos puntos cercanos deberían tener también imágenes bajo la dinámica que son cercanas (aunque no necesariamente tan cercanas debido a inestabilidad dinámica). Por lo tanto, se espera que los puntos negros aparezcan típicamente como segmentos de línea cortos paralelos a la diagonal. Si además hay muchos puntos aislados, estos más bien indican cercanía coincidental y, por tanto, una fuerte componente de ruido en los datos, o una dimensión de embebido insuficiente. Si hay principalmente puntos dispersos, el componente determinista está ausente o al menos es débil. Los puntos negros en altas dimensiones también son negros en embebidos de menor dimensión. Por lo tanto, el número relativo de líneas aumenta al incrementar la dimensión de embebido, ya que desaparecen los puntos aislados.

\newpage

\textbf{Ejemplo 3.7 (Recurrence plot de datos de un circuito de resonancia eléctrica).} En la Fig. 3.7 mostramos un \textit{recurrence plot} de los datos de un láser NMR en dimensiones de embebido 2 y 4.

La irregularidad en la disposición de los segmentos de línea indica caos. Si la señal fuera periódica, estos deberían formar un patrón periódico. La estacionariedad de toda la serie temporal requiere que la densidad de segmentos de línea sea uniforme en el plano $i-j$. Véase el Ejemplo 13.8 para un ejemplo no estacionario. El espaciado temporal entre líneas paralelas también puede dar pistas sobre la existencia de órbitas periódicas inestables dentro del atractor caótico.

El inconveniente esencial de esta útil herramienta visual es el hecho de que para una serie temporal de longitud $N$ la matriz $M_{ij}$ tiene $N^2$ elementos, por lo que se debe dibujar un enorme número de píxeles. Una versión comprimida del \textit{recurrence plot} fue bautizada \textit{meta-recurrence plot} por Manuca \& Savit (1996). Ellos subdividen la serie temporal en segmentos y definen los elementos de la matriz $M_{ij}$ para representar distancias entre los segmentos de las series $i$ y $j$. Una posible medida de distancia es el \textit{cross prediction error}. Véase la discusión en Schreiber (1999).

\textbf{Ejemplo 3.8 (Meta-recurrence plot de datos de ECG).} Para datos de ECG, se pueden alinear segmentos a lo largo de cada latido y luego tomar una distancia euclidiana. Definimos un segmento tomando los últimos 100 ms antes y los 500 ms después del cruce ascendente en el complejo QRS (véase Apéndice B.7). La matriz de distancias resultante se representa en escala de grises en la Fig. 3.8. La no-estacionariedad es claramente visible. Más importante aún, casi siempre se pueden reconocer episodios estacionarios como bloques claros en la diagonal.

\section*{Lecturas recomendadas}

Las fuentes originales sobre embebidos en espacio de fases son Takens (1981), Casdagli \textit{et al.} (1991), y Sauer \textit{et al.} (1991). Discusiones sobre la elección adecuada de parámetros de embebido se encuentran, por ejemplo, en Fraser \& Swinney (1986), Liebert \& Schuster (1989), Liebert \textit{et al.} (1991), Buzug \& Pfister (1992), Kennel \textit{et al.} (1992), y en Kugiumtzis (1996). Véase también Grassberger \textit{et al.} (1991). Los \textit{recurrence plots} fueron introducidos por Eckmann \textit{et al.} (1987). Hay un buen recuento más reciente en Casdagli (1997).

\newpage

\section{Ejercicios}

\begin{enumerate}
	\item Cree diferentes retratos de fase bidimensionales de una serie temporal del mapa de Hénon (Ejemplo 3.2) usando tiempos de retardo $\tau = 1, 2, \ldots$ (aquí el tiempo de muestreo es $\Delta t = 1$). ¿Qué figura ofrece la información más clara? Reescriba el mapa en coordenadas de retardo con retardo unidad.

\newpage

	\item Cree una serie escalar de datos de flujo numéricamente mediante la integración del sistema de Lorenz [Lorenz (1963)]:
	\begin{equation}
		\begin{aligned}
			\dot{x} &= \sigma (y - x),\\
			\dot{y} &= r x - y - x z,\\
			\dot{z} &= - b z + x y,
		\end{aligned}
	\end{equation}
	con los parámetros $\sigma = 10$, $r = 28$ y $b = 8/3$. Este modelo fue diseñado para describir el movimiento convectivo de tipo Rayleigh–Bénard, donde $x$ es la velocidad del fluido, $y$ es la diferencia de temperatura entre el fluido ascendente y el descendente, y $z$ es la desviación del perfil de temperatura respecto de la linealidad. Para un sistema láser gobernado por las mismas ecuaciones y datos experimentales, véase el Apéndice B.1. Como integrador, use, por ejemplo, una rutina de Runge–Kutta con un paso pequeño. Muestree los datos a una tasa tal que registre en promedio unas 25 mediciones escalares de la coordenada $x$ durante una sola vuelta de la trayectoria en una hoja del atractor. Represente el atractor en coordenadas de retardo bidimensionales con diferentes retardos temporales. Convénzase por inspección visual de que la reconstrucción es mejor cuando el retardo es aproximadamente un cuarto del tiempo medio de ciclo. Calcule la función de autocorrelación y la información mutua con retardo temporal (que se definirá en la Sección 9.2) para confirmar esta impresión.

\newpage

	\item Integre numéricamente el sistema de Lorenz y realice una sección de Poincaré. Registre $(y,z)$ cada vez que la coordenada $x$ sea igual a cero y su derivada sea negativa.

\newpage

	\item Use la serie temporal de la variable $x$ del Ejercicio 3.2 y recoja todos los 	máximos. Represente la serie de máximos con retardo uno y compárela con el atractor del Ejercicio 3.3.
\end{enumerate}

\newpage

\chapter{Determinismo y predictibilidad}

En este capítulo discutiremos la noción de la predictibilidad de un sistema que evoluciona en el tiempo o, estrictamente hablando, de una señal emitida por tal sistema. Predecir valores futuros de alguna magnitud es un problema clásico en el análisis de series temporales, pero la importancia conceptual del problema de la predicción no se limita a quienes desean enriquecerse conociendo las tasas de cambio de mañana. Incluso si, en cambio, usted está interesado en describir, entender o clasificar señales, acompáñenos en estas páginas. En este libro nos ocupamos de la detección y cuantificación de estructuras posiblemente complicadas en una señal. Queremos poder convencer a otros de que las estructuras que encontramos son reales y no solo fluctuaciones. El argumento más convincente para la existencia de un patrón es si puede usarse para dar una mejor predicción. Es una condición necesaria para que una teoría sea compatible con los datos conocidos, pero no es suficiente. Para ser aceptada, una teoría debe predecir con éxito algo que pueda verificarse posteriormente. En el análisis de series temporales, podemos tomar este requisito de calidad predictiva de manera bastante literal.

\bigskip

La mayoría de los conceptos que introduciremos más adelante para describir datos de series temporales pueden interpretarse en cierta medida como medidas indirectas de predictibilidad. Debido a su naturaleza indirecta, algunas conclusiones seguirán siendo controvertidas, especialmente si las estructuras son bastante débiles. La capacidad estadísticamente significativa de predecir la señal mejor que otras técnicas será entonces una afirmación más convincente de estructura no lineal y determinista que varios dígitos dudosos de la dimensión fractal. Los lectores cuyo interés principal sea la predicción deberían ver este capítulo únicamente como una introducción a conceptos que serán elaborados más adelante en el capítulo sobre modelado y predicción, Capítulo 12.

\section{Fuentes de predictibilidad}

Una señal que no cambia es trivial de predecir: la última observación es un pronóstico perfecto para la siguiente. Incluso si la señal cambia, este puede ser un método razonable de predicción, y una señal para la cual esto se cumple se denomina \textit{persistente}. Una señal que cambia periódicamente con el tiempo también es fácil de predecir una vez que se ha observado un ciclo completo. Los números aleatorios independientes también son fáciles: no es necesario esforzarse porque el esfuerzo no ayuda de todos modos. La mejor predicción es simplemente el valor medio. Lo interesante son las señales que se encuentran entre estos extremos; no son periódicas pero contienen algún tipo de estructura que puede aprovecharse para obtener mejores predicciones.

Antes de mencionar las estructuras más comunes que pueden explotarse para realizar predicciones, establezcamos cómo cuantificamos el éxito de las predicciones que hacemos. La medida de error que más se utiliza es el \textit{error cuadrático medio} (\textit{root mean squared, rms prediction error}). Si predecimos el resultado de las mediciones en los tiempos $n$ como $\hat{s}_n$, mientras que las mediciones reales son $s_n$, entonces el error rms de predicción es

\begin{equation}
e = \sqrt{\langle (\hat{s}_n - s_n)^2 \rangle} ,
\end{equation}

donde $\langle \cdot \rangle$ denota el promedio sobre todas las pruebas que hemos realizado. Puede verificarse fácilmente que para números aleatorios independientes este error se minimiza con $\hat{s}_n = \text{const.} = \langle s_n \rangle$. Otras medidas de error son el \textit{error absoluto medio} ($\langle |\hat{s}_n - s_n| \rangle$), o el \textit{error logarítmico} ($\langle |\ln \hat{s}_n - s_n| \rangle$). Las funciones de error o de costo pueden ser mucho más complicadas, en particular cuando hay dinero en juego. Si viaja a un país donde las tarjetas de crédito no se aceptan y predice que necesita una cantidad de \$100 en efectivo, entonces estaría en problemas: tendría que transferir más dinero por cable, podría perder su vuelo a casa y acabar en prisión antes de que alguien lo ayude. Aunque los \$100 en su bolsillo sean inofensivos, el costo máximo de perderlos es de \$100.

Si por alguna razón utiliza una función de costo más sofisticada, la medida de error no es necesariamente la mejor predicción para su propósito. La predictibilidad está aumentada por el conocimiento de las funciones de error específicas involucradas.

\bigskip

La media y la distribución de probabilidad son características estáticas de los datos y no involucran correlaciones en el tiempo. Dependiendo de la fuerza y el tipo de correlaciones, podemos mejorar las predicciones considerablemente. Las fuentes de predictibilidad más comunes y estudiadas son las correlaciones lineales en el tiempo. Supongamos un proceso completamente aleatorio con una función de autocorrelación que decae lentamente. La aleatoriedad nos impide hacer predicciones precisas, incluso con un conocimiento absoluto del estado presente. La fuerte correlación nos asegura que la siguiente observación estará dada por una combinación lineal de las observaciones anteriores, con una incertidumbre de media cero,

\begin{equation}
\hat{s}_{n+1} = \bar{s} + \sum_{j=1}^{m} a_j (s_{n-m+j} - \bar{s}), 
\end{equation}

donde el valor medio $\bar{s}$ de la serie temporal se separa convenientemente. Como mostramos en la Sección 2.3, la función de autocorrelación y el espectro de potencia están íntimamente relacionados por el \textit{teorema de Wiener–Khinchin} (véase Sección 2.3), que son objetos destacados de la inferencia estadística lineal. Su información puede explotarse exhaustivamente mediante predictores lineales. Sin embargo, debido a la naturaleza aleatoria asumida de los datos, los modelos subyacentes de \textit{auto-regresión} y \textit{promedios móviles} son estocásticos.

Otro tipo de correlación temporal está presente en los sistemas dinámicos no lineales deterministas. Recordemos que los sistemas dinámicos puros se describen ya sea por mapas temporales discretos,  

\begin{equation}
\mathbf{x}_{n+1} = \mathbf{F}(\mathbf{x}_n), \quad n \in \mathbb{Z}, 
\end{equation}

o por ecuaciones diferenciales ordinarias de primer orden,  

\begin{equation}
\frac{d}{dt} \mathbf{x}(t) = \mathbf{f}(\mathbf{x}(t)), \quad t \in \mathbb{R}.
\end{equation}

Ambas variantes son descripciones matemáticas del hecho de que todos los estados futuros de un sistema están \textit{inequívocamente} determinados al especificar su estado presente en algún tiempo $n$ (respectivamente $t$). Esto implica que también existe una función determinista de predicción. Sin embargo, como veremos más adelante, cualquier inexactitud en nuestro conocimiento del estado presente también evolucionará con el tiempo y, en el caso de un sistema caótico, crecerá exponencialmente. En este último caso, ya no podemos usar la evolución temporal para calcular estados arbitrariamente lejanos en el futuro, ya que nunca podemos conocer una cantidad física con incertidumbre cero.  

\bigskip

No obstante, incluso para un sistema caótico, la incertidumbre se amplifica solo a una tasa finita, y aún podemos esperar hacer predicciones razonables a corto plazo explotando la ley determinista de evolución. En contraste con las correlaciones reflejadas por la función de autocorrelación, aquellas impuestas por la dinámica determinista no lineal pueden ser visibles únicamente al usar estadísticas no lineales. Estas suelen llamarse \textit{correlaciones no lineales}, y se deben emplear nuevas técnicas para explotarlas en las predicciones. Por supuesto, existen muchos datos que no son ni caos determinista ni procesos estocásticos no lineales. No siguen ninguno de estos dos paradigmas, y solo en ciertos casos podremos hacer algo con ellos. Véase la discusión en la Sección 12.1.

\section{Algoritmo simple de predicción no lineal}

Seamos ahora más específicos y desarrollemos un algoritmo de predicción simple que explote las estructuras deterministas en la señal. Como con la mayoría de los algoritmos en este campo, procederemos suponiendo que los datos se originan en un sistema dinámico subyacente. Una vez que hayamos establecido el algoritmo, tendremos que estudiar su comportamiento en casos más realistas.

Un conjunto de datos determinista, con valores vectoriales, muestreado en tiempos discretos en su espacio de estados, está perfectamente descrito por la Ec. (4.3). Desafortunadamente, solo si conocemos la aplicación $\mathbf{F}$ podemos usar esto como un esquema de predicción. Conocer $\mathbf{F}$ es una expectativa poco realista cuando trabajamos con datos del mundo real. Con $\mathbf{F}$ desconocida, tenemos que hacer algunas suposiciones sobre sus propiedades. Usando la suposición mínima de que la aplicación $\mathbf{F}$ es continua\footnote{Usualmente podemos permitir que la aplicación esté compuesta de algunas piezas continuas.}, podemos construir un esquema de predicción muy simple. Para predecir el estado futuro $\mathbf{x}_{N+1}$, dado el estado presente $\mathbf{x}_N$, buscamos en una lista de todos los estados pasados $\mathbf{x}_n$ con $n < N$ aquel más cercano a $\mathbf{x}_N$ respecto a alguna norma. Si el estado en el tiempo $n_0$ era similar al presente (y por tanto cercano en el espacio de fases), la continuidad de $\mathbf{F}$ garantiza que $\mathbf{x}_{n_0+1}$ también estará cerca de $\mathbf{x}_{N+1}$. Dado que por suposición $n_0 < N$, también $n_0 + 1 \leq N$, y por lo tanto podemos leer $\mathbf{x}_{n_0+1}$ del registro de datos. Si hemos observado el sistema durante mucho tiempo, habrá estados en el pasado que sean arbitrariamente cercanos al estado presente, y nuestra predicción, $\hat{\mathbf{x}}_{N+1} = \mathbf{x}_{n_0+1}$, estará arbitrariamente cerca de la verdad. Así, en teoría tenemos un algoritmo de predicción muy sencillo y elegante. Usualmente se le conoce como el “método de los análogos de Lorenz”, ya que fue propuesto como un método de predicción por Lorenz (1969).

\bigskip

Ahora enfrentemos la realidad paso a paso. Incluso si la suposición de que existe un sistema determinista subyacente es correcta, normalmente no medimos los estados reales $\mathbf{x}_n$. En su lugar, observamos una (o unas pocas) cantidades que dependen funcionalmente de estos estados. Más comúnmente tenemos mediciones escalares

\[
s_n = s(\mathbf{x}_n), \qquad n = 1, \ldots, N,
\tag{4.5}
\]

en las cuales, la mayoría de las veces, la función de medición $s$ es tan desconocida como $\mathbf{F}$. Obviamente no podemos invertir $s$, pero podemos usar una \textit{reconstrucción con retardos} (véase la Sección 3.2) para obtener vectores equivalentes a los originales:

\[
\mathbf{s}_n = (s_{n-(m-1)\tau}, s_{n-(m-2)\tau}, \ldots, s_{n-\tau}, s_n).
\tag{4.6}
\]

Este procedimiento introducirá dos parámetros ajustables en el método de predicción (que en principio está libre de parámetros): el tiempo de retardo $\tau$ y la dimensión de incrustación $m$. El método resultante sigue siendo muy simple. Para todas las mediciones $s_1, \ldots, s_N$ disponibles hasta ahora, se construyen los vectores de retardo correspondientes $s_{(m-1)\tau+1}, \ldots, s_N$. Para predecir una medición futura $s_{N+\Delta n}$, se busca el vector de incrustación $s_{n_0}$ más cercano a $s_N$ y se usa $s_{n_0+\Delta n}$ como predictor. Este esquema ha sido usado en pruebas de determinismo por Kennel \& Isabelle (1992), por ejemplo. Hasta donde sabemos, la idea se remonta a Pikovsky (1986). El método usado por Sugihara \& May (1990) está estrechamente relacionado.

\bigskip

Como un siguiente paso hacia la realidad, ya no podemos ignorar el hecho de que toda medición de una cantidad continua solo es válida hasta cierta resolución finita. Esta resolución está limitada principalmente por las fluctuaciones en el equipo de medición, que con suerte son aleatorias, y por el hecho de que los resultados eventualmente se almacenan en alguna forma discretizada. En cualquier caso, la resolución finita, llamemos a su tamaño típico $\sigma$, implica que buscar el punto más cercano en el pasado ya no es lo mejor que podemos hacer, ya que las distancias entre puntos están contaminadas con una incertidumbre del orden de $\sigma$. Todos los puntos dentro de una región en el espacio de fases de radio $\sigma$ deben ser considerados igualmente buenos predictores \textit{a priori}. En lugar de elegir uno de ellos arbitrariamente, proponemos tomar la media aritmética de las predicciones individuales basadas en estos valores. La resolución estimada de las mediciones introduce el tercer parámetro en nuestro esquema de predicción.

\bigskip

Ahora estamos listos para proponer el siguiente algoritmo de predicción. Dada una serie temporal escalar $s_1, \ldots, s_N$, se elige un tiempo de retardo $\tau$ y una dimensión de incrustación $m$ y se forman los vectores de retardo $s_{(m-1)\tau+1}, \ldots, s_N$ en $\mathbb{R}^m$. Para predecir un tiempo $\Delta n$ hacia adelante desde $N$, se elige el parámetro $\epsilon$ del orden de la resolución de las mediciones y se forma una vecindad $\mathcal{U}_\epsilon(s_N)$ de radio $\epsilon$\footnote{Por simplicidad, proponemos el uso de la norma máxima: un punto pertenece a 
\(\mathcal{U}_{\epsilon}(s_{N})\) si ninguna de sus coordenadas difiere en más de 
\(\epsilon\) de la coordenada correspondiente de \(s_{N}\).
} alrededor del punto $s_N$. Para todos los puntos $s_n \in \mathcal{U}_\epsilon(s_N)$, es decir, todos los puntos más cercanos que $\epsilon$ a $s_N$, se buscan las “predicciones” individuales $s_{n+\Delta n}$. La predicción es entonces el promedio de todas estas predicciones individuales:

\[
\hat{s}_{N+\Delta n} = \frac{1}{|\mathcal{U}_\epsilon(s_N)|} \sum_{s_n \in \mathcal{U}_\epsilon(s_N)} s_{n+\Delta n}.
\tag{4.7}
\]

Aquí $|\mathcal{U}_\epsilon(s_N)|$ denota el número de elementos de la vecindad $\mathcal{U}_\epsilon(s_N)$. En el caso en que no encontremos ningún vecino más cercano que $\epsilon$, podríamos simplemente aumentar el valor de $\epsilon$ hasta encontrar alguno. Una implementación simple pero efectiva de este esquema es el programa \texttt{zeroth} incluido en el paquete TISEAN y discutido en la Sección A.4.1. Más adelante en este libro describiremos métodos más refinados de predicción y modelado. En este escenario más general, resultará que el esquema descrito aquí se basa en una aproximación de orden cero de la dinámica, es decir, es un predictor localmente constante en lugar de un modelo lineal local o incluso global no lineal. Aún más interesante, se demostrará que para ciertos procesos no lineales estocásticos, en particular cadenas de Markov, este método también funciona. Si se necesitan más de unas pocas predicciones, por ejemplo, porque el error de predicción promedio debe determinarse con precisión, esta implementación es muy ineficiente. Para cada predicción, se consideran \textit{todos} los vectores en el espacio de incrustación (excluyendo únicamente aquellos en la vecindad temporal del que se va a predecir con el fin de evitar conflictos de causalidad). Se prueban para verificar su cercanía, y la mayoría de ellos se rechazan por estar demasiado lejos. Por lo tanto, es recomendable utilizar un método eficiente para encontrar los vecinos más cercanos en el espacio de fases, véase la Sección A.1.1.

\newpage

\section{Verificación de una predicción exitosa}

Si realizas una predicción que resultó ser cierta, tu éxito podría haber sido simplemente una coincidencia. Para tener significancia, debes estudiar los errores de muchas predicciones partiendo de estados reales independientes. Este enfoque puede ser muy tedioso. Imagina que pudieras predecir el 1 de noviembre si habrá nieve en la víspera de Navidad (digamos en Copenhague; es demasiado fácil predecirlo en Montreal). Sólo después de un par de años será estadísticamente significativo saber si tu predicción es útil. En la práctica, preferirías probar tu método en un registro de los últimos cien años o algo así, pero surge un problema: los eventos “futuros” que estás prediciendo ya son conocidos, y tendrás que asegurarte de no usar ninguna información existente sobre los resultados que deseas predecir. Si utilizas el pasado para optimizar tu método de predicción, eventualmente tendrás en cuenta todas las coincidencias curiosas que ocurrieron en el último siglo. Sin embargo, ya no es apropiado citar el rendimiento del predictor en este registro como el error de predicción. Tal error se denominaría un \textit{error de predicción in-sample}. 

\bigskip

Si alguien intenta convencerte de algo mostrando un error de predicción impresionante, asegúrate siempre de que sea un \textit{error de predicción out-of-sample}! (De lo contrario, se debería hablar más bien de una “post-dicción”). La predicción out-of-sample significa que el conjunto de datos dado se divide en un \textit{training set} y un \textit{test set}. Los parámetros del predictor (p. ej., dimensión de incrustación, tiempo de retardo, tamaño de vecindad en la Ecuación (4.7)), se optimizan minimizando el error de predicción en el conjunto de entrenamiento. El número que se utiliza para evaluar su rendimiento debe ser el error de predicción en el conjunto de prueba. El predictor de la Ecuación (4.7) contiene una ambigüedad con respecto a esto: ¿la vecindad \(\mathcal{U}_{\epsilon}(s_{N})\) debe tomarse del conjunto de entrenamiento o del conjunto de prueba? El concepto de causalidad nos permite buscar vecinos en toda la historia del estado actual \(s_{N}\), ignorando la división en conjunto de entrenamiento y de prueba. Sin embargo, la historia implica que también los “futuros” de los vecinos, los elementos \(s_{n+\Delta n}\), deben estar en el pasado del estado actual, es decir, \(n + \Delta n < N\). Esto es importante para predicciones a largo plazo de datos altamente muestreados con gran \(\Delta n\).

\bigskip

El hecho de que los errores \textit{in-sample} a menudo sean mucho más pequeños que los errores \textit{out-of-sample} es más relevante cuanto más parámetros se ajustan para optimizar el rendimiento en el conjunto de entrenamiento. El predictor no lineal simple, Ecuación (4.7), casi puede ser llamado un modelo libre de parámetros, ya que los pocos parámetros se fijan inicialmente. No obstante, citaremos errores \textit{out-of-sample} en los ejemplos. Estos han sido estimados a partir de datos que fueron reservados mientras se construía el predictor. Una excepción notable donde el uso de errores de predicción \textit{in-sample} puede estar justificado son las pruebas de no linealidad, en las cuales los errores no se evalúan por sí solos, sino que se comparan con conjuntos de datos \textit{surrogate}. Una discusión más profunda sobre la relevancia estadística de cantidades tales como los errores de predicción será tratada en el Capítulo 7. A menudo, la base de datos es demasiado pequeña para dividirla en conjuntos de entrenamiento y de prueba. Entonces, se emplea una variante de esto, llamada \textit{leave-one-out statistics} o \textit{complete cross-validation}: para cada predicción individual, el predictor utiliza información de todo el conjunto de datos, excluyendo únicamente la información que está temporalmente correlacionada con la predicción real que se va a realizar. En el predictor simple, Ecuación (4.7), esto significa que aquellos vectores de retardo cuyos índices temporales cumplen $|n - N| < n_{corr}$ son excluidos del vecindario $\mathcal{U}_\epsilon(s_N)$, donde $n_{corr}$ debe fijarse de acuerdo con el horizonte de predicción $\Delta n$ y las correlaciones en el conjunto de datos.

\newpage

\textbf{Ejemplo 4.1 (Predicción de datos de láser NMR).} Probemos el algoritmo propuesto en un conjunto de datos experimental pero altamente determinista. En la introducción, mencionamos el experimento de láser NMR llevado a cabo en el Instituto de Física de la Universidad de Zúrich. El conjunto de datos que usamos para producir la Fig. 1.2 consistió en 38,000 puntos, exactamente un punto por ciclo de conducción (véase el Apéndice B.2 para más información). Dividimos el conjunto de datos en dos partes. Los primeros 37,900 puntos forman la base de datos para el algoritmo de predicción y los últimos 100 puntos forman el conjunto de prueba. En cada uno de estos últimos realizamos predicciones para $\Delta n = 1, \ldots, 35$ pasos de tiempo hacia adelante. Para cada horizonte de predicción $\Delta n$, calculamos el error cuadrático medio de predicción $e_{\Delta n} = \sqrt{\langle (s_{n+\Delta n} - \hat{s}_{n+\Delta n})^2 \rangle}$, el cual mide la desviación promedio de las predicciones respecto a los valores reales medidos posteriormente. Este error de predicción se muestra como diamantes en la Fig. 1.2 (véase también la Fig. 4.1). Utilizamos el programa \texttt{zeroth} (véase Apéndice A.4.1) con dimensión de incrustación $m = 3$, tiempo de retraso $\tau = 1$ y tamaño de vecindad $\epsilon = 50$, lo cual es aproximadamente el doble de la amplitud del error de medición en este conjunto de datos. Nótese que, para el régimen de pequeña amplitud, el incremento en el error está bien descrito por una exponencial $e_{\Delta n} \approx \epsilon e^{0.3 \Delta n}$, lo que sugiere que el error de predicción está dominado por la divergencia caótica de trayectorias inicialmente cercanas. Más adelante estimaremos (en el Capítulo 5) el exponente de Lyapunov máximo de este conjunto de datos en $\lambda = 0.3$. Sin embargo, no recomendamos el uso del incremento en el error de predicción como método para estimar el exponente de Lyapunov, ya que el promedio se realiza de manera incorrecta.

\bigskip

\textbf{Ejemplo 4.2 (Predicción de la tasa de respiración humana).} También podemos intentar predecir datos que no son tan claramente deterministas como los datos de láser NMR en el Ejemplo 4.1. Sin embargo, debemos ser cuidadosos al interpretar los resultados. Si existen correlaciones lineales en los datos, el predictor no lineal las captará en cierta medida. Así, si un algoritmo basado en la teoría del caos es capaz de predecir mejor que el azar, esto por sí solo no puede tomarse como una prueba de caos determinista en los datos. En el Capítulo 7 discutiremos cómo hacer afirmaciones más significativas en este sentido. Ahora intentemos pronosticar valores futuros del conjunto de datos de la tasa de respiración humana mencionado en la introducción (Ejemplo 1.2, Fig. 1.4; véase también el Apéndice B.5). Nuevamente dividimos el conjunto de datos en dos partes, esta vez con una base de datos que contiene 1500 puntos. Para cada horizonte de predicción, el error de predicción se estima mediante un promedio de 100 ensayos. La Figura 4.2 muestra el error cuadrático medio (rms) de predicción para el mejor predictor lineal que pudimos encontrar (diamantes), y para el simple predictor no lineal introducido en esta sección (cruces). El predictor lineal tiene cinco coeficientes, mientras que el no lineal usa vecindades de $\epsilon = 1000$ en dos dimensiones. Los resultados para ambos no dependen mucho de los parámetros elegidos, y en consecuencia no se realizó ninguna optimización sistemática de estos parámetros. Así, los errores de predicción pueden considerarse como errores fuera de muestra. Observamos que el método no lineal es capaz de hacer predicciones ligeramente mejores en tiempos cortos (hasta tres unidades de tiempo), mientras que el método lineal es marginalmente mejor a partir de ahí. Nótese el mucho mayor error de predicción relativo (es decir, error como fracción de la desviación estándar de la señal) que en el ejemplo del láser NMR. Existe un componente mucho más fuerte efectivamente aleatorio en este conjunto de datos. Revisaremos este ejemplo en el Capítulo 7, donde lo investigaremos con respecto a posibles estructuras no lineales.

\bigskip

La calidad de una predicción no lineal solo puede evaluarse mediante la comparación con algún punto de referencia dado por técnicas estándar. Llamemos la atención sobre una forma común de presentar los resultados de la predicción que consideramos no siempre instructiva y, a veces, engañosa. A menudo se muestran gráficos en los que una serie temporal aparece junto con una o varias predicciones a unos pocos pasos adelante. Si la serie temporal está autocorrelacionada (como es típico en un proceso de tiempo continuo), no resulta particularmente difícil predecir unos pocos pasos hacia adelante. Como ejemplo extremo, es fácil predecir el clima dentro de diez minutos a partir de ahora. En analogía con esto, no nos sorprende que las predicciones lineales locales en Abarbanel (1996) sobre el volumen del Gran Lago Salado en Utah resulten bastante convincentes. Sin embargo, en este caso podríamos verificar que el horizonte de predicción de un paso, comparado con el tiempo de autocorrelación, es tan corto que el error de predicción puede superarse fácilmente con el modelo lineal 
\[
s_{n+1} = 2s_n - s_{n-1},
\] 
que simplemente extrapola la tendencia de las dos últimas mediciones. Esquemas de predicción más sofisticados y muchos temas relacionados con el modelado y la predicción se discutirán en el Capítulo 12.

\newpage

\section{Errores de predicción cruzada: examinando la estacionariedad}

Resulta que el esquema de predicción localmente constante descrito en la Sección 4.2 es uno de los indicadores más robustos de estructura no lineal en una serie temporal y es particularmente sensible al comportamiento determinista. Por ello, lo utilizaremos más adelante en la Sección 7.1 como una estadística discriminante en una prueba de no linealidad. En efecto, este esquema simple funciona con muchos menos datos que, por ejemplo, los predictores lineales locales (véase la Sección 12.3.2) y, a diferencia de los ajustes no lineales globales (véase la Sección 12.4), no requiere una elección elaborada de funciones base para cada nuevo conjunto de datos. Estas propiedades hacen que el algoritmo sea adecuado como una sonda tosca pero sensible siempre que busquemos (no lineal) predictibilidad. Como ejemplo de las muchas situaciones en que tal sonda puede resultar útil, usémosla para estudiar la no estacionariedad en un conjunto de datos largo.

\bigskip

La idea es que, en un sistema donde diferentes dinámicas actúan en distintos momentos, los registros de datos de episodios diferentes no son igualmente útiles como bases para predicciones. Dividamos los datos disponibles en segmentos $S_i, \quad i = 1, \dots, I$, de tal forma que cada segmento sea lo suficientemente largo para permitir predicciones no lineales simples, digamos $N_s = N/I \approx 500$. Ahora, para cada par de segmentos $S_i$ y $S_j$, calculemos el error cuadrático medio de predicción usando $S_i$ como conjunto de entrenamiento (es decir, encontrar vecinos en $S_i$) pero utilizando $S_j$ como conjunto de prueba (es decir, sobre el cual se realizan las predicciones). El error de predicción cruzada como función de $i$ y $j$ revela entonces qué segmentos difieren en su dinámica. Si el error de predicción para un par de segmentos $S_i, \; S_j$ es mayor que el promedio, los datos en $S_i$ obviamente proporcionan un mal modelo para los datos en $S_j$. Esperamos que las entradas diagonales $i = j$ sean sistemáticamente más pequeñas: representan errores \textit{in-sample} puesto que el conjunto de entrenamiento y el de prueba son idénticos. El uso de predicciones mutuas como prueba de estacionariedad fue desarrollado en Schreiber (1997). El paquete TISEAN contiene un programa llamado \texttt{nstatc} que implementa este método.

\bigskip

\textbf{Ejemplo 4.3 (No estacionariedad de la tasa de respiración)} El conjunto de datos introducido en el Apéndice B.5 contiene un registro de la frecuencia respiratoria de un paciente humano durante casi toda una noche (aproximadamente 5 horas). Por supuesto, las condiciones probablemente cambiaron a lo largo de ese tiempo. Como primer enfoque a la cuestión de qué episodios de la señal muestran dinámicas similares, dividimos la grabación en 34 segmentos no solapados $S_i$ de 1000 puntos (500 s) cada uno. Luego utilizamos cada segmento, por turno, como base de datos para predicciones con \texttt{zeroth}. Usamos un retardo unitario y embebimientos tridimensionales. El tamaño de los vecindarios se fijó en un cuarto de la varianza de los valores en la base de datos. Para cada uno de los 34 conjuntos de datos, calculamos 34 errores de predicción, prediciendo valores en el segmento $S_j$. En la Fig. 4.3 mostramos el error cuadrático medio de predicción como una superficie sobre el plano $i-j$. Excepto por los errores esperadamente más bajos en $i = j$, vemos que existe una transición alrededor de la mitad de la grabación: los segmentos de la primera mitad son inútiles para predecir en la segunda mitad, y viceversa. La idea de dividir un conjunto de datos en partes y considerar las fluctuaciones de una cierta estadística calculada en cada parte está muy extendida, y la única cuestión es qué estadística utilizar. Hemos elegido la predictibilidad no lineal, ya que es muy apropiada en el caso en que los datos poseen fuertes correlaciones no lineales, y produce resultados significativos incluso para datos de procesos estocásticos. La principal ventaja de esta elección radica en que permite calcular correlaciones cruzadas como en la Fig. 4.3. Las variaciones de los errores de predicción cruzada son mucho mayores y, por tanto, más significativas que las fluctuaciones de los errores in-sample, que aparecen en la diagonal de la Fig. 4.3. Volveremos sobre este punto en el Capítulo 13.

\newpage

\section{Reducción no lineal simple de ruido}

Un pariente cercano de la predicción es la reducción de ruido. En lugar de predecir valores futuros, ahora queremos reemplazar mediciones ruidosas por valores mejores que contengan menos ruido. Para la predicción, no teníamos información sobre la cantidad a pronosticar aparte de las mediciones precedentes. Para la reducción de ruido, estamos en una mejor situación a este respecto. Tenemos una medición ruidosa desde la cual empezar y, a menos que debamos procesar los datos inmediatamente en tiempo real, también disponemos de los valores futuros. Por otro lado, esto proporciona un criterio de referencia: los reemplazos que propongamos deben contener errores que, en promedio, sean menores que la amplitud inicial del ruido, de lo contrario todo el procedimiento carece de valor. La reducción de ruido significa que se intenta descomponer cada valor de una serie temporal en dos componentes, uno que supuestamente contiene la señal y otro que contiene fluctuaciones aleatorias. Así, siempre asumimos que los datos pueden concebirse como una superposición aditiva de dos componentes diferentes que deben distinguirse mediante algún criterio objetivo. La herramienta estadística clásica para obtener esta distinción es el espectro de potencia. El ruido aleatorio tiene un espectro plano, o al menos uno amplio, mientras que las señales periódicas o cuasiperiódicas tienen líneas espectrales agudas. Una vez que ambos componentes han sido identificados en el espectro, un filtro de Wiener (Sección 2.4) puede usarse para separar la serie temporal en consecuencia.

\bigskip

Este enfoque falla para dinámicas caóticas deterministas porque la salida de tales sistemas generalmente conduce a espectros de banda ancha y, por lo tanto, posee propiedades espectrales que suelen atribuirse al ruido aleatorio. Incluso si partes del espectro pueden asociarse claramente con la señal, una separación entre señal y ruido fracasa en la mayoría de las partes del dominio de frecuencias. Los sistemas dinámicos caóticos deterministas son de particular interés porque el determinismo ofrece un criterio alternativo para distinguir la señal y el ruido (lo cual, por supuesto, no es determinista). La manera en que explotaremos la estructura determinista seguirá de cerca lo que hemos hecho para las predicciones no lineales. Para mayor comodidad, planteemos el problema directamente en coordenadas de retardo. Supongamos que la evolución temporal de la señal es determinista mediante un mapa $f$ que desconocemos. Todo lo que sabemos proviene de mediciones ruidosas de esta señal:  

\begin{equation}
s_n = x_n + \eta_n, \qquad x_n = f(x_{n-m}, \ldots, x_{n-1}) .
\end{equation}

Aquí $\eta_n$ se supone como un ruido aleatorio con autocorrelaciones que decaen rápidamente y sin correlaciones con la señal $x_n$. Nos referiremos a $\sigma = \sqrt{\langle \eta^2 \rangle}$ como la amplitud del ruido o nivel absoluto de ruido. ¿Cómo podemos recuperar la señal $x_n$ dada únicamente la secuencia ruidosa $s_n$? Tras lo que hemos aprendido sobre predicciones, a primera vista la solución parece obvia: para limpiar un valor particular de la serie temporal, eliminamos las mediciones y las reemplazamos por una predicción $\hat{x}_n$ basada en mediciones previas, es decir:  

\begin{equation}
\hat{x}_n = \hat{f}(s_{n-m}, \ldots, s_{n-1}) .
\end{equation}

Si se tiene un carácter más conservador, puede preferirse usar alguna combinación lineal de las mediciones ruidosas y sus predicciones. Pero incluso entonces, este esquema fallará para sistemas caóticos. Recordemos que las mediciones previas también son ruidosas, y los errores que contienen serán amplificados por la dinámica caótica. El reemplazo propuesto será en promedio mucho peor que las mediciones ruidosas originales. (¡Cuán mucho peores depende de la tasa de expansión local promedio de la dinámica —esperar un factor de 1.5 a 3.0 para mapas típicamente caóticos! Véase la Sección 10.2 para una discusión más cuantitativa). Finalmente, lo que funciona es resolver la relación implícita 

\begin{equation}
x_n - f(x_{n-m}, \ldots, x_{n-1}) = 0
\end{equation}

para una de las coordenadas intermedias, digamos $x_{n-m/2}$, para un valor par de $m$. Esto estabiliza el reemplazo tanto en las direcciones estables como en las inestables. Esto se ilustra en la Fig. 4.4. Por supuesto, no conocemos la función $f$. Incluso si la conociéramos, no podríamos resolverla en general para uno de sus argumentos. Por lo tanto, elegimos una aproximación muy tosca por ahora y reemplazamos la función $f$ (resuelta para un argumento) por una función localmente constante. Esto es exactamente lo que hicimos en el método de predicción no lineal simple descrito antes. En otras palabras, para obtener una estimación $\hat{x}_{n_0-m/2}$ del valor de $x_{n_0-m/2}$, formamos vectores de retraso 

\[
\mathbf{s}_n = (s_{n-m+1}, \ldots, s_n)
\]

y determinamos aquellos que son cercanos a $\mathbf{s}_{n_0}$. El valor promedio de $s_{n-m/2}$ es entonces usado como el valor corregido $\hat{x}_{n_0-m/2}$: 

\begin{equation}
\hat{x}_{n_0-m/2} = \frac{1}{|U_\epsilon(\mathbf{s}_{n_0})|} \sum_{\mathbf{s}_n \in U_\epsilon(\mathbf{s}_{n_0})} s_{n-m/2}.
\end{equation}

De nuevo, por simplicidad numérica usamos la norma máxima. Solo cuando la dimensión de incrustación es realmente grande (digamos, $m > 20$) y queremos tratar datos con alta amplitud de ruido, una norma euclídea es más útil para la identificación de aquellos vecinos en el conjunto de datos ruidosos que también serían buenos vecinos para los datos libres de ruido. Debido a la supuesta independencia entre señal y ruido,

\begin{equation}
(\mathbf{s}_n - \mathbf{s}_k)^2 = \sum_{l=0}^{m-1} (s_{n+l} - s_{k+l})^2 = \sum_{l=0}^{m-1} (x_{n+l} + \eta_{n+l} - x_{k+l} - \eta_{k+l})^2 \approx (\mathbf{x}_n - \mathbf{x}_k)^2 + 2m\sigma^2,
\end{equation}

donde la aproximación mejora cuanto mayor es $m$. Así, la distancia euclídea al cuadrado de dos vectores de retraso ruidosos es idéntica a la distancia euclídea al cuadrado de los vectores libres de ruido, más una constante. Por lo tanto, los vectores ruidosos más cercanos a $\mathbf{s}_n$ son, en buena aproximación, exactamente aquellos cuyos contrapartes sin ruido también serían los más cercanos a la versión libre de ruido de $\mathbf{s}_n$, y estos son los que necesitamos para lograr una reducción de ruido exitosa.

\bigskip

No sorprendentemente, la Ecuación (4.11) es muy similar a la (4.7). A diferencia del caso de predicción, aquí $U_\epsilon(s_{n_0})$ nunca está vacío, sin importar cuán pequeño sea el valor de $\epsilon$ que elijamos: siempre contiene al menos a $s_{n_0}$. Esto es bueno de saber ya que tenemos que hacer alguna elección de $\epsilon$ al usar este algoritmo. Está garantizado que si elegimos un $\epsilon$ demasiado pequeño, lo peor que puede ocurrir es que el único vecino encontrado sea $s_{n_0}$ mismo. Esto, sin embargo, da la estimación $\hat{x}_{n_0-m/2} = s_{n_0-m/2}$, lo cual simplemente significa que no se realiza ninguna corrección en absoluto. Hemos iniciado una discusión importante. ¿Cómo podemos asegurarnos de que todo lo que hacemos es reducir el ruido sin distorsionar la señal? Dado que no conocemos la señal limpia, todo lo que podemos pedir en la práctica es que el efecto de cualquier posible distorsión sea más aceptable que el efecto del ruido anterior. En la mayoría de los casos, esto significará que el error debido a la distorsión debe ser considerablemente menor que la amplitud del ruido. Obviamente, será necesario adquirir experiencia con la reducción de ruido no lineal utilizando ejemplos donde la señal verdadera sea conocida. Como esto suele ser el caso únicamente con datos numéricos generados en computadora, dichos datos predominan en la literatura original. Estos estudios han demostrado que el esquema no lineal simple, así como esquemas más avanzados que describiremos en la Sección 10.3, son de hecho capaces de reducir ruido en datos deterministas sin causar distorsiones indebidas. Antes de limpiar datos experimentales, el lector más escéptico puede desear probar el algoritmo en datos artificiales de propiedades similares para ganar confianza. Además, recomendamos como prueba mínima de fiabilidad estudiar las propiedades estadísticas de lo que el algoritmo ha eliminado como ruido, es decir, de la serie temporal de diferencias entre señal ruidosa y señal después de la reducción de ruido, $\hat{\eta}_n = s_n - \hat{x}_n$. Su función de autocorrelación debería decaer rápidamente, no debería haber correlaciones cruzadas significativas entre $\hat{\eta}_n$ y la señal $\hat{x}_n$, y su distribución marginal debería cumplir con las expectativas de lo que el ruido de medición aparenta (en particular, su desviación estándar no debería exceder la amplitud de ruido asumida). Estudios con señales verdaderas conocidas sugieren que una buena elección para el tamaño de los vecindarios está dada por 2--3 veces la amplitud de ruido. Si bien en algunos casos es fácil estimar el nivel de ruido aproximado a partir de un gráfico de los datos, también podemos estimarlo utilizando la suma de correlaciones, como se discutirá en la Sección 6.7. Si no estamos seguros sobre el nivel de ruido, deberíamos subestimarlo. Como mencionamos antes, esto en el peor de los casos debilitará el rendimiento, pero no introducirá artefactos de filtrado.

\bigskip

Para la dimensión de incrustación $m$ parece razonable usar un valor tal que la dinámica sea determinista tanto en las coordenadas futuras como en las pasadas. Aquí, la elección conservadora es un valor mayor de $m$, con lo cual eliminamos algunos de los vecinos más dudosos. Dentro del paquete TISEAN, el esquema descrito aquí está implementado como el programa \texttt{lazy}, que se discute en el Apéndice A.8.

\newpage

\textbf{Ejemplo 4.4 (Reducción no lineal simple de ruido en datos de láser NMR).} Como tributo a aquellos científicos experimentales muy razonables que se niegan a convencerse solo porque un método funciona en el mapa de Hénon, el conjunto de datos que elegimos como nuestro primer ejemplo es al menos en parte real. La señal “limpia” se toma como 5000 valores de los datos del láser NMR que mostramos antes (véase el Apéndice B.2 para más detalles). Contienen aproximadamente un 1.1\% de ruido de medición. Ahora añadimos artificialmente “ruido de medición” consistente en números aleatorios que han sido filtrados para tener el mismo espectro de potencia que los propios datos. La amplitud rms del ruido se elige como el 10\% de la amplitud de los datos. Este tipo de ruido se llama \textit{ruido en banda} y es completamente indistinguible de la señal para cualquier filtro espectral. Ahora usamos el esquema de reducción no lineal simple de ruido para eliminar de nuevo el ruido de medición. El resultado se muestra en la Fig. 4.5. Por supuesto, nuestro algoritmo no podrá distinguir entre el ruido real y el adicional, lo que limita la evaluación de su rendimiento. Aun así, podemos reducir la amplitud del componente aleatorio añadido en un factor de aproximadamente 1.7, usando valores de $\epsilon$ correspondientes a unas 2–3 veces la amplitud del ruido.

\bigskip

\textbf{Ejemplo 4.5 (Reducción no lineal simple de ruido en datos de láser infrarrojo lejano).} Pasemos ahora a un ejemplo realmente real. En la Fig. 4.6 (panel izquierdo) mostramos una ampliación por un factor lineal de dos de un retrato de fase del conjunto de datos A de la competencia del Instituto Santa Fe; véase Weigend \& Gershenfeld (1993). Los datos (Apéndice B.1) representan la intensidad fluctuante de un láser infrarrojo lejano y pueden describirse aproximadamente mediante tres ecuaciones diferenciales ordinarias acopladas (Hübner et al. (1993)). Dado que los nuevos datos son números enteros de 8 bits, el error de medición es al menos el error de discretización, que puede suponerse distribuido uniformemente en el intervalo $[-0.5,0.5]$. La reducción de ruido valdría la pena incluso si pudiéramos eliminar solo este error y así aumentar el rango dinámico de los datos. Aplicamos el algoritmo de reducción de ruido suponiendo una amplitud de ruido (absoluta) de una unidad y elegimos $\epsilon = 2$, es decir, definimos como vecinos aquellos puntos que no están más alejados de dos unidades en cada una de las siete direcciones del espacio de incrustación. El conjunto de datos resultante se muestra en la Fig. 4.6 (panel derecho) con la misma magnificación que los datos originales. Todo lo que podemos decir en este punto es que la estructura en los datos parece más ordenada después de la reducción de ruido. Esta es, por supuesto, una afirmación insatisfactoria, pero necesitaremos herramientas mucho más potentes para verificar y cuantificar la reducción de ruido en datos reales. En particular, necesitaremos una estimación del nivel de ruido antes y después de la reducción de ruido. Tal estimación es posible sobre la base de la integral de correlación (Capítulo 6, y Sección 10.2).

\bigskip

Hagamos una observación sobre la implementación práctica de este sencillo método de reducción no lineal de ruido. Una codificación directa de la Ec. (4.11) resultaría en un algoritmo que consume mucho tiempo, ya que queremos limpiar todos los puntos y no solo unos pocos. La principal carga computacional consiste en formar los vecindarios-$\epsilon$ para cada uno de los puntos que deseamos corregir. Podemos ahorrar órdenes de magnitud en tiempo de cómputo si aplicamos un algoritmo rápido de búsqueda de vecinos en lugar de probar todos los pares posibles y descartar aquellos que no están lo suficientemente cerca. De hecho, reducimos el recuento asintótico de operaciones de $O(N^2)$ a $O(N \ln N)$ si usamos un árbol binario, o incluso a $O(N)$ si utilizamos un método asistido por cajas. Schreiber (1995) describe en detalle cómo se hace esto. El paquete TISEAN utiliza métodos de búsqueda de vecinos asistidos por cajas siempre que sea apropiado, donde se encuentran implementaciones de este y de esquemas de reducción de ruido más refinados.

\bigskip

\section*{Lecturas recomendadas}

Fuentes generales para el problema de la predicción de series de tiempo son los volúmenes editados por Casdagli \& Eubank (1992) y Weigend \& Gershenfeld (1993). Artículos clásicos en el campo incluyen a Farmer \& Sidorowich (1987, 1988), Casdagli (1989) y Sugihara \& May (1990). Referencias más específicas sobre métodos avanzados de predicción se darán en el Capítulo 12. El método simple de reducción de ruido presentado en la Sección 4.5 fue propuesto por Schreiber (1993).

\newpage

\section{Ejercicios}

\begin{enumerate}
\item Pruebe el programa \texttt{zeroth} (véase Sección A.4.1) en algunos conjuntos de datos artificiales. Comience con los dos conjuntos de datos linealmente no correlacionados creados en el Ejercicio 2.2. ¿Cómo dependen los errores de predicción del tiempo de predicción en los dos casos?
\newpage
\item Pruebe el programa anterior en 5000 iteraciones del mapa de Hénon (Ejemplo 3.2), añadiendo diferentes cantidades de ruido. Compare el error de predicción con el caso (artificial) en que usa el propio mapa de Hénon para hacer predicciones. Como referencia, en caso de que el resultado le parezca desconcertante: obtenemos un error medio de predicción de aproximadamente 0.1 para un ruido gaussiano de amplitud 0.05 si elegimos vecindarios de radio entre 0.05 y 0.15.
\end{enumerate}



\newpage

\chapter{Inestabilidad: Exponentes de Lyapunov}

\section{Dependencia sensible de las condiciones iniciales}

La característica más destacada del caos es la impredecibilidad del futuro, a pesar de una evolución temporal determinista. Esto ya se había evidenciado en la Fig. 1.2: el error promedio cometido al predecir el resultado de una medición futura aumenta rápidamente con el tiempo, y en este sistema la capacidad de predicción se pierde casi por completo después de solo 20 pasos temporales. Sin embargo, afirmamos que estos datos experimentales pueden describirse muy bien como un sistema determinista de baja dimensión. ¿Cómo podemos explicar esta aparente contradicción?

\bigskip

\textbf{Ejemplo 5.1 (Divergencia de trayectorias del láser NMR)} En la Fig. 5.1 mostramos varios segmentos de la serie temporal del láser NMR (los mismos datos usados en la Fig. 1.2; véase el Apéndice B.2) que inicialmente están muy próximos entre sí. Con el paso del tiempo, estos segmentos se separan y finalmente se vuelven no correlacionados. Por lo tanto, es imposible predecir la posición de la trayectoria más allá de, digamos, diez pasos temporales, incluso conociendo la posición de otra trayectoria que en ese momento inicial estaba extremadamente cerca. (Esto está muy en el espíritu del esquema de predicción de la Sección 4.2.)

\bigskip

El ejemplo anterior ilustra que nuestra experiencia cotidiana “causas similares producen efectos similares” no es válida en sistemas caóticos, excepto durante períodos cortos, y que solo una reproducción matemáticamente exacta de algún evento produciría el mismo resultado debido al determinismo. Nótese que esto no tiene nada que ver con alguna influencia no observada desde fuera del sistema (aunque en los datos experimentales siempre está presente), sino que puede encontrarse en cualquier modelo matemático de un sistema caótico. Esta impredecibilidad es una consecuencia de la inestabilidad inherente de las soluciones, reflejada en lo que se denomina \textit{dependencia sensible de las condiciones iniciales}. Las pequeñas desviaciones entre las “condiciones iniciales” de todas las trayectorias mostradas en la Fig. 5.1 se amplifican después de unos pocos pasos de tiempo, y cada detalle no observado del estado en el tiempo $n_0$ es importante para el recorrido preciso de la trayectoria en el espacio de fases.

\bigskip

Una investigación más cuidadosa de esta inestabilidad conduce a dos conceptos diferentes, aunque relacionados. Un aspecto, que no queremos desarrollar en la primera parte del libro, es la pérdida de información relacionada con la impredecibilidad. Esto se cuantifica mediante la \textit{entropía de Kolmogorov–Sinai} y será introducido en el Capítulo 11. El otro aspecto es de naturaleza geométrica simple, a saber: que trayectorias cercanas se separan muy rápidamente o, más precisamente, de manera exponencial con el tiempo.

\newpage

\section{Divergencia exponencial}

El hecho de que las trayectorias diverjan con el paso del tiempo no sería en sí mismo muy dramático si ocurriera de forma muy lenta, como es típico en los sistemas predominantemente periódicos. Por tanto, hablamos de caos únicamente si esta separación ocurre de manera exponencialmente rápida. El exponente promediado adecuadamente de este crecimiento es característico del sistema subyacente a los datos y cuantifica la intensidad del caos. A este se le llama el \textit{exponente de Lyapunov}\footnote{Si las tasas de divergencia se promedian en intervalos de tiempo muy cortos, se pueden definir \textit{exponentes locales de Lyapunov} que permiten un análisis resuelto en tiempo o espacio de fases. Sin embargo, estos exponentes locales no son invariantes simples y los resultados suelen ser difíciles de interpretar. En cualquier caso, conviene recordar que el exponente de Lyapunov global es un promedio sobre una cantidad que puede mostrar fuertes fluctuaciones.}.  

\bigskip

\textbf{Ejemplo 5.2 (Separación exponencial de trayectorias, datos del láser NMR).} 
La Fig. 5.2 muestra otra perspectiva de lo que se demostró en la Fig. 5.1: las distancias escalares entre los elementos de segmentos inicialmente cercanos de la serie temporal del láser NMR (Apéndice B.2) se representan en una escala logarítmica como función del tiempo. En primer lugar, casi todas las curvas son más o menos paralelas. Esto confirma que estos datos experimentales realmente obedecen a una evolución temporal determinista. Aparte de una gran cantidad de estructura no explicada, se observa claramente un aumento lineal general en este gráfico semilogarítmico.  

\bigskip

En el Capítulo 11 veremos que se pueden definir tantos exponentes de Lyapunov para un sistema dinámico como dimensiones tenga su espacio de fases. Aquí queremos restringirnos al más importante de ellos, el \textit{exponente de Lyapunov máximo} $\lambda$. Sean $\mathbf{s}_{n_1}$ y $\mathbf{s}_{n_2}$ dos puntos en el espacio de fases cuya distancia $\|\mathbf{s}_{n_1} - \mathbf{s}_{n_2}\| = \delta_0 \ll 1$. Denotemos por $\delta_{\Delta n}$ la distancia, un tiempo $\Delta n$ después, entre las dos trayectorias que emergen de estos puntos: $\delta_{\Delta n} = \|\mathbf{s}_{n_1+\Delta n} - \mathbf{s}_{n_2+\Delta n}\|$. Entonces, $\lambda$ se determina mediante la relación  
\begin{equation}
\delta_{\Delta n} \simeq \delta_0 e^{\lambda \Delta n}, \qquad \delta_{\Delta n} \ll 1, \qquad \Delta n \gg 1.
\end{equation}
Si $\lambda$ es positivo, esto significa una divergencia exponencial de trayectorias cercanas, es decir, caos. Naturalmente, dos trayectorias no pueden separarse más allá del tamaño del atractor, por lo que la ley de la Ecuación (5.1) solo es válida durante tiempos $\Delta n$ para los cuales $\delta_{\Delta n}$ se mantiene pequeña. De lo contrario, se observa una violación de la ecuación (5.1) en forma de una saturación de la distancia. Debido a este hecho, una definición matemáticamente más rigurosa deberá implicar un primer límite $\delta_0 \to 0$, de modo que se pueda tomar un segundo límite $\Delta n \to \infty$ sin involucrar efectos de saturación. Solo en este segundo límite el exponente $\lambda$ se convierte en una cantidad bien definida e invariante. Volveremos sobre esto más adelante, en particular en la Sección 11.2.

\bigskip

En los sistemas disipativos también puede encontrarse un exponente de Lyapunov máximo negativo, lo que refleja la existencia de un punto fijo estable. Dos trayectorias que se aproximan al punto fijo también se acercan entre sí de manera exponencialmente rápida. Si el movimiento converge hacia un ciclo límite, dos trayectorias sólo pueden separarse o aproximarse entre sí más lentamente que de forma exponencial. En este caso, el exponente de Lyapunov máximo es cero y el movimiento se denomina \textit{marginalmente estable}. Si un sistema predominantemente determinista es perturbado por ruido aleatorio, en las escalas pequeñas puede caracterizarse mediante un proceso de difusión, con $\delta_{\Delta n}$ creciendo como $\sqrt{\Delta n}$. Por tanto, el exponente de Lyapunov máximo es infinito.\footnote{A veces se dice que el exponente del movimiento difusivo es cero, porque las distancias entre dos trayectorias que comienzan en el mismo punto crecen como $\sqrt{t}$. Sin embargo, la naturaleza afín de las trayectorias brownianas sugiere observar el comportamiento en el límite de tiempos cortos, donde la derivada de $\sqrt{t}$ respecto a $t$ diverge. Además, $\lambda = \infty$ para el movimiento browniano concuerda con el hecho de que la entropía de un proceso estocástico también es infinita.} De acuerdo con la definición matemática, esto es cierto sin importar cuán pequeño sea el componente de ruido. Si asumimos que sólo existe un error aditivo de medición, podemos intentar estimar el exponente de Lyapunov de un sistema determinista subyacente asumido (véase la Tabla 5.1 para un resumen). Los exponentes de Lyapunov tienen unidades de tiempo inverso y proporcionan una escala temporal típica para la divergencia o convergencia de trayectorias cercanas. Al calcular valores numéricos, se debe tener en cuenta la normalización adecuada: ¿queremos computar esta escala de tiempo en unidades del índice temporal de las mediciones o en unidades de tiempo real, por ejemplo, en segundos? Cuando los datos se transforman en datos tipo mapa mediante el método de la superficie de sección de Poincaré, es necesario conocer el tiempo promedio entre dos puntos sucesivos del mapa de Poincaré, $\tau_P$. Entonces, los exponentes de Lyapunov calculados para los datos del mapa están relacionados con los de los datos del flujo mediante
\begin{equation}
\lambda_{\text{map}} = \lambda_{\text{flow}} \, \tau_P.
\end{equation}

\newpage

En los párrafos anteriores hablábamos del exponente de Lyapunov del sistema como una cantidad característica. Pero el sistema solo está representado por las mediciones, y podríamos realizar las mediciones sobre el mismo sistema de otra manera. Más aún, podríamos reescalar, desplazar o procesar de otra forma los datos, y tenemos mucha libertad en la elección del método preciso de reconstrucción del espacio de estados. ¿Cómo afectaría el exponente de Lyapunov a todas estas modificaciones? La respuesta es simplemente que los exponentes de Lyapunov no dependen de ellas en absoluto. Los exponentes de Lyapunov son \textit{invariantes} bajo todas estas transformaciones, siempre que sean suaves. La razón básica es que describen el comportamiento a largo plazo. Cualquier reparametrización suave e invertible del espacio de fases solo puede modificar las proporciones de distancia por factores finitos. Tales factores desaparecen en el límite cuando $\Delta n \to \infty$. Esto justifica nuestro interés en este número, ya que todos los demás obtendrán exactamente los mismos resultados con sus propias observaciones, siempre que el sistema físico subyacente sea realmente el mismo. 

\bigskip

Ahora bien, puede que no estés convencido por la Fig. 5.2 y te preguntes si realmente existe un crecimiento exponencial. De hecho lo hay, pero no hay que olvidar que la Ec. (5.1) es válida en el verdadero espacio de estados del sistema. El proceso de medición es una proyección, y la proyección de las distancias puede reducirse simplemente debido al ángulo de observación. Ya en una incrustación retardada, la figura análoga se ve mucho más clara. Además, la tasa local de divergencia varía a lo largo del atractor y el verdadero exponente es un promedio apropiado sobre lo que se observa en la Fig. 5.2.

\section{Medición del exponente máximo a partir de datos}

Dado que un exponente de Lyapunov máximo positivo es una fuerte señal de caos, resulta de gran interés determinar su valor para una serie temporal dada. El primer algoritmo con este propósito fue propuesto por \textit{Wolf et al.} (1985). Desafortunadamente, su uso requiere mucho cuidado, ya que es fácil obtener resultados erróneos. Este algoritmo no permite verificar la presencia de divergencia exponencial, sino que simplemente asume su existencia y, por lo tanto, produce un exponente finito incluso para datos estocásticos, donde el verdadero exponente es infinito. Aunque el algoritmo de Wolf no utiliza más que una reconstrucción por retardos del espacio de fases, existe otra clase de algoritmos que también implican una aproximación de la dinámica determinista subyacente. Las referencias relevantes son \textit{Sano \& Sawada} (1985) y \textit{Eckmann et al.} (1986). Este método es muy eficiente siempre que los datos permitan una buena aproximación de la dinámica. Será explicado con mayor detalle en el Capítulo 11.Aquí queremos describir un algoritmo introducido recientemente por \textit{Rosenstein et al.} (1993) y por \textit{Kantz} (1994) de manera independiente. Este algoritmo prueba directamente la divergencia exponencial de trayectorias cercanas, y por tanto nos permite decidir si realmente tiene sentido calcular un exponente de Lyapunov para un conjunto de datos dado.

\bigskip

En la Fig. 5.2 observamos un crecimiento exponencial promedio superpuesto con fuertes fluctuaciones. Estas fluctuaciones tienen su origen en diferentes aspectos de la dinámica determinista. Como ya se mencionó, la proyección involucrada en la medición puede hacer que las distancias parezcan contraerse temporalmente, aunque crezcan en el verdadero espacio de estados. Además, en el verdadero espacio de estados, las distancias no crecen en todas partes del atractor con la misma tasa, y localmente incluso pueden disminuir. Por lo tanto, el exponente de Lyapunov es un promedio de estas \textit{tasas locales de divergencia} a lo largo de todos los datos. Finalmente, los datos experimentales están generalmente contaminados por ruido. Su influencia puede minimizarse utilizando estadísticas de promediado adecuadas al calcular el exponente. Para obtener el crecimiento exponencial promedio de las distancias, se debe proceder de la siguiente manera. Se elige un punto $s_{n_0}$ de la serie temporal en el espacio embebido y se seleccionan todos los vecinos con distancia menor que $\epsilon$. Se calcula el promedio sobre las distancias de todos los vecinos al punto de referencia de la trayectoria como función del tiempo relativo $\Delta n$, es decir, el promedio (aritmético) sobre las curvas mostradas en la Fig. 5.2. El logaritmo de la distancia promedio en el tiempo $\Delta n$ es una tasa de expansión efectiva sobre el intervalo $\Delta n$ (más el logaritmo de la distancia inicial), conteniendo todas las fluctuaciones deterministas debidas a la proyección y a la dinámica. Repitiendo esto para muchos valores de $n_0$, las fluctuaciones de las tasas de expansión efectivas se promediarán. Así, se debe calcular:

\begin{equation}
S(\Delta n) = \frac{1}{N} \sum_{n_0=1}^{N} \ln \left( \frac{1}{|\mathcal{U}_\epsilon(s_{n_0})|} \sum_{s_n \in \mathcal{U}_\epsilon(s_{n_0})} |s_{n_0+\Delta n} - s_{n+\Delta n}| \right).
\end{equation}

Los puntos de referencia $s_{n_0}$ son vectores embebidos como de costumbre; $\mathcal{U}_\epsilon(s_{n_0})$ es la vecindad de $s_{n_0}$ con diámetro $\epsilon$. Nótese que el último elemento de $s_{n_0}$ es $s_{n_0}$, por lo tanto, $s_{n_0+\Delta n}$ está fuera del intervalo cubierto por el vector de retardo $s_{n_0}$. Como a priori no se conocen ni la dimensión mínima de embebido $m$ ni la distancia óptima $\epsilon$, se debe calcular $S(\Delta n)$ para varios valores de ambos parámetros. El tamaño de la vecindad debe ser lo más pequeño posible, pero lo suficientemente grande como para que en promedio cada punto de referencia tenga al menos unos pocos vecinos. De lo contrario, se podrían omitir partes del atractor y así calcular un valor erróneo (véase el Ejemplo 5.4). El programa \texttt{lyap\_k} del paquete \texttt{TISEAN} (véase la Sección A.5) realiza los cálculos necesarios para generar las curvas $S(\Delta n)$. Si para cierto rango de $\Delta n$ la función $S(\Delta n)$ muestra un incremento lineal \textit{robusto}, su pendiente es una estimación del exponente de Lyapunov máximo $\lambda$ por paso temporal, el cual debe convertirse a las unidades deseadas. Enfatizamos la palabra \textit{robusto}, ya que encontrar un aumento lineal para valores aislados de $m$ y $\epsilon$ no puede considerarse una firma positiva de divergencia exponencial. A continuación, se presentarán algunos resultados típicos que ilustrarán este punto. A veces se encuentran incrementos lineales para ciertos valores de $\epsilon$ y $m$, pero sin líneas rectas para otros casos. Antes de citar un valor numérico para $\lambda$, se debe intentar comprender las desviaciones de las curvas respecto al comportamiento lineal (ruido, dimensión de embebido demasiado pequeña, pocos vecinos, saturación debido a $\epsilon$ demasiado grande, separación tipo ley de potencias, etc.).

\bigskip

\textbf{Ejemplo 5.3 (Exponente de Lyapunov de los datos del láser NMR).} Nuestros resultados para los datos del láser NMR (Apéndice B.2) se muestran en la Fig.~5.3. Dado que este es un resultado muy típico y, según nuestro criterio, una forma razonable de representar $S$, dediquemos un poco más de tiempo a analizarlo. Las curvas se muestran para dimensiones de incrustación $m = 2$ (panel superior) y $m = 3, 4, 5$ (panel inferior). Usamos $\epsilon = 15, 30, 60, 120,$ y $240$. Representamos $e^{S(\Delta n)}$ en una escala logarítmica (lo cual equivale a representar $S(\Delta n)$ en una escala lineal) con el fin de poder usar unidades físicas en el eje. La base de datos consistió en 38,000 mediciones escalares, pero para cada curva extendimos la primera suma de la Ec.~(5.2) solo lo suficiente como para que al menos 500 puntos de referencia tuvieran más de 10 vecinos. Las curvas con $m = 2$ (panel superior) presentaron un salto bastante pronunciado desde $\Delta n = 0$ hasta $\Delta n = 1$. Esto significa que $m = 2$ no produce una incrustación adecuada. Esto se confirma observando nuevamente el atractor en la Fig.~3.6. En una representación bidimensional, las diferentes ramas del atractor se intersectan. Naturalmente, las imágenes de puntos cercanos en las proximidades de estas intersecciones pueden estar muy alejadas entre sí. Esta es una consecuencia de una proyección no resuelta. Pasemos ahora a discutir el panel inferior. Si observas con atención, verás cinco grupos de curvas, cada uno correspondiente a uno de los cinco valores de $\epsilon$ que pueden verse para $\Delta n = 0$, pero debido a los efectos que se describirán más adelante, los tres grupos inferiores se superponen fuertemente. Cada grupo individual representa los valores $m = 3, 4, 5$, lo que muestra que el resultado no depende de la dimensión de incrustación tan pronto como $m$ es lo suficientemente grande (podríamos perder solo unos pocos vecinos en la vecindad inicial al pasar de $m$ a $m + 1$). Para $\Delta n$ positivos crecientes, $S(\Delta n)$ crece hasta que se satura alrededor de un valor $e^{S(\Delta n \rightarrow \infty)} \approx 2000$, que corresponde aproximadamente a la distancia media absoluta entre dos vectores de incrustación arbitrarios en el atractor. En realidad, vemos que el rango de escalamiento termina antes, en $e^S \approx 1000$. Por lo tanto, no deberíamos usar un valor de $\epsilon$ demasiado grande al estimar el exponente de Lyapunov. La información más útil proviene de las curvas intermedias, para las cuales la distancia inicial de las trayectorias es lo suficientemente pequeña como para que queden suficientes pasos de tiempo para que ocurra la divergencia exponencial antes de que los efectos de saturación terminen la parte lineal. Finalmente, ¿qué sucede con las curvas inferiores? Muestran un rango casi lineal con la misma pendiente que los otros grupos, pero con fluctuaciones mucho más fuertes debido a los vecindarios escasos. Desde $\Delta n = 0$ hasta $\Delta n = 1$, todas las curvas presentan un salto hacia arriba. Así, nuevamente tenemos el problema de los falsos vecinos más cercanos, pero ahora para todos los valores de $m$, y la explicación es diferente. Los datos son ruidosos, y si $\epsilon$ es menor que el nivel de ruido, todavía se encuentran algunos vecinos, pero la distancia real puede ser mayor que $\epsilon$. La distancia de las imágenes sería del orden del nivel de ruido, incluso si no existiera divergencia exponencial. A partir de ahí, sin embargo, la divergencia exponencial domina sobre el ruido y el aumento lineal deseado se hace visible (comparar con la Fig.~10.5: después de la reducción de ruido, estas curvas también se comportan correctamente). 

\bigskip

A partir del cálculo de \( S(\Delta n) \) hemos aprendido varias cosas: tenemos una idea aproximada del nivel de ruido, sabemos que \( m = 2 \) es demasiado pequeño para formar un buen embebido, y hemos observado un rango donde las distancias aumentan exponencialmente. Usando una regresión lineal para las partes lineales de las curvas, podemos determinar el valor numérico del exponente de Lyapunov junto con el error del ajuste, \( \lambda = 0.3 \pm 0.01 \). Una estimación más rigurosa del error estadístico y de los posibles errores sistemáticos requeriría un análisis mucho más exhaustivo.

\newpage

\textbf{Ejemplo 5.4 (Exponente de Lyapunov del láser infrarrojo lejano).} A partir de los datos de flujo del experimento de láser tipo Lorenz del PTB Braunschweig (véase el Apéndice B.1), la serie temporal es un registro de la potencia de salida de un láser caótico. Los datos presentan oscilaciones de corto plazo con una amplitud que aumenta lentamente. Cuando la amplitud alcanza un valor crítico, el sistema colapsa y las oscilaciones comienzan nuevamente. En este punto, la información de fase se pierde, y la nueva amplitud también resulta difícil de predecir. Estos son casos con la mayor inestabilidad, mientras que el aumento suave de la amplitud es bastante predecible. Debido a esta fuerte inhomogeneidad en las tasas locales de divergencia, se esperan dificultades al estimar el exponente de Lyapunov. Para evitar dificultades adicionales debidas a la discretización relativamente gruesa de los datos, utilizamos los datos con reducción de ruido producidos en el Ejemplo 4.5. Las curvas \( e^{S(\Delta n)} \) se muestran en la Fig. 5.4. A primera vista, el resultado es decepcionante. Las curvas presentan fuertes oscilaciones, y aun interpolando entre ellas, las partes lineales son menos evidentes. Sus pendientes son ligeramente diferentes para distintos valores de \( \epsilon \). Nuestra explicación es la siguiente: en primer lugar, las oscilaciones tienen claramente su origen en la periodicidad de la señal. En segundo lugar, como ya se indicó, la tasa de expansión es mayor cerca de las disrupciones y muy pequeña en otros lugares. Desafortunadamente, cuanto menor es \( \epsilon \), menos partes inestables se muestrean, ya que no se encuentran vecinos. (Verificamos esto explícitamente al graficar todos los puntos con/sin vecinos en tres dimensiones). Por lo tanto, para valores pequeños de \( \epsilon \) se calcula el promedio de \( \lambda \) no sobre la medida invariante, sino sobre un componente menos caótico de la misma. Con esta comprensión de nuestro resultado, tenemos suficiente evidencia de caos determinista, aunque debemos ser cautelosos con nuestras afirmaciones. El valor numérico de \( \lambda \) es mucho menos preciso que en el ejemplo anterior, y como estimación proponemos \( 0.025 < \lambda < 0.03 \). El problema es que el conjunto de datos contiene solo unos 20 eventos de disrupción, lo cual no es suficiente para una buena estimación de su contribución. Realizamos el mismo análisis después de aplicar una sección de Poincaré para obtener datos tipo mapa, aplicando un método avanzado de reducción no lineal de ruido y eliminando la asimetría de los datos. Esto confirma el valor citado anteriormente, aunque las curvas \( S(\Delta n) \) siguen siendo algo insatisfactorias.

\bigskip

El problema que encontramos en el último ejemplo es común en sistemas donde la inestabilidad está fuertemente localizada en el espacio de estados. Más adelante veremos que esta también es una dificultad en métodos más refinados, y no solo al calcular exponentes de Lyapunov. Además de eso, el algoritmo propuesto aquí presenta algunas deficiencias para los datos de flujo, ya que tiende a mezclar la expansión no exponencial a lo largo de la dirección marginalmente estable de un flujo con la expansión exponencial que nos interesa. Evidentemente, el límite $\Delta n \to \infty$ no puede aproximarse muy bien con este procedimiento, lo cual impediría distinguir entre ambos efectos. Para una mejor comprensión de los resultados “nulos”, recomendamos enfáticamente el Ejercicio 5.3 que se presenta más adelante. Cuando las distancias aumentan de manera difusiva, no se puede observar incrementos lineales de $S(\Delta n)$. Esto corresponde al hecho de que el exponente máximo de Lyapunov es infinito o cero. Cuando $\Delta n$ se grafica en una escala logarítmica, puede detectarse tanto un aumento de las distancias $\delta_n \propto \Delta n^{1/2}$ —típico de la difusión— como $\delta_n \propto \Delta n$, como en los sistemas periódicos estables.

\section*{Lecturas adicionales.}  
El algoritmo para calcular el exponente de Lyapunov máximo fue desarrollado por Kantz (1994) y por Rosenstein \textit{et al.} (1993). Un método relacionado fue introducido anteriormente por Kurths y Herzel (1987). El algoritmo de Wolf \textit{et al.} (1986) es bastante conocido y ha sido ampliamente utilizado, pero debido a su inestabilidad y a la imposibilidad de distinguir la divergencia exponencial de la divergencia causada por el ruido, no se recomienda su uso. Para referencias más avanzadas sobre cómo calcular más de un exponente de Lyapunov máximo, véase el Capítulo 11.

\newpage

\section{Ejercicios}

\begin{enumerate}
    \item Crea una serie temporal escalar sintética de datos de un mapa iterando el siguiente mapa tridimensional [Baier \& Klein (1990)]:
    \[
    \begin{cases}
        x_{n+1} = 1.76 - y_n^2 - 0.1z_n, \\
        y_{n+1} = x_n, \\
        z_{n+1} = y_n.
    \end{cases}
    \]
    Elige como observable el cuadrado de la coordenada \(x\). Calcula el exponente de Lyapunov máximo utilizando el método descrito en este capítulo (programa \texttt{lyap\_k} del paquete \textit{TISEAN}) para series de longitudes 500 y 5000, respectivamente, y compara los resultados con el valor exacto \(\lambda \approx 0.225\).
	\newpage
    \item Usa un generador de números aleatorios (o la utilidad \texttt{addnoise} de \textit{TISEAN}) para añadir ruido uniformemente distribuido a las series temporales anteriores (ruido de medición). Repite el análisis.
	\newpage
    \item Crea la trayectoria de un movimiento cuasi-browniano en tu computadora. Sea \(\eta_n\) la salida de un generador de números aleatorios, y define \(s_n = s_{n-1} + \eta_n\) (para obtener un movimiento browniano verdadero los números aleatorios deben ser gaussianos). Intenta calcular el “exponente de Lyapunov máximo” para esta serie temporal de longitud \(> 10\,000\) (si no introduces un retardo adicional, \(\Delta n\) debería variar hasta unos 1000). ¿Cómo interpretas el resultado? Para una mejor comprensión, grafica \(S(\Delta n)\) frente al logaritmo de \(\Delta n\).
\end{enumerate}


\chapter{Autosimilitud: dimensiones}
\section{Geometría del atractor y fractales}

En el capítulo anterior discutimos el lado dinámico del caos, que se manifiesta en la dependencia sensible de la evolución de un sistema con respecto a sus condiciones iniciales. Este extraño comportamiento en el tiempo de un sistema caótico determinista tiene su contrapartida en la geometría del conjunto en el espacio de fases formado por las trayectorias (no transitorias) del sistema, el atractor.

\bigskip

Los atractores de sistemas caóticos disipativos (el tipo de sistemas que nos interesa) generalmente tienen una geometría muy complicada, lo que llevó a llamarlos \emph{extraños}. Sin embargo, los conjuntos extraños también pueden ocurrir sin disipación en configuraciones más generales. Como ya señalamos en el Capítulo 3, un sistema descrito por ecuaciones diferenciales autónomas (un flujo) no puede ser caótico en menos de tres dimensiones. Con el mismo argumento de que las trayectorias no deben intersectarse en un sistema determinista, podemos concluir que no solo el espacio de fases, sino también el atractor de un flujo caótico debe ser de más de dos dimensiones. Sin embargo, ligeramente más de dos dimensiones es suficiente y el movimiento en un fractal de dimensiones $2 + \epsilon$ realmente puede ser caótico. Como veremos, los atractores extraños con dimensiones fraccionarias son típicos de los sistemas caóticos. Los sistemas de tipo mapa pueden mostrar caos con atractores de dimensiones mayores que dos. Se asignan dimensiones no enteras a objetos geométricos que exhiben una especie de autosimilitud y que muestran estructura a todas las escalas de longitud.

\bigskip

\textbf{Ejemplo 6.1 (Autosimilitud del atractor del láser NMR)} Tal autosimilitud se demuestra en la Fig. 6.1 para un atractor reconstruido a partir de la serie temporal del láser NMR, Apéndice B.2. Desafortunadamente, el ruido omnipresente oculta las estructuras a pequeña escala, de modo que por debajo de una cierta escala de longitud, la autosimilitud no trivial es reemplazada por la autosimilitud trivial de una distribución suave en dos dimensiones. Por esta razón, mostramos el atractor después de que se haya aplicado un algoritmo de reducción de ruido no lineal, lo que permite varios pasos de magnificación antes de que se agote la resolución de los datos. Si las figuras te impresionan tanto como a nosotros, probablemente querrás aprender más sobre la reducción no lineal de ruido en la Sección 10.3.

\bigskip

El ejemplo matemático estándar de un conjunto auto-similar es un conjunto de Cantor. Es exactamente auto-similar cuando se escala por un factor dado. Los atractores extraños se comportan de manera diferente, como vemos en la Fig. 6.1. Encontramos auto-similitud exacta solo localmente con factores de escala dependientes de la posición. Pero para fines prácticos, esto es aún más útil, ya que hay una invariancia de escala en el sentido estadístico para factores de escala arbitrarios.

\newpage

\section{Dimensión de correlación}

Existen varias maneras de cuantificar la auto-similitud de un objeto geométrico mediante una dimensión. Por supuesto, necesitamos la definición que coincida con la noción usual de dimensión cuando se aplica a objetos no fractales: un conjunto finito de puntos es de dimensión cero, las líneas tienen dimensión uno, las superficies dos, etc. Pero, por el momento, cortemos la historia y propongamos una definición que sea de particular interés en aplicaciones prácticas donde el objeto geométrico debe ser reconstruido a partir de una muestra finita de datos que probablemente contenga algunos errores también. Esta noción, llamada \textit{dimensión de correlación}, fue introducida por Grassberger y Procaccia (1983) y (1983a).

Primero definamos la suma de correlación para una colección de puntos \(\mathbf{x}_i\), en algún espacio vectorial para que sea la fracción de todos los pares posibles de puntos que están más cerca de una distancia dada \(\epsilon\) en una norma particular. La fórmula básica (que debe ser modificada para aplicaciones prácticas) es

\begin{equation}
C(\epsilon) = \frac{2}{N(N-1)} \sum_{i=1}^N \sum_{j=i+1}^N \Theta(\epsilon - || \mathbf{x}_i - \mathbf{x}_j ||),
\end{equation}

donde \(\Theta\) es la función escalón de Heaviside, \(\Theta(x) = 0 \ \text{si} \ x \leq 0\) y \(\Theta(x) = 1 \ \text{para} \ x > 0\). La suma simplemente cuenta los pares \((\mathbf{x}_i, \mathbf{x}_j)\) cuya distancia es menor que \(\epsilon\). En el límite de una cantidad infinita de datos (\(N \rightarrow \infty\)) y para valores pequeños de \(\epsilon\), esperamos que \(C\) escale como una ley de potencia, \(C(\epsilon) \sim \epsilon^D\), y podemos definir la dimensión de correlación \(D\) por

\begin{equation}
d(N, \epsilon) = \frac{\partial \ln C(\epsilon, N)}{\partial \ln \epsilon}, 
\end{equation}

\begin{equation}
D = \lim_{\epsilon \to 0} \lim_{N \to \infty} d(N, \epsilon).
\end{equation}

Podemos verificar fácilmente que esta definición da las dimensiones integrales correctas para objetos geométricos regulares. (Ver los Ejercicios 6.1 y 6.2.) Es igualmente obvio que los límites en los que tenemos que basarnos nos meterán en problemas siempre que tengamos una muestra finita en lugar de una distribución completa: \(N\) está limitado por el tamaño de la muestra, y el rango de elecciones significativas de \(\epsilon\) está limitado por la precisión de los datos y por la inevitable cercanía de los vecinos más cercanos a distancias de escala pequeña. Antes de dar más detalles sobre cómo estimar la suma de correlación y la dimensión a partir de series temporales, dejemos que algunos comentarios generales que sentimos son importantes. Más específicamente, en cuanto a los requisitos prácticos para la estimación de la dimensión, se darán a continuación.

\bigskip

Las dimensiones, los exponentes de Lyapunov, etc., son formas de cuantificar las propiedades de una señal. Cada uno de estos conceptos es una de las muchas maneras de convertir una secuencia de datos en un solo número (o unos pocos). ¿Qué tiene de especial que escojamos los exponentes de Lyapunov o las dimensiones y no otro concepto? Esta elección está guiada por la esperanza de que el resultado mejore nuestro conocimiento sobre el \textit{sistema subyacente}, más que simplemente comprimir muchas mediciones en un solo número. Esto no se puede lograr con la mayoría de las estadísticas que podemos definir ad hoc.

\bigskip

Una cantidad que pueda servir para este propósito no debe depender significativamente del procedimiento de medición, las coordenadas elegidas, etc. De acuerdo con su definición, las dimensiones y otras cantidades cumplen con este requisito de \textit{invarianza}. Si queremos comunicar si un objeto geométrico es una superficie o un volumen al dar su dimensión, este número no debe depender de los detalles del procedimiento, la resolución de los datos, etc. No podemos relajar este requisito fundamental para los objetos fractales solo porque sean menos intuitivos. Es importante señalar que, aunque la dimensión de correlación es invariante, la suma de correlación en una escala dada \(\epsilon_0\) no lo es. Por lo tanto, no podemos simplemente tomar \(d(N, \epsilon_0)\) como una estimación de \(D\). Si queremos establecer un valor \(D\) (con sus márgenes de error) como la dimensión de correlación de un sistema, debemos asegurarnos de que este valor sea robusto bajo variaciones razonables de la técnica de medición y el procesamiento de datos.

\bigskip

Como observación conceptual, insistimos en que, en primer lugar, la dimensión de correlación es una herramienta para cuantificar la auto-similitud cuando se sabe que está presente. Cuando se aplica para este propósito, proporciona un algoritmo seguro y estable. El concepto es mucho menos adecuado y debe usarse con mucha más cautela si la auto-similitud sigue siendo incierta, es decir, si no se está seguro de que los datos sean de un sistema determinista de baja dimensión.

\newpage

\section{Suma de correlación a partir de una serie temporal}

Las definiciones anteriores de la suma y dimensión de correlación implican vectores del espacio de fases como las ubicaciones de puntos en un atractor. Así, dada una serie temporal escalar, primero debemos reconstruir un espacio de fases auxiliar mediante un procedimiento de incrustación. El más popular y completamente satisfactorio para este propósito es la técnica de incrustación por retardos que hemos utilizado en varias ocasiones anteriormente (fue introducida en la Sección 3.2). Técnicamente, el procedimiento de reconstrucción implica la elección de un tiempo de retardo $\tau$. Mientras que elecciones absurdas de $\tau$ hacen inútil el algoritmo de correlación, debemos insistir en que los resultados son invariantes bajo cambios \textit{razonables} en el procedimiento de incrustación. Cómo hacer tales elecciones razonables se discute en la Sección 3.3. Un valor para el tiempo de retardo que produzca un retrato de fase convincente debería servir también para la suma de correlación.


\bigskip

Una vez que se han reconstruido los vectores de incrustación $s_n$, la estimación de la dimensión de correlación se realiza en dos pasos. El primero es determinar la suma de correlación $C(m, \epsilon)$, Ec.~(6.1), para el rango de $\epsilon$ disponible y para varias dimensiones de incrustación $m$. Luego inspeccionamos $C(m, \epsilon)$ en busca de signos de autosimilaridad. Si estas señales son lo suficientemente convincentes, podemos calcular un valor para la dimensión. Ambos pasos requieren cierto cuidado con el fin de evitar resultados erróneos o engañosos. Queremos interpretar los datos (en forma de vectores de incrustación) como una muestra aleatoria extraída de una distribución de probabilidad subyacente, la \textit{medida natural} sobre el atractor en cuestión. Es la geometría autosimilar de esta distribución lo que queremos cuantificar. Para estimar la suma de correlación de la distribución subyacente, tenemos que extrapolar información obtenida a partir de la muestra finita. El estimador directo, Ec.~(6.1), resulta estar sesgado hacia dimensiones demasiado pequeñas cuando los pares que entran en la suma no son estadísticamente independientes. Para series temporales con autocorrelaciones distintas de cero, no puede asumirse independencia. Por ejemplo, los vectores de incrustación en tiempos sucesivos suelen estar también cercanos en el espacio de fases debido a la evolución continua del tiempo. A este tipo de correlaciones las llamamos \textit{temporales}. Véase la Fig.~6.2 para una ilustración de este problema en un flujo determinista.


\bigskip

Detengámonos un poco más en la distinción entre correlaciones temporales y geométricas. Las correlaciones temporales más importantes son causadas por el hecho de que los datos cercanos en el tiempo también están cercanos en el espacio, un hecho que no solo es cierto para sistemas deterministas puramente, sino también para muchos procesos estocásticos impulsados aleatoriamente. Superficialmente, los sistemas con esta propiedad parecen 'menos aleatorios' que los independientes extraídos de una distribución, pero esto no es lo que aquí entendemos por 'determinista'. Más que continuidad en el tiempo, nos interesa la suavidad en el espacio de fases, es decir, que estados presentes similares evolucionen hacia estados similares en el futuro cercano. Desafortunadamente, no solo el ojo sino también muchos métodos cuantitativos de análisis de series temporales pueden verse sesgados por correlaciones temporales. En el caso de la dimensión de correlación, esto conduce a una subestimación seria a menos que se tenga el cuidado necesario. Incluso señales estocásticas de dimensión infinita pueden conducir a estimaciones finitas (y aún de baja dimensión), como ha sido demostrado, por ejemplo, por Theiler (1986) y (1991). Es muy probable que la mayoría de las estimaciones de dimensión publicadas para mediciones de campo sean demasiado bajas, porque confunden la coherencia temporal con estructura geométrica. Más adelante daremos algunos ejemplos de sumas de correlación espurias.

\bigskip

Tan grave como el problema de las correlaciones temporales es para la estimación de la suma de correlación, la solución es simple \textit{Theiler1986}. En la Ec.~(6.1) tenemos que excluir aquellos pares de puntos que están cercanos, no por la geometría del atractor, sino simplemente porque están correlacionados. Usualmente, este es el caso cuando están cercanos en el tiempo. Así, la segunda suma comienza solo después de que haya transcurrido un intervalo de tiempo típico $t_{\min} = n_{\min} \Delta t$:

\begin{equation}
C(m, \epsilon) = \frac{2}{(N - n_{\min})(N - n_{\min} - 1)} \sum_{i=1}^{N} \sum_{j=1, j \neq i + t_{\min}}^{N} \Theta(\epsilon - \| \mathbf{s}_i - \mathbf{s}_j \|) \,.
\end{equation}

Aquí hemos reemplazado los vectores del espacio de estados por los vectores de incrustación por retardos $\mathbf{s}_n$, donde la dimensionalidad $m$ entra como un parámetro. Nótese que $t_{\min}$ no es necesariamente igual al tiempo de correlación agregada, definido, por ejemplo, por el tiempo en que la función de autocorrelación ha decaído a $1/e$. Como buscar pares cercanos requiere estadísticas amplias, las correlaciones entre estos no deben hacer que se perciban como ``similares'' en promedios. Por tanto, debemos ser generosos con la elección de $t_{\min}$: también perdemos algunos pocos ``verdaderos'' vecinos, pero la pérdida de estadísticas es soportable — es solo una fracción de aproximadamente $2 n_{\min} / N$ de todos los pares. Estudiaremos los efectos de no tener en cuenta las correlaciones temporales con más detalle en un ejemplo real. Esto establece un procedimiento estándar para calcular la suma de correlación.

\newpage

\textbf{Ejemplo 6.2 (Suma de correlación de datos de láser NMR).} Consideremos los datos del láser NMR que hemos encontrado en varias ocasiones anteriormente. (Véase el Apéndice B.2 para un resumen). Después de inspeccionar visualmente los datos (en el muestreo estroboscópico) en busca de no estacionariedades evidentes, observar retratos de fase con diferentes retardos y en distintas porciones de los datos disponibles, nos aventuramos a calcular la suma de correlación Ec.(6.3), para este conjunto de datos. Solo interpretaremos el resultado sujeto a verificaciones de consistencia adicionales descritas más adelante en este capítulo. Por ahora debemos elegir un tiempo de retardo para la incrustación, un rango de dimensiones de incrustación interesantes y un tiempo de correlación $t_{\min}$ con el fin de descartar vecinos temporales que podrían afectar negativamente el resultado. Para este conjunto de datos tipo mapa, no hay razón para elegir un tiempo de retardo distinto de 1. Calculemos entonces $C(m, \epsilon)$ en incrustaciones de dimensión de 2 a 7. Si siete resulta insuficiente, podemos repetir el cálculo con valores mayores. Es muy probable que estemos del lado seguro si descartamos todos los pares más cercanos de 500 pasos en el tiempo (véase el Ejemplo 6.6 más adelante). Los tiempos de correlación de los datos son mucho más cortos, pero podemos permitirnos fácilmente la pérdida resultante del 2.5\% de los pares para fines estadísticos. La Fig.6.3 muestra las sumas de correlación $C(m, \epsilon)$ obtenidas con estas elecciones. Como es típico en datos experimentales deterministas de baja dimensión, encontramos algo similar a una ley de potencias para $C(m, \epsilon)$, aunque solo dentro de un pequeño rango de escalas de longitud $\epsilon$, aquí en la región $100 < \epsilon < 1000$. El comportamiento de ley de potencias de $C(m, \epsilon)$ como firma de autosimilitud puede encontrarse mejor graficando la pendiente $D(m, \epsilon)$ de una gráfica doble logarítmica de $C(m, \epsilon)$ contra $\epsilon$, esta última también en escala logarítmica. Esto se ha hecho en la Fig.~6.4. Esta representación muestra mucho más claramente la \textit{meseta} de $D(m, \epsilon) = \partial \ln C(m, \epsilon) / \partial \ln \epsilon$ que corresponde a la ley de potencias deseada para $C$. Podemos encontrar fácilmente que el valor de meseta de $D$ no cambia mucho con la dimensión de incrustación $m$ tan pronto como $m > 2$, aunque aún hay algunas fluctuaciones presentes. Sin embargo, la curva estimada $D(m, \epsilon)$ es lo suficientemente característica como para sugerir que los datos son una muestra tomada de un atractor extraño de dimensión $D < 2$. Evidencia más convincente será presentada una vez que hayamos aplicado reducción de ruido no lineal a los datos; véase el Ejemplo 10.3.

\newpage

\section{Interpretación y dificultades}

Habrás notado que hasta ahora hemos evitado en gran medida hablar sobre la \textit{dimensión de correlación} de un conjunto de datos. La razón es que, desde nuestro punto de vista, la cantidad que puede calcularse con un procedimiento más o menos automatizado es la suma de correlación $C(m, \epsilon)$, mientras que una \textit{dimensión} solo puede asignarse como resultado de una cuidadosa \textit{interpretación} de estas curvas y de cualquier otra información disponible. Mientras que la suma de correlación está bien definida para cualquier conjunto finito de puntos, la dimensión de este conjunto finito es cero y solo tiene un valor no trivial bajo la suposición de una medida subyacente. Así, una estimación de dimensión implica una extrapolación no trivial desde un conjunto de datos finito hacia un atractor que supuestamente existe.En esta sección explicaremos las características principales que podemos esperar encontrar en una suma de correlación típica para una muestra tomada de un atractor extraño. Este conocimiento nos ayudará a encontrar una región de escalas de longitud caracterizadas por autosimilitud — si es que tal región está presente.


\bigskip

Veamos un resultado típico del algoritmo de correlación aplicado a una muestra experimental de un sistema dinámico de baja dimensión. En particular, queremos estudiar las pendientes locales $D(m, \epsilon)$ que se grafican para los datos del láser NMR, Ejemplo 6.2, en la Fig.~6.4. Recuerda que las diferentes curvas representan diferentes dimensiones de incrustación $m$. Se pueden distinguir claramente cuatro regiones diferentes que ahora discutiremos en orden desde escalas grandes hasta escalas pequeñas. En el \textit{régimen macroscópico}, el escalamiento o la autosimilitud se destruyen obviamente por el corte impuesto por el tamaño finito del atractor. Para valores grandes de $\epsilon$ (aquí por encima de $\epsilon > 1000$), la suma de correlación no muestra un comportamiento de ley de potencias ya que las estructuras macroscópicas del atractor determinan su valor, y su exponente local de escalamiento depende tanto de $m$ como de $\epsilon$. En escalas de longitud más pequeñas, puede encontrarse un verdadero \textit{rango de escalamiento}. Para un conjunto autosimilar, el exponente local de escalamiento es constante, y lo mismo se cumple para todas las dimensiones de incrustación mayores que $n_{\min} > D$. Si esta meseta es lo suficientemente convincente, este exponente de escalamiento puede utilizarse como estimación de la dimensión de correlación del conjunto fractal. En la Fig.~6.4, las curvas son relativamente planas para $100 < \epsilon < 3000$ pero el valor de $D$ aumenta con $m$ para $\epsilon > 1000$. Por tanto, podemos hablar de una \textit{meseta} en casi todo el rango $100 < \epsilon < 1000$.

\bigskip

En escalas más pequeñas, usualmente podemos discernir un \textit{régimen de ruido}. Si los datos tienen ruido, entonces por debajo de una escala de longitud de unas pocas veces el nivel de ruido, el método detecta que los puntos de datos no están confinados a la estructura fractal sino que están dispersos por todo el espacio de estados disponible. Así, los exponentes de escalamiento locales aumentan y eventualmente alcanzan el valor verdadero $D(m, \epsilon) = m$, la dimensión de incrustación. El comportamiento es diferente por errores puros de discretización. Entonces, los vectores del espacio de fases reconstruidos están confinados a una cuadrícula regular en el espacio de dimensión $m$ y así parecen tener dimensión cero cuando se observan en escalas de longitud más pequeñas que la resolución de los datos. Por encima de estas escalas de longitud, pueden ocurrir oscilaciones extrañas. Añadir ruido artificial blanco a datos toscamente discretizados antes de aplicar el algoritmo de correlación (Möhler et al. (1989), Theiler (1990a)) ha demostrado de manera más convincente que parte de la información perdida se debe a errores de discretización, y no a fallos del propio algoritmo de correlación. Véase también el Ejemplo 6.4 a continuación.

\bigskip

A escalas aún más pequeñas, enfrentaremos una \textit{falta de vecinos} dentro de las $\epsilon$-vecindades. Los errores estadísticos se convierten en el efecto dominante y las curvas comienzan a fluctuar enormemente. Esta escala es menor que la distancia promedio entre puntos en el espacio de incrustación (que es aproximadamente $N^{-1/m}$ en el régimen de ruido) y se caracteriza por el hecho de que, por ejemplo, menos de 50 pares contribuyen a $C(m, \epsilon)$. Dado que este es un libro de texto, los datos mostrados en la Fig.~6.4 deben considerarse como un ejemplo de libro. Con diferentes datos encontrarás desviaciones del esquema descrito anteriormente. Discutamos primero aquellas desviaciones que pueden esperarse \textit{incluso si} la señal subyacente del conjunto de datos es determinista y de baja dimensión. Lo especialmente confuso acerca de las estimaciones de dimensión (no solo del algoritmo de correlación) es que existen conjuntos de datos que \textit{no} son deterministas, pero que se asemejan aproximadamente a datos de baja dimensión, ya que violan ciertos requisitos de los datos. Afortunadamente, con el cuidado suficiente, uno puede evitar por completo ser engañado por estos problemas, como mostraremos en otro punto. Una señal determinista de baja dimensión, medida con un error relativo mayor a aproximadamente 2--3\%, no mostrará una región de escalamiento significativa de $C(m, \epsilon)$ (ni una meseta de $D(m, \epsilon)$), ya que los artefactos del ruido se superponen con los artefactos de la forma global del atractor. Como veremos en la Sección 6.7, la suma de correlación se ve afectada por el ruido hasta escalas de longitud que son al menos tres veces el nivel del ruido. Incluso en los casos más optimistas, no se puede esperar un escalamiento por encima de aproximadamente una quinta parte de la extensión del atractor, y una sola octava no es suficiente para una región de escalamiento útil.

\newpage

\textbf{Ejemplo 6.3 (Suma de correlación del flujo de Taylor–Couette).} Al observar algún sistema que varía en el tiempo, el nivel mínimo de ruido alcanzable depende de varios factores. Obviamente, un circuito electrónico o un sistema láser es más fácil de controlar que, por ejemplo, el flujo turbulento en un tanque. Así, los datos en este ejemplo, aunque más ruidosos que, por ejemplo, los datos del láser NMR que usamos anteriormente, deben considerarse como un conjunto de datos de calidad excepcionalmente alta. El fenómeno físico observado aquí es el movimiento de un fluido viscoso entre cilindros coaxiales, llamado \textit{flujo de Taylor–Couette}. Los datos fueron proporcionados por Thorsten Buzug y Gerd Pfister; véase Buzug \textit{et al.} (1990a) y Pfister \textit{et al.} (1992), Apéndice B.4. Debido al mayor nivel de ruido y a la aparente mayor dimensión del conjunto de datos, solo puede discernirse una pequeña región de escalamiento en la Fig.6.5. 

\bigskip

\textbf{Ejemplo 6.4 (Suma de correlación de datos de láser infrarrojo lejano).} En el Ejemplo 4.5 realizamos una reducción de ruido no lineal en los datos de láser infrarrojo lejano, Apéndice B.1. La resolución de este conjunto de datos está predominantemente limitada por la discretización de 8 bits, mientras que otros errores de medición parecen ser bastante pequeños. Este es un buen ejemplo para estudiar el efecto de mediciones con grano grueso sobre la suma de correlación, lo cual hacemos en la Fig.~6.6, donde mostramos $D(m, \epsilon)$ para una incrustación de dimensión seis. La discretización hace que el exponente efectivo de escalamiento $D(m, \epsilon)$ tienda a cero en escalas menores que la unidad (diamantes). Obviamente, una cuadrícula hace que el espacio de fases sea discreto y, por tanto, de dimensión cero. Pero observamos artefactos incluso en escalas de longitud mucho mayores, digamos para $\epsilon < 10$. Una solución rápida es desplazar los puntos fuera de la cuadrícula agregando números aleatorios no correlacionados y uniformemente distribuidos en el intervalo $[-0.5, 0.5]$ (cruces). Aun así, los errores son del mismo orden de magnitud, pero tienen propiedades mucho más agradables y conducen a una curva más suave para $D(m, \epsilon)$. Un mejor remedio es pasar los datos por un algoritmo de reducción de ruido no lineal. Esto tiene sentido ya que la suma de correlación original sugiere que los datos podrían al menos describirse efectivamente mediante una dimensión baja. Los cuadrados en la Fig.~6.6 indican la suma de correlación obtenida después de que se aplicó el método de reducción de ruido simple de la Sección 4.5 (véase el Ejemplo 4.5, en particular la Fig.~4.6). El nivel de ruido se reduce visiblemente. Aunque esta curva se ve bastante bien, recuerda que una gráfica de $D(m, \epsilon)$ para una sola dimensión de incrustación $m$ no es suficiente para establecer una región de escalamiento o incluso para determinar una dimensión. También debemos verificar que $d$ se satura al aumentar $m$. 

\bigskip

Incluso datos deterministas de muy buena calidad pueden carecer de una región de escalamiento porque las estructuras macroscópicas impiden el escalamiento hacia escalas pequeñas. De hecho, la teoría no ofrece indicación alguna de cómo debe comportarse asintóticamente el escalamiento $C(m, \epsilon) \propto \epsilon^D$, cf. Ec.~(6.2), por lo que debe aproximarse a escalas de longitud finitas. Un ejemplo destacado es el de datos biperiódicos que forman un \textit{torus} en el espacio de fases. Un torus es una pieza de tubo cerrado que forma un anillo. No importa qué tan grande sea el anillo, la representación adecuada con dimensión dos (de la superficie del tubo) no puede encontrarse por encima de aproximadamente la mitad del diámetro del tubo. Un tubo delgado incluso puede parecer unidimensional en escalas intermedias de longitud. (Véase también el Ejercicio 6.5.) De la discusión anterior debe quedar inmediatamente claro que una región de escalamiento siempre debe establecerse por inspección. Incluso los conjuntos de datos del mismo experimento pueden variar tanto en su estructura geométrica que no recomendamos el uso de ningún procedimiento automatizado para determinar el exponente de escalamiento $D$ como estimación de la dimensión de correlación. Hasta ahora solo hemos considerado el caso en que los datos son efectivamente una muestra de un atractor de baja dimensión, y nos hemos preguntado bajo qué circunstancias se puede esperar estimar su dimensión a partir de la suma de correlación. Por supuesto, normalmente no sabemos de antemano si el caos determinista está presente en el sistema. Estrictamente hablando, estamos usando el algoritmo —de hecho incluso la incrustación retardada— fuera de su hábitat natural. La esperanza es, sin embargo, encontrar firmas que reconozcamos como típicas de sistemas deterministas y así justificar el enfoque \textit{a posteriori}, un concepto peligroso. Desafortunadamente, existen conjuntos de datos que \textit{no} son deterministas pero que, no obstante, llevan a algo que se parece a una \textit{meseta} en la curva $D(m, \epsilon)$ si no se tiene suficiente cuidado. La siguiente sección te mostrará cómo tener más precaución. Personas poco críticas encuentran líneas rectas en prácticamente cada gráfica doble logarítmica. Algunos autores ni siquiera buscan escalamiento, sino que \textit{definen} una cierta gama de escalas de longitud como ``región de escalamiento''. Con este método siempre se obtiene un número que puede llamarse \textit{dimensión}, pero ese número es bastante insignificante. En particular, un número pequeño no implica determinismo ni caos. Por supuesto, no debes seguir esta mala práctica, pero iríamos un paso más allá y te aconsejaríamos no creer \textit{ninguna} estimación de dimensión publicada si no viene acompañada de una gráfica de escalamiento convincente.

\newpage

\section{Correlaciones temporales, no estacionariedad y diagramas espacio–tiempo}

Poco después de su publicación en 1983, la dimensión de correlación se volvió bastante popular y, desde entonces, muchos trabajos han reportado dimensiones finitas o incluso muy bajas para una gran variedad de sistemas. [Ejemplos usando datos experimentales incluyen a Nicolis \& Nicolis (1984), Fraedrich (1986), Essex \textit{et al.} (1987), Kurths \& Herzel (1987), Tsonis \& Elsner (1988), y Babloyantz (1989).] Generalmente, estas afirmaciones se hacían cuando los autores lograban ajustar una línea recta a una porción de la gráfica log-log de la suma de correlación $C(m, \epsilon)$. Algunos autores no notaban que las curvas que estaban ajustando con líneas rectas no eran del todo rectas, sino más bien con forma de S. Solo más tarde, varios autores encontraron ejemplos de conjuntos de datos no deterministas que producían sumas de correlación con tramos casi rectos, lo cual permitiría interpretar que $C(m, \epsilon)$ sigue una ley de potencias; véanse Theiler (1986) y (1991), Osborne \textit{et al.} (1986), y Osborne \& Provenzale (1989). Esto parecía contradecir el hecho de que los datos estocásticos deberían tener dimensión infinita. Un examen más detallado muestra que todos estos ejemplos, y de hecho muchas de las estimaciones de dimensión publicadas para datos de series temporales, padecen del mismo problema: \textit{correlaciones temporales}. La suma de correlación es un caso particular de un concepto abstracto, la \textit{integral de correlación}, que es funcional sobre la medida invariante subyacente. La estimación solo es imparcial si la muestra finita que usamos es una muestra no correlacionada de puntos en el espacio de dimensión $m$ de acuerdo con la medida invariante en dicho espacio. Los datos de series temporales son dinámicamente correlacionados y claramente violan este requisito. Sin embargo, si la violación es muy débil, como en datos de una señal caótica altamente mixta, su efecto será casi invisible. Si las correlaciones son fuertes, pueden arruinar todo. Como señalamos anteriormente, en muchos procesos como mapas con cambios abruptos, los pares de puntos que fueron medidos dentro de un intervalo de tiempo corto también son cercanos en el espacio de fases, lo cual introduce un sesgo cuando se estima la suma de correlación a partir de su definición, Ec.~(6.1). Una vez que conocemos el intervalo de tiempo típico dentro del cual los datos están correlacionados, el problema puede resolverse fácilmente rechazando los pares que están demasiado cercanos en el tiempo, como se describió en la Sección~6.3 anteriormente; en particular, compara la Ec.~(6.1) con la Ec.~(6.3). Cuando encuentres un artículo donde la dimensión fractal de datos de flujo caótico es cercana pero menor que dos (lo cual es una contradicción), lo más probable es que no se hayan excluido de la suma de correlación los puntos cercanos en el tiempo.

\bigskip

\textbf{Ejemplo 6.5 (Cálculo incorrecto de la suma de correlación del flujo de Taylor–Couette).}  Las correlaciones temporales están presentes en casi todos los conjuntos de datos. Aunque su efecto sobre la suma de correlación es más severo para datos cortos y muestreados densamente, nunca debemos dejar de tomar la precaución de descartar vecinos temporales. Como advertencia, repitamos el cálculo de la suma de correlación para los datos del flujo de Taylor–Couette (Ejemplo 6.3; véase también el Apéndice B.4). Usamos los primeros 10,000 puntos de datos y la fórmula original, Ec.~(6.1), que solo excluye los pares idénticos de la suma (Fig.~6.7). La pregunta que queremos abordar en esta sección es cómo detectar estas correlaciones temporales y cómo estimar un valor seguro para el tiempo de correlación $t_{\min}$. Este problema fue resuelto por Provenzale \textit{et al.} (1992) mediante la introducción del \textit{diagrama espacio–tiempo} (space time separation plot). La idea es que en presencia de correlaciones temporales, la probabilidad de que un par de puntos tenga una distancia menor que $\epsilon$ no depende únicamente de $\epsilon$, sino también del tiempo transcurrido entre las dos mediciones.

\bigskip

Esta dependencia puede detectarse graficando el número de pares como función de dos variables: la separación temporal $\Delta t$ y la distancia espacial $\epsilon$. Puedes realizar un diagrama de dispersión, o trazar una superficie o líneas de contorno con tu programa de visualización favorito. Técnicamente, para cada separación temporal $\Delta t$ se crea un histograma acumulado de las distancias espaciales $\epsilon$. Recuerda normalizar adecuadamente. Puedes utilizar el programa \texttt{stp} del paquete \texttt{TISEAN}, véase la Sección A.3.

\newpage

\textbf{Ejemplo 6.6 (Diagrama espacio–tiempo para los datos del láser NMR).} En la Fig.~6.8 mostramos un diagrama espacio–tiempo para los datos del láser NMR (Apéndice B.2). Puedes leer el gráfico de la siguiente manera: las líneas de contorno indican la distancia que debes recorrer para encontrar una fracción dada de pares, dependiendo de su separación temporal $\Delta t$. Solo para valores de $\Delta t$ donde las líneas de contorno son planas (al menos en promedio, pueden oscilar), la correlación temporal no causa artefactos en la suma de correlación. Para los datos del láser NMR, este es el caso aproximadamente para $\Delta t > 100$. Por lo tanto, realmente estuvimos en lo correcto al elegir $t_{\min} = 500$ al calcular la suma de correlación en el Ejemplo~6.2. 

\bigskip

\textbf{Ejemplo 6.7 (Diagrama espacio–tiempo para el flujo de Taylor–Couette).} El diagrama espacio–tiempo para los datos del flujo de Taylor–Couette (Apéndice B.4) parece un poco más complicado. La frecuencia dominante en el sistema también se refleja en las curvas de nivel mostradas en la Fig.~6.9. Mientras la oscilación no sea problemática, siempre que el período de observación sea mucho mayor que la longitud de un ciclo, debemos preocuparnos por los primeros 20–40 pasos de tiempo, donde los niveles aumentan de forma consistente. Si no descartamos al menos los pares que están separados por menos de 20 pasos de tiempo, entonces obtenemos la suma de correlación espuria mostrada en la Fig.~6.7. Para estar seguros, usualmente elegimos $t_{\min} = 500$ o más — la pérdida en estadística es marginal. 

\bigskip

Todos los ejemplos de conjuntos de datos que exhiben un escalamiento espurio de $C(m, \epsilon)$ (véanse las referencias anteriores) no habrían sido confundidos con caos de baja dimensión si las correlaciones temporales se hubieran manejado correctamente, excluyendo los pares cercanos en el tiempo de la suma de correlación. El diagrama espacio–tiempo nos indica cuántos pares deben excluirse. Las correlaciones temporales están presentes mientras las curvas de nivel no se saturen. ¿Pero qué pasa si las curvas nunca se saturan? Esto es precisamente lo que ocurre con datos que tienen gran potencia en las frecuencias bajas, como el ruido $1/f$ o el movimiento browniano. En este caso, \textit{todos} los puntos del conjunto de datos están correlacionados temporalmente y no hay forma de determinar un \textit{atractor} a partir de la muestra. Una situación similar se presenta si el conjunto de datos es demasiado corto: entonces casi no quedan pares tras eliminar los correlacionados temporalmente. Si abordamos el problema desde otro punto de vista, los tiempos de correlación del orden de la longitud de la muestra (curvas de nivel no saturantes) significan que los datos no muestrean suficientemente el fenómeno observado. Si quieres saber qué programa de televisión se transmite todos los días, no basta con ver la televisión un solo día. Esta deficiencia en los datos también puede verse como una forma de no estacionariedad de la medición. Es muy común en ciencia que los métodos solo sean aplicables si se cumplen ciertos requisitos. La dimensión de correlación no es la excepción. Si se aplica correctamente, nunca da \textit{resultados incorrectos} (como dimensiones finitas) para ruido $1/f$, por ejemplo. Simplemente no da \textit{ningún resultado}. Cuando usamos el diagrama espacio–tiempo, obtenemos una advertencia cuando no se puede aplicar el algoritmo.

\bigskip

No obstante, podrías pensar que una dimensión finita espuria de un conjunto completamente correlacionado podría ser característica del proceso subyacente. De hecho, esto es correcto para ciertos \textit{procesos no recurrentes}, como la difusión libre: el conjunto en el espacio formado por una trayectoria browniana típica en dos o más dimensiones tiene una dimensión de Hausdorff de 2 [Mandelbrot (1985)]. Este valor es, de hecho, el que se obtiene en un cálculo numérico de la dimensión de correlación sobre escalas de longitud adecuadas (dependiendo de la duración de la serie temporal), si se usa $n_{\min} = 1$ en la Ec.~(6.3). La trayectoria browniana es un fractal bastante extraño, ya que una vez incrustada en un espacio $m$-dimensional, no se restringe a una variedad de menor dimensión (en particular, no a una bidimensional, como podría esperarse por su dimensionalidad). Las generalizaciones de los procesos de difusión (\textit{difusión anómala}) producen geometrías con propiedades similares pero diferentes valores de dimensión. La falta de recurrencia se manifiesta también en el espectro de potencia, mediante una potencia divergente en frecuencias bajas. Para ciertos procesos de ruido coloreado con un espectro tipo $f^{-\alpha}$ con $1 < \alpha < 3$, la dimensión (de correlación) tomará el valor finito $D_2 = 2 / (\alpha - 1)$ [Osborne \& Provenzale (1989), Theiler (1991)], si se usa para caracterizar una sola trayectoria en lugar de la distribución probabilística suave subyacente. Una caracterización más adecuada de tales procesos es el análisis en términos de \textit{exponentes de Hurst}, como se discutirá en la Sección 6.8.

\newpage

\section{Consideraciones prácticas}

Obviamente, tendremos que usar una computadora para obtener la suma de correlación en la práctica. (Véase, sin embargo, el Ejercicio 6.2). Como siempre, hay formas inteligentes y no tan inteligentes de implementar la Ec.~(6.3). Solo cuando se tiene acceso a una computadora realmente rápida se puede permitir usar la solución no tan inteligente e implementar la fórmula como dos bucles anidados sin pensarlo mucho. Debe quedar claro a estas alturas que necesitas una cantidad considerable de puntos para estimar $C(m, \epsilon)$ en un rango suficientemente amplio de escalas de longitud. Unos pocos cientos no son suficientes para obtener un resultado estadísticamente significativo en las escalas de longitud pequeñas. Para $N$ puntos de datos, la doble suma en la definición de $C(m, \epsilon)$ contiene aproximadamente $N^2/2$ términos, lo cual hace que sea muy costoso computacionalmente calcular $C(m, \epsilon)$ para todos los conjuntos de datos, salvo aquellos que son demasiado pequeños de todos modos. Dado que el número de pares cercanos determina la menor escala de longitud accesible, no podemos darnos el lujo de desperdiciar parte de los valiosos datos limitando el rango de una o ambas sumas (aunque eso es lo que algunas personas hacen). Es más adecuado observar que \textit{solo} en las escalas de longitud más pequeñas necesitamos la máxima cantidad posible de estadísticas. Para distancias mayores, fácilmente podemos encontrar suficientes pares. Esta observación ha llevado a varios autores a la idea de usar un método de búsqueda de vecinos cercanos rápido para encontrar \textit{todos} los pares cercanos dentro del conjunto de datos, y luego aplicar el algoritmo de forma más ingenua para las escalas de longitud mayores. Este no es el lugar adecuado para discutir las distintas posibilidades para encontrar vecinos cercanos eficientemente, un problema que ocurre con frecuencia en el análisis de series temporales no lineales. Daremos algunas referencias en la literatura y describiremos el método utilizado en el paquete TISEAN en la Sección A.1.1.

\bigskip

Las Figs.~6.3 y 6.4 se obtuvieron usando este algoritmo, de hecho usando la rutina \texttt{d2} del paquete TISEAN; véase la Sección A.6. Computamos la suma de correlación $C(m, \epsilon)$ para dos valores de $\epsilon$ en cada octava, es decir, valores adyacentes que difieren por un factor de $\sqrt{2}$. Aunque debe quedar claro de inmediato cómo se obtiene un diagrama como el de la Fig.~6.3, también debemos calcular derivadas numéricas para la Fig.~6.4. Suprimimos las fluctuaciones estadísticas evitando diferencias finitas entre valores adyacentes de $\epsilon$ y estimando en su lugar:

\begin{equation}
D(m, \epsilon) \approx \frac{\ln C(m, 2\epsilon) - \ln C(m, \epsilon/2)}{\ln 2\epsilon - \ln \epsilon/2}
= \frac{1}{2} \log_2 \frac{C(m, 2\epsilon)}{C(m, \epsilon/2)}.
\end{equation}

Cabe señalar que el algoritmo se comporta de manera un tanto problemática cuando el conjunto de datos contiene muchos pares \textit{idénticos}, situación que solo puede ocurrir si los datos están muy groseramente discretizados. El programa no verifica esto, ya que conjuntos de datos de este tipo no son adecuados para cálculos de dimensión de todos modos. Por supuesto, formalmente, puntos en una red discreta en el espacio forman nuevamente un conjunto de dimensión cero, incluso si fueran infinitamente muchos. De ahí que tengamos nuevamente el problema de extrapolar correctamente la dimensión desde escalas mayores. Para lograrlo, debemos minimizar los efectos de los artefactos sobre la suma de correlación en escalas grandes. Podemos, por ejemplo, procesar tales datos mediante un algoritmo de reducción de ruido no lineal antes de estimar la suma de correlación. Si eso no es posible, podemos añadir ruido blanco de una magnitud al menos comparable con la resolución de los datos.

\newpage

\section{Una aplicación útil: determinación del nivel de ruido mediante la integral de correlación}

El algoritmo de correlación es una herramienta para analizar las propiedades de escalamiento de conjuntos de puntos en el espacio de fases. Los atractores de sistemas deterministas muestran un escalamiento tipo ley de potencias de $C(m, \epsilon)$. Es característico de estos sistemas que el exponente en la ley de potencias sea invariante bajo transformaciones de coordenadas y, en particular, no dependa de la dimensión de incrustación $m$ una vez que $m$ es suficientemente grande para garantizar una reconstrucción adecuada. Para los sistemas estocásticos la situación es diferente. Los vectores retardados formados a partir de una señal aleatoria no se restringen a una variedad de baja dimensión, sino que llenan todas las direcciones disponibles en el espacio de fases. Aparte de efectos de borde, el exponente de escalamiento es siempre igual a $m$.\footnote{Si las correlaciones temporales decaen suficientemente rápido, véase la última sección.} Consideremos ahora el caso típico de una señal determinista que es medida con resolución finita debido a errores de medición y digitalización de la salida. Muy por encima del nivel de ruido, los errores pueden ser ignorados y el exponente de escalamiento correcto $D(m, \epsilon) \approx d$ puede observarse dentro de cierto rango de escalamiento. A escalas pequeñas comparables con la amplitud de los errores aleatorios, este escalamiento se rompe y el comportamiento típico para datos estocásticos, $D(m, \epsilon) = m$, domina en su lugar. Esta diferencia cualitativa llevó a la distinción entre una \textit{región de escalamiento} y una \textit{región de ruido} en la Sección 6.4 anterior. (Véase también el Ejemplo 6.2 y las gráficas correspondientes.)

\bigskip

En esta sección estudiaremos las consecuencias para las curvas $D(m, \epsilon)$ con más detalle, lo cual nos permitirá medir la intensidad del ruido (el error promedio de la medida) presente en un conjunto de datos determinista. Tal estimación será muy útil en la interpretación de resultados de métodos no lineales. Si somos capaces de medir el éxito de una reducción de ruido no lineal, podremos entender mejor los límites de predictibilidad del conjunto de datos y contrastarlos con lo que aprendimos acerca del ruido en la Sección 5.3 al estudiar las trayectorias. Comparemos las sumas de correlación $C(m, \epsilon)$ y $C(m+1, \epsilon)$ obtenidas con incrustaciones en $m$ y en $m+1$ dimensiones respectivamente. Supongamos que $m$ es suficientemente grande para una reconstrucción fiel. Si los datos fueran completamente deterministas y sin ruido, las dos curvas diferirían únicamente por un factor constante a través del rango de escalamiento. Esta constante está determinada por un promedio del ritmo local de expansión, relacionado con la entropía de correlación (Sección 11.4). Los puntos en el espacio de dimensión $m$ están también cerca en el espacio de dimensión $m+1$ — de hecho podríamos predecir la $(m+1)$-ésima coordenada a partir de los vectores de incrustación en dimensión $m$. En la práctica, todas las coordenadas, en particular la $(m+1)$-ésima, están contaminadas con ruido. Por tanto, incluso por debajo del nivel de ruido, la suma en dimensión $m+1$ siente la presencia del ruido en el conjunto de datos. El siguiente resultado intuitivo para la dependencia en $\epsilon$ de la suma de correlación puede derivarse con un argumento heurístico tomado de Schreiber (1993a):

\begin{equation}
C(m+1, \epsilon) \propto C(m, \epsilon) \, C_\text{noise}(\epsilon).
\end{equation}

Aquí $C_\text{noise}(\epsilon)$ es la suma de correlación de la distribución unidimensional del ruido, suponiendo que el tiempo de autocorrelación del ruido es más corto que el intervalo de incrustación. Si el ruido $\eta$ proviene de una distribución continua $\mu(\eta)$ entonces:

\begin{equation}
C_\text{noise}(\epsilon) = \int_{-\infty}^{\infty} \mathrm{d}\eta\, \mu(\eta) \int_{\eta - \epsilon}^{\eta + \epsilon} \mathrm{d}\eta'\, \mu(\eta'),
\end{equation}

lo cual puede evaluarse cuando se conoce $\mu(\eta)$. Para ruido gaussiano con media cero y amplitud $\sigma$, tenemos $\mu(\eta) \propto \exp(-\eta^2/2\sigma^2)$ y:

\begin{equation}
D_\text{Gauss}(\epsilon) = \frac{\epsilon \exp(-\epsilon^2/4\sigma^2)}{\sigma \sqrt{\pi} \, \text{erf}(\epsilon/2\sigma)}.
\end{equation}

Para ruido uniforme en el intervalo $-\sigma/2 < \eta < \sigma/2$, el resultado es:

\begin{equation}
D_\text{uniform}(\epsilon) = \frac{2 - 2\epsilon/\sigma}{2 - \epsilon/\sigma} \, \Theta(\sigma - \epsilon).
\end{equation}

Sin embargo, cualquier filtrado o procedimiento de incrustación que no sea la incrustación por retardo simple hará que la mayoría de los ruidos se comporten más como una distribución gaussiana, según el teorema del límite central. De hecho, encontramos usualmente que cuando $m$ coordenadas retardadas producen una incrustación adecuada, entonces la diferencia entre las curvas $D(m+1, \epsilon)$, $D(m+2, \epsilon)$, \dots{} se describe bien con la Ec.~(6.7), si elegimos $\sigma$ apropiadamente. Finalmente, señalamos que el error por discretización no equivale a ruido blanco uniforme. Induce correlaciones geométricas espurias y hace que la señal parezca tener dimensión cero en escalas pequeñas.

\newpage

\textbf{Ejemplo 6.8 (Nivel de ruido de los datos del flujo de Taylor–Couette).} Como ejemplo, grafiquemos los datos de la Fig.~6.5, es decir, las pendientes locales $D(m, \epsilon)$ de la suma de correlación para los datos del flujo de Taylor–Couette (Apéndice B.4), de una manera diferente. En lugar de $D(m, \epsilon)$ para $m = 2, \dots, 10$, mostramos en la Fig.~6.10 la diferencia $D(m+1, \epsilon) - D(m, \epsilon)$ para $m = 5, \ldots, 10$ (para valores menores de $m$ no tenemos una buena incrustación). Además, nos concentramos en la región de transición entre el régimen de ruido y el rango de escalamiento usando una escala lineal en $\epsilon$. Luego probamos diferentes estimaciones del nivel de ruido $\sigma$ y comparamos los datos con las curvas correspondientes $D_{\text{Gauss}}(\epsilon)$ dadas por la Ec.(6.7). 

\bigskip

En el ejemplo anterior utilizamos un ajuste visual para extraer una estimación del nivel de ruido $\sigma$. Por supuesto, podemos reemplazar este ajuste visual por un esquema de mínimos cuadrados. Consideramos que el método visual es suficientemente bueno, ya que con datos reales debemos esperar desviaciones considerables respecto a la forma esperada debido a la presencia de ruido no gaussiano. Sin embargo, incluso cuando hay tales desviaciones, a menudo se puede obtener información útil sobre el ruido presente en el conjunto de datos. En lo anterior, estudiamos el efecto del ruido en la suma de correlación de forma cuantitativa, con el resultado de que la desviación de la función de escalamiento observada $D(m, \epsilon)$ con respecto a la verdadera está dada por la fórmula $D(m+1, \epsilon) = D(m, \epsilon) + D_{\text{noise}}(\epsilon),$ o bien,
\begin{equation}
D(m, \epsilon) = D(m_0, \epsilon) + (m - m_0) D_{\text{noise}}(\epsilon),
\end{equation}
siempre que $m_0$ sea suficiente para una incrustación adecuada de los datos. Para la mayoría de los tipos de ruido, la forma de $D_{\text{noise}}(\epsilon)$ es similar a la dada por la Ec.~(6.7). Otras referencias sobre el efecto del ruido en la integral de correlación son Diks (1996) y Oltmans \& Verheijen (1997).

\bigskip

Una consecuencia importante de este resultado analítico sobre la forma de la integral de correlación es que una pequeña cantidad de ruido ya oculta el comportamiento de escalamiento: incluso con $\epsilon = 3\sigma$, la dimensión efectiva aumenta visiblemente con $m$ más allá de $m_0$, concretamente, en una cantidad de 0.2 por cada dimensión adicional. Esto significa que, dado que el mejor caso de escalamiento solo puede esperarse hasta, como máximo, una cuarta parte de la extensión del atractor (véanse los ejemplos en este libro) y hasta tres veces el nivel de ruido, un conjunto de datos con 2\% de ruido puede ofrecer a lo sumo una región de escalamiento de solo dos octavas. Nótese que esta es una visión muy optimista: usualmente tenemos que utilizar al menos $D(m+2, \epsilon)$ para verificar que $D(m, \epsilon)$ se satura con $m$ dentro del rango de escalamiento. Ahora podrías estar confundido por el último ejemplo, donde se encontró un nivel de ruido del 5\%, ¡pero se estimó una dimensión finita! Bueno, si vuelves a la Fig.~6.5 verás que la región de escalamiento es realmente muy pequeña, tal como se anticipó por los resultados de esta sección.

\newpage

\section{Señales multi-escala o auto-similares}

En esta sección queremos discutir algunos fenómenos de escalamiento y fractalidad donde la señal misma forma un objeto geométrico complicado, mientras que el concepto de un atractor en el espacio de fases no puede aplicarse exitosamente. Los sistemas dinámicos gobernados por ecuaciones diferenciales ordinarias, tales como la Ec. (3.2), producen señales continuas y diferenciables. Varios procesos estocásticos y deterministas pueden crear señales que se comportan de forma mucho menos regular: aunque continuas, no son diferenciables. Tales señales son ya sea fractales o son invariantes a escala de alguna otra manera y requieren un tratamiento especial, en particular porque a menudo no son recurrentes y por lo tanto deben considerarse efectivamente no estacionarias: muchos de los procesos para los que el análisis siguiente es adecuado no crean una medida invariante cuando se consideran en espacios de incrustación más altos (ver la Sección 6.5). Una señal escalar puede ser \textit{fractal} en el sentido de que su gráfica, como función del tiempo, tiene una dimensión de Hausdorff no trivial. Esto significa que la longitud de la gráfica vista a una resolución finita aumenta indefinidamente a medida que se incrementa la resolución en el eje del tiempo, ya que se resuelven más y más fluctuaciones. Su dimensión puede entonces estar entre uno y dos. Desde el punto de vista matemático, este comportamiento es bastante exótico y no es trivial crear una gráfica de este tipo. Las \textit{funciones de Weierstrass} son de este tipo. Debemos enfatizar que tales señales usualmente no pueden resultar de sistemas caóticos bidimensionales excepto en el caso de difusión mejorada determinista en potenciales periódicos, ver Zaslavsky (1993) y Klafter \& Zumofen (1993), donde, sin embargo, el espacio de fases no está acotado.

\bigskip

\textbf{Ejemplo 6.9 (Turbulencia hidrodinámica)} La turbulencia plenamente desarrollada se caracteriza por la invariancia de escala (con un límite superior e inferior). La energía se introduce en el sistema en las escalas de longitud más grandes, donde se crean grandes remolinos. La disipación (y por tanto la transformación de energía cinética en calor) ocurre en las escalas de longitud más pequeñas (escala viscosa). Las escalas intermedias se denominan \textit{rango inercial}. Una jerarquía de remolinos en todas las escalas representa, en el espacio real, la cascada de energía que conecta las escalas grandes con las pequeñas. Todo esto da lugar a leyes de escalamiento para todos los observables, como las correlaciones del campo de velocidad o el campo de disipación de energía. Una serie temporal experimental refleja esta estructura complicada.  Imaginemos la medición de la velocidad local en una posición fija dentro del fluido. Debido a la mezcla constante del sistema, la sonda registrará todas las escalas. Por lo tanto, \textit{la señal misma} forma un fractal, pero no existe un atractor de baja dimensión.

\newpage

\subsection{\textit{Leyes de escalamiento}}

La turbulencia es un proceso determinista de dimensión infinita descrito por una ecuación diferencial parcial. Los procesos estocásticos pueden producir señales de tipo similar. Recordemos las propiedades de la difusión estándar (proceso de \textit{Wiener}) con el fin de introducir un nuevo concepto relacionado con la autosimilitud, la \textit{autoafinidad}. En el proceso de Wiener, una partícula en la posición \(x(t)\) tiene una probabilidad $P(\Delta x, \Delta t) = \frac{\exp(-\Delta x^2 / 2D\Delta t)}{\sqrt{2\pi D \Delta t}}$ de moverse a la posición \(x + \Delta x\) en el intervalo de tiempo \(\Delta t\). La varianza de los incrementos escala con el intervalo de tiempo como $\langle \Delta x^2 \rangle \propto \Delta t^{1/2}.$ Supongamos que hemos creado una gráfica \(x(t)\) con precisión infinita. Ahora bien, si la observamos muy de cerca, es decir, con buena resolución tanto en tiempo como en espacio, observaremos fluctuaciones divergentes, ya que para las “pendientes” se encuentra fácilmente: $\langle \Delta x^2 \rangle / \Delta t \propto \Delta t^{-1/2}.$ Por otro lado, si observamos desde muy lejos, reduciendo así las escalas en tiempo y espacio por el mismo factor, vemos una curva casi suave (véase la Figura 6.11).

\bigskip

En el límite de alta resolución, la gráfica llena el plano y tiene dimensión dos, mientras que en el límite opuesto, degenera hacia una línea horizontal y por tanto tiene dimensión uno. Así, en cada uno de los dos límites encontramos autosimilitud. Finalmente, podemos encontrar una vista independiente de escala si escalamos espacio y tiempo con factores de escala diferentes, o más precisamente, si reescalamos el espacio con factores que son la raíz cuadrada de los factores de escala en el tiempo. Estos objetos se denominan \textit{autoafines} y son auto-similares sólo en ciertos límites. Pasemos ahora a leyes de difusión anómalas, es decir, leyes donde la raíz cuadrada de la relación entre incrementos y los intervalos de tiempo se sustituye por leyes de potencias no triviales:

\begin{equation}
\langle \Delta x^2 \rangle \propto \Delta t^{2H},
\end{equation}

con un \textit{exponente de Hurst} \(H\) distinto de \(1/2\). El movimiento browniano (o proceso de Wiener) tiene la propiedad de que la distribución de los incrementos \(\Delta x\) posee una varianza finita para cada \(\Delta t\) fija, y los incrementos no están correlacionados en pasos de tiempo sucesivos. Se puede mostrar fácilmente que una generalización de la Ecuación (6.10) hacia \(H \ne 1/2\) requiere relajar al menos una de estas dos propiedades. Si insistimos en incrementos no correlacionados y por tanto en la propiedad de Markov, debemos admitir incrementos con varianza no acotada, aunque aún requerimos que cada incremento individual tenga varianza finita con probabilidad uno. Esto genera una distribución de probabilidad cuyo comportamiento para grandes \(\Delta x\) es $p(\Delta x) \propto \Delta x^{-1/H}$ con \(1/2 \leq H \leq 1\), es decir, una distribución de probabilidad con colas pesadas. Este tipo de proceso, que se llama \textit{vuelo de Lévy}, crea una difusión no estándar.

\bigskip

Un proceso estocástico más natural puede encontrarse si se permiten correlaciones entre los incrementos con una memoria infinita. En este caso, la varianza de los incrementos puede permanecer finita. A este tipo de proceso no markoviano se le denomina \textit{movimiento browniano fraccionario} y genera exponentes \(0 < H < 1\). Podemos encontrar una explicación intuitiva de lo que sucede: los efectos de memoria inducen correlaciones entre incrementos sucesivos. Si están positivamente correlacionados, la difusión se ve aumentada (\textit{persistencia}); si están anti-correlacionados, la difusión se suprime (\textit{anti-persistencia}). El escalamiento a nivel microscópico se traduce directamente al nivel macroscópico. Por lo tanto, para incrementos \(x_t - x_{t_0}\) observados en intervalos de tiempo macroscópicos también se obtiene:
\begin{equation}
\langle (x_t - x_{t_0})^2 \rangle \propto (t - t_0)^{2H},
\end{equation}
donde \(0 < H < 1\) en el caso del movimiento browniano fraccionario y \(1/2 < H < 1\) para los vuelos de Lévy. La difusión ordinaria cumple con \(H = 1/2\), y el movimiento balístico (es decir, \(x(t) \approx at\) con \(a \ne 0\)) tiene \(H = 1\).

\bigskip

\textbf{Ejemplo 6.10 (Fluctuaciones de un humano erguido completamente quieto).} Toda persona en posición erguida y normal realiza involuntariamente algunas oscilaciones de pequeña amplitud, lo cual induce un movimiento complicado del centro de presión bajo las plantas de los pies en el plano \(x\)-\(y\) (se mueve desde los dedos hacia el talón, de izquierda a derecha). La posición de este centro de presión como función del tiempo puede ser medida mediante un dispositivo llamado \textit{plataforma de fuerza}. Resulta que el análisis más adecuado de esta señal es el cálculo de las fluctuaciones mediante la Ecuación (6.11), la cual produce una ley de potencias bien definida y respalda la idea de que las oscilaciones son de naturaleza aleatoria (Fig. 6.12). Los intentos de encontrar una dinámica determinista de baja dimensión han fallado, y no se ha encontrado un desacuerdo significativo con la hipótesis nula de que el proceso es gaussiano y estocástico lineal. Una observación muy interesante es el cruce desde una difusión anómala acelerada con un exponente cercano a \(3/4\) hacia una difusión con un exponente cercano o menor que \(1/2\) en una escala de tiempo de aproximadamente 2 segundos. Es plausible que esto sea una firma del mecanismo de control postural, el cual trata de suprimir grandes fluctuaciones y necesita algo de tiempo para generar una respuesta adecuada. Naturalmente, esto se traduce en una anti-correlación con la posición real (compárese esto con el movimiento browniano fraccionario). Obtuvimos las series de tiempo de Steve Boker \cite[Apéndice B.9]{Bertenthal1996}. Véase también \cite{CollinsDeLuca1995} para un estudio de este fenómeno.

\newpage

La autosimilitud de la señal refleja correlaciones temporales de largo alcance en la serie temporal. Al estudiar el cúmulo en el espacio de incrustación con retardo formado por estos datos, se observa que, para toda serie temporal finita, también existe una estructura estadísticamente autosimilar. Ya mencionamos el resultado obtenido por Mandelbrot, que indica que un cúmulo de difusión tiene dimensión dos en cualquier dimensión de incrustación igual o mayor que dos. La dimensión que refleja las correlaciones temporales aquí está, como era de esperarse, relacionada con el exponente de Hurst por \( D_H = \frac{1}{H} \).

\bigskip

Sin embargo, al estimar numéricamente esta dimensión, es crucial \textit{no descartar puntos cercanos en el tiempo}, ya que, estrictamente hablando, todo el cúmulo está correlacionado temporalmente. En este sentido, no es una contradicción a nuestro razonamiento habitual que un proceso estocástico estacionario cree un conjunto cuya dimensión sea idéntica a la dimensión de incrustación. El aspecto intrincado de estos objetos fractales es que no pueden ser incrustados en una variedad de menor dimensión del espacio en el que se están considerando. Mientras que el atractor fractal de un sistema dinámico típicamente puede (al menos localmente) ser incrustado en una variedad cuya dimensión es el siguiente entero mayor que la dimensión fractal, esto no es cierto para los procesos discutidos aquí. Dado que el exponente \( H \) refleja correlaciones temporales, podrías preguntarte si puede encontrarse en el espectro de potencias. En efecto, el espectro de potencias de dichas señales muestra una caída de ley de potencias hacia frecuencias altas, dada por \( \beta = 1 + 2H \). Por lo tanto, el espectro de potencias de una trayectoria Browniana se comporta como \( 1/k^2 \), donde \( k \) es el número de onda.

\newpage

\subsection{\textit{Análisis de fluctuación sin tendencia}}

No suele ser recomendable estimar el exponente de Hurst a partir de los datos medidos por la Ec.~(6.11). Los datos reales obtenidos de procesos con exponentes Hurst no triviales, como los que resultan en un exponente trivial, es decir, \(H = 1\), reflejan el aumento “balístico” de la cantidad debido a la tendencia o, más a menudo, no muestran una escala de tamaño. Se ha demostrado que en este caso en particular, un ajuste aditivo de la tendencia es razonable. El método fue desarrollado por Peng et al. (1994) para el análisis de ADN y también se aplica a muchos fenómenos dinámicos, como fenómenos fisiológicos, clima, y economía. \footnote{En noviembre de 2002, cuando se realizó la revisión final de este libro, la base de datos de ISI contenía más de 150 artículos que hacen uso de DFA en series temporales.} 

\bigskip

Antes de pensar en eliminar la tendencia, a menudo se necesita pasar a un sistema diferente. Considere, por ejemplo, las mediciones de velocidad de un río turbulento. Las velocidades registradas son definitivamente finitas (una velocidad de más de 30 m/s en la superficie de la Tierra es realmente excepcional), por lo que la escala tal como la describe la Ec.~(6.11) estaría restringida a intervalos de tiempo muy cortos, más allá de los cuales esta cantidad se saturaría hasta alcanzar la varianza de los datos. Para extender el escalado, debe realizarse el análisis de tales mediciones con una precisión temporal finita, como los incrementos de un segundo, tal como en procesos de difusión. Esto se logra mediante la transformación
\begin{equation}
y(t) = \int_{0}^{t} (x(t') - \bar{x}) dt',
\end{equation}
donde \(\bar{x} = \frac{1}{T} \int_0^T x(t') dt'\) es el promedio de \(x(t)\) sobre toda la serie temporal (donde el promedio está implícitamente contenido en el detrend, tal como se describe a continuación, pero el detrend lineal es más que restar el valor promedio global). El resultado del análisis de fluctuación será un exponente de \(y(t)\) o el cual puede ser tomado como el exponente de Hurst para \(y(t)\) o bien para el proceso subyacente \(x(t)\). En lugar de realizar directamente un análisis de fluctuaciones para \(y(t)\), se introducen diferentes niveles de detrendado sustractivo. Como se explicará en el Capítulo 13, restar la tendencia no lineal suele ser razonable en muchos casos, pero las fluctuaciones pueden afectar a los datos. Las tendencias se pueden eliminar aplicando un suavizado adecuado al valor medio de \(y(t)\). Discutimos aquí el suavizado mediante polinomios de orden \(k\) sobre el mismo intervalo en el que se encuentran las fluctuaciones o el proceso.

\bigskip

Sea \( g^{(0)}(t, t') \) un polinomio de orden \( k \) con la propiedad de que \( g \) minimiza
\begin{equation}
\int_{t}^{t'} \left[ g(t,t') - y(t) \right] dt',
\end{equation}
para \( k = 0 \), la constante \( g^{(0)}(t) \) satisface el problema de minimizar el valor medio de \( y(t) \) en el intervalo de tiempo considerado. Para \( k = 1 \), \( g^{(1)}(t') \) es una función lineal de \( y(t') \) para \( t' \in [t, t'] \) (ver la Fig. 6.13), y así sucesivamente.
\begin{equation}
\frac{1}{\tau} \int_{t}^{t+\tau} \left[ y(t') - g^{(k)}(t,t') \right]^2 dt' = \nu^{(k)}(\tau),
\end{equation}
donde \( \nu^{(k)}(\tau) \) es la fluctuación correspondiente. En este trabajo con una variable de tiempo discreta, todos los integrales se corresponden con las sumas respectivas. Además, para obtener resultados más precisos, uno considera \( \langle \nu^{(k)}(\tau) \rangle \), promediado sobre intervalos de tiempo (pueden superponerse) de longitud \( \tau \) dentro de la serie de tiempo de longitud \( T \). Esto es compatible con el hecho de que si \( x(t) \) cumple con la Ec.~(6.11), entonces

\begin{equation}
\langle \nu^{(k)}(\tau) \rangle \propto \tau^{2+2H}
\end{equation}

Si los datos \( x(t) \) no son estacionarios y contienen tendencias y el escalado se destruye por \( x \) y \( y \), el detrendado con un polinomio de orden suficientemente alto puede restablecer el escalado para \( \langle \nu^{(k)}(\tau) \rangle \). Si la señal misma ya es auto-similar, como en el caso de un proceso de difusión, el inicio de la integración puede abandonarse (haciendo que \( y(t) = x(t) \)), y el análisis resultante de DFA es idéntico al de \( H \) asociado con \( x \). \footnote{Si \( x \) es ruido blanco, el \( y(t) \) integrado forma un proceso de difusión con \( \alpha = 1/2 \), y por lo tanto se puede asociar un exponente de Hurst de \( H = \alpha - 1 = -1/2 \) al ruido blanco. Dado que el resultado más pequeño del análisis de fluctuaciones es \( \alpha = 0 \), no se pueden determinar exponentes negativos. La integración en la Ec.~(6.12) aumenta el exponente en una unidad.} 

\newpage

Dado que \( \tau = O(T) \) solo unos pocos intervalos de tiempo constituyen el promedio en la Ec.~(6.14), los resultados para \( \tau \) grandes son problemáticos, de modo que solo se tiene \( t \leq T/4 \) para \( \tau \) pequeños. Para valores de \( \tau \) cercanos al intervalo de muestreo de los datos, el detrendado introduce un sesgo hacia escalas pequeñas de \( \langle \nu^{(k)}(\tau) \rangle \) y por lo tanto una violación adicional del escalado. Como extremo, para un detrendado de orden \( k \), un intervalo de tiempo que contenga \( k + 1 \) puntos de muestra da lugar a una fluctuación exacta, violando evidentemente la Ec.~(6.14). En parte, esto puede ser compensado aplicando la normalización \( 1 / (l(\tau) - k - 1) \) en la versión discreta de la Ec.~(6.13), donde \( l(\tau) \) denota el número de muestras cubiertas por el intervalo de tiempo \( \tau \). Una estrategia ampliamente utilizada que muestra resultados más precisos fue sugerida por Kantelhardt et al. (2001): construir surrogados de \( x(t) \) mediante una reordenación aleatoria. Sin sesgo, su análisis de fluctuaciones sin tendencia debería dar \( \alpha = 1/2 \) en toda la gama de escalas. La violación del escalado en las estimaciones numéricas refleja la estadística y la sensibilidad del procedimiento de estimación numérica, y define la función de corrección \( c(\tau) := t^{1/2} / \langle \nu^{(k)}_{\text{data}}(\tau) \rangle \). El resultado del análisis de los datos originales \( x(t) \) puede ser corregido multiplicando \( \langle \nu^{(k)}_{\text{data}}(\tau) \rangle \) por \( c(\tau) \). Si no hubiera violación del escalado, esta corrección simplemente implicaría la normalización adecuada para que las integrales de fluctuaciones sean independientes del orden de \( k \) del detrendado.

\bigskip

Al principio, te sorprenderá que el detrendado no modifique el exponente de escalado, ya que la señal puede ser no lineal y necesitar una eliminación de tendencia sustantiva. En particular, podrías preguntarte si un proceso de Brownian Path cuyos dos puntos finales están casi en el mismo valor (lo cual es el caso para el detrendado lineal) no se ve afectado por la eliminación de tendencias. Sin embargo, esto se debe al hecho de que el comportamiento de escalado de los datos se debe a las fluctuaciones en el sistema, no a su tendencia. La ausencia de violación del escalado es una consecuencia importante del análisis de fluctuaciones sin tendencia.


\end{document}
